\documentclass[11pt,oneside]{article}

% =====================================================================
%  Polynomial--Logarithmic Transseries
%  Canonical consolidated volume.  Editorial merge (2026-09-02) of the six
%  independently written arrivals filed on 2026-09-01; see the provenance
%  appendix for the absorbed sources and their SHA-256 receipts.
%
%  House style follows the canonical corpus volumes: A4, Libertinus with an
%  lmodern fallback, article class carrying \part/\section/\subsection.
%  All mathematical notation is normative from
%  docs/FabiusFunction_Mathematical_Notation_Catalogue.
% =====================================================================

\usepackage[a4paper,margin=26mm,headheight=15pt]{geometry}
\usepackage[T1]{fontenc}
\usepackage[utf8]{inputenc}
\IfFileExists{libertinus.sty}{\usepackage{libertinus}}{\usepackage{lmodern}}
\usepackage{microtype}
\usepackage{amsmath,amssymb,amsthm,mathtools,mathrsfs,bm}
% Canonical mathematical notation for active FabiusFunction LaTeX documents.
%
% This file is intentionally a plain TeX input rather than a LaTeX package.
% Every standalone document loads it by an explicit path relative to that
% document's directory; included fragments inherit the definitions from their
% root document.  Keep this file independent of document layout and metadata.
%
% Contract:
%   * one command has one mathematical meaning throughout the active corpus;
%   * names encode semantic roles rather than a document-local choice of letter;
%   * visually similar but differently normalized objects have different print;
%   * every command below is defined normatively in the notation catalogue.
%
% The active corpus excludes docs/archive/.  Historical sources may retain
% their original notation only inside an explicitly labelled quotation.

\makeatletter
\@ifundefined{mathbb}{%
  \PackageError{fabius-notation}{The amssymb package must be loaded first}%
    {Load amsmath and amssymb before inputting fabius-notation.tex.}%
}{}
\@ifundefined{operatorname}{%
  \PackageError{fabius-notation}{The amsmath package must be loaded first}%
    {Load amsmath before inputting fabius-notation.tex.}%
}{}
\makeatother

% Version stamp used by the catalogue and by static audits.
\newcommand{\FabiusNotationVersion}{2026-08-31}

% Number systems.  NaturalNumbers is explicitly the nonnegative convention.
\newcommand{\NaturalNumbers}{\mathbb{N}_{0}}
\newcommand{\PositiveIntegers}{\mathbb{N}_{>0}}
\newcommand{\IntegerNumbers}{\mathbb{Z}}
\newcommand{\RationalNumbers}{\mathbb{Q}}
\newcommand{\RealNumbers}{\mathbb{R}}
\newcommand{\ComplexNumbers}{\mathbb{C}}
\newcommand{\ExtendedRealNumbers}{\overline{\mathbb{R}}}

% The three principal Fabius--Rvachev functions.  Subscripts are deliberate:
% a bare F and a calligraphic F have several unrelated uses in the corpus.
\newcommand{\FabiusBounded}{F_{\mathrm{bd}}}
\newcommand{\FabiusGlobal}{F_{\mathrm{gl}}}
\newcommand{\FabiusQuantile}{Q_{F}}
\newcommand{\FabiusClampedQuantile}{Q_{F,\mathrm{cl}}}
\newcommand{\RvachevUp}{\operatorname{up}}

% Constants, elementary operators, and delimiters.
\newcommand{\EulerE}{\mathrm{e}}
\newcommand{\ImaginaryUnit}{\mathrm{i}}
\newcommand{\Differential}{\mathop{}\!\mathrm{d}}
\newcommand{\AbsoluteValue}[1]{\left\lvert #1\right\rvert}
\newcommand{\Norm}[1]{\left\lVert #1\right\rVert}
\newcommand{\Floor}[1]{\left\lfloor #1\right\rfloor}
\newcommand{\Ceiling}[1]{\left\lceil #1\right\rceil}
\newcommand{\FractionalPart}[1]{\left\{#1\right\}}
\newcommand{\PositivePart}[1]{\left(#1\right)_{+}}
\newcommand{\IndicatorOf}[1]{\mathbf{1}_{#1}}
\newcommand{\SetOf}[1]{\left\{#1\right\}}
\newcommand{\InnerProduct}[1]{\left\langle #1\right\rangle}
\newcommand{\CardinalityOf}[1]{\left\lvert #1\right\rvert}
\newcommand{\DefinitionEquals}{\mathrel{\coloneqq}}
\newcommand{\Sign}{\operatorname{sgn}}
\newcommand{\IdentityOperator}{\operatorname{id}}
\newcommand{\SpanOperator}{\operatorname{span}}
\newcommand{\TraceOperator}{\operatorname{tr}}
\newcommand{\RankOperator}{\operatorname{rank}}
\newcommand{\DiagonalOperator}{\operatorname{diag}}
\newcommand{\ResidueOperator}{\operatorname*{Res}}
\newcommand{\LeastCommonMultiple}{\operatorname{lcm}}
\newcommand{\OrderOperator}{\operatorname{ord}}
\newcommand{\PrincipalLogarithm}{\operatorname{Log}}
\newcommand{\FallingFactorial}[2]{\left(#1\right)^{\underline{#2}}}
\newcommand{\RisingFactorial}[2]{\left(#1\right)^{\overline{#2}}}

% Support, topology, convolution, and metric notation.
\newcommand{\SupportOperator}{\operatorname{supp}}
\newcommand{\SupportOf}[1]{\SupportOperator\!\left(#1\right)}
\newcommand{\TopologicalSupportOf}[1]{\operatorname{tsupp}\!\left(#1\right)}
\newcommand{\Convolution}{\mathbin{*}}
\newcommand{\MetricDistance}{\operatorname{dist}}
\newcommand{\TotalVariationDistance}{d_{\mathrm{TV}}}

% Probability.  Equality and convergence in law are not metric distance.
\newcommand{\Probability}{\mathbb{P}}
\newcommand{\Expectation}{\mathbb{E}}
\newcommand{\Variance}{\operatorname{Var}}
\newcommand{\Covariance}{\operatorname{Cov}}
\newcommand{\LawOperator}{\operatorname{Law}}
\newcommand{\LawOf}[1]{\LawOperator\!\left(#1\right)}
\newcommand{\EqualInLaw}{\mathrel{\overset{\mathrm{law}}{=}}}
\newcommand{\ConvergesInLaw}{\mathrel{\xrightarrow{\mathrm{law}}}}

% Fourier and integral transforms.  The printed labels expose normalization.
% FourierTwoPi uses exp(-2*pi*i*x*xi); FourierAngular uses exp(-i*x*omega).
\newcommand{\FourierTwoPi}[1]{\widehat{#1}_{\,2\pi}}
\newcommand{\FourierAngular}[1]{\widehat{#1}_{\,\mathrm{ang}}}
\newcommand{\LaplaceTransformSymbol}{\mathcal{L}}
\newcommand{\MellinTransformSymbol}{\mathcal{M}}
\newcommand{\LaplaceTransformOf}[1]{\LaplaceTransformSymbol\!\left[#1\right]}
\newcommand{\MellinTransformOf}[1]{\MellinTransformSymbol\!\left[#1\right]}
\newcommand{\MomentGeneratingFunction}[1]{M_{#1}}
\newcommand{\CharacteristicFunction}[1]{\varphi_{#1}}

% The two standard sinc functions must never share one symbol.
\newcommand{\SincRad}{\operatorname{sinc}_{\mathrm{rad}}}
\newcommand{\SincPi}{\operatorname{sinc}_{\pi}}

% Binary arithmetic and Thue--Morse notation.
\newcommand{\BinaryDigitSum}{\operatorname{wt}_{2}}
\newcommand{\ThueMorseSignSymbol}{\varepsilon}
\newcommand{\ThueMorseSign}[1]{\varepsilon_{#1}}
\newcommand{\PadicValuation}[1]{\nu_{#1}}
\newcommand{\TwoAdicValuation}{\nu_{2}}

% q-series notation.  Every base is an explicit argument.
\newcommand{\QPochhammer}[3]{\left(#1;#2\right)_{#3}}
\newcommand{\GaussianBinomial}[3]{%
  \genfrac{[}{]}{0pt}{}{#1}{#2}_{#3}%
}
\newcommand{\QInteger}[2]{\left[#1\right]_{#2}}

% Named special functions.
\newcommand{\Polylogarithm}{\operatorname{Li}}
\newcommand{\LambertW}{W}
\newcommand{\GammaFunction}{\Gamma}
\newcommand{\RiemannZeta}{\zeta}

% Asymptotics.  BigO and LittleO are operator tokens for existing formula
% grammar; prefer the At forms so that the limiting regime is printed.
\newcommand{\BigO}{\mathcal{O}}
\newcommand{\LittleO}{o}
\newcommand{\BigOAt}[2]{\mathcal{O}_{#2}\!\left(#1\right)}
\newcommand{\LittleOAt}[2]{o_{#2}\!\left(#1\right)}

\usepackage{booktabs,longtable,tabularx,array,multirow}
\usepackage{graphicx}
\usepackage{xcolor}
\usepackage{enumitem}
\usepackage{fancyhdr}
\usepackage{titlesec}
\usepackage{float}
\usepackage{xurl}
\usepackage{listings}
\usepackage[most]{tcolorbox}
\usepackage{hyperref}
\usepackage[nameinlink,noabbrev,capitalise]{cleveref}

\definecolor{linkblue}{RGB}{24,72,124}
\definecolor{deepgreen}{RGB}{23,102,69}
\definecolor{deepred}{RGB}{140,42,42}
\definecolor{deepgold}{RGB}{122,91,19}
\definecolor{shadegray}{RGB}{246,247,249}
\definecolor{rulegray}{RGB}{195,201,208}
\definecolor{palered}{RGB}{251,236,236}
\definecolor{palegold}{RGB}{251,245,230}
\definecolor{palegreen}{RGB}{233,245,243}

\hypersetup{
  hypertexnames=false,% ed.: avoid duplicate destinations across parts
  colorlinks=true,
  linkcolor=linkblue,
  citecolor=deepgreen,
  urlcolor=linkblue,
  pdftitle={Polynomial-Logarithmic Transseries: Algebra, Composition, Reversion, and the Lambert W Archetype},
  pdfauthor={Vladimir Reshetnikov; ProveIt formal-development synthesis},
  pdfsubject={Consolidated volume on the polynomial-logarithmic transseries scale, its arithmetic, composition, compositional inversion, and the Lambert W expansion},
  pdfkeywords={transseries, polynomial-logarithmic scale, Hahn series, series reversion, Lagrange-Burmann, Lambert W, Wright omega, Stirling numbers, Dulac series}
}

\pagestyle{fancy}
\fancyhf{}
\fancyhead[L]{\small Polynomial--logarithmic transseries}
\fancyhead[R]{\small\nouppercase{\leftmark}}
\fancyfoot[C]{\thepage}
\renewcommand{\headrulewidth}{0.4pt}

\setcounter{tocdepth}{2}
\setcounter{secnumdepth}{3}
\setlist[itemize]{leftmargin=1.45em,itemsep=0.25ex,topsep=0.55ex}
\setlist[enumerate]{leftmargin=1.85em,itemsep=0.35ex,topsep=0.55ex}
\allowdisplaybreaks[2]

\theoremstyle{plain}
\newtheorem{theorem}{Theorem}[section]
\newtheorem{proposition}[theorem]{Proposition}
\newtheorem{lemma}[theorem]{Lemma}
\newtheorem{corollary}[theorem]{Corollary}
\theoremstyle{definition}
\newtheorem{definition}[theorem]{Definition}
\newtheorem{example}[theorem]{Example}
\newtheorem{algorithm}[theorem]{Algorithm}
\newtheorem{exercise}[theorem]{Exercise}
\theoremstyle{remark}
\newtheorem{remark}[theorem]{Remark}

\crefname{exercise}{Exercise}{Exercises}
\Crefname{exercise}{Exercise}{Exercises}
\crefname{algorithm}{Algorithm}{Algorithms}
\Crefname{algorithm}{Algorithm}{Algorithms}

\newtcolorbox{warningbox}[1][]{%
  enhanced,breakable,colback=palered,colframe=deepred,
  boxrule=0.6pt,arc=1.2mm,left=1.8mm,right=1.8mm,top=1.2mm,bottom=1.2mm,
  title={A common trap},fonttitle=\bfseries,#1}
\newtcolorbox{summarybox}[1][]{%
  enhanced,breakable,colback=palegold,colframe=deepgold,
  boxrule=0.6pt,arc=1.2mm,left=1.8mm,right=1.8mm,top=1.2mm,bottom=1.2mm,
  title={Section summary},fonttitle=\bfseries,#1}
\newtcolorbox{checkpointbox}[1][]{%
  enhanced,breakable,colback=palegreen,colframe=deepgreen,
  boxrule=0.6pt,arc=1.2mm,left=1.8mm,right=1.8mm,top=1.2mm,bottom=1.2mm,
  title={Structural checkpoint},fonttitle=\bfseries,#1}

\lstdefinestyle{pseudo}{%
  basicstyle=\ttfamily\small,columns=fullflexible,frame=single,
  rulecolor=\color{rulegray},backgroundcolor=\color{shadegray},
  showstringspaces=false,keepspaces=true,
  xleftmargin=0.4em,xrightmargin=0.4em,aboveskip=0.8em,belowskip=0.8em}
\lstset{style=pseudo}

% ---------------------------------------------------------------------
%  Document-local shorthands.  Every one of them expands to a command
%  defined normatively in the notation catalogue; none introduces a new
%  mathematical meaning.  See the notation appendix.
% ---------------------------------------------------------------------
\newcommand{\K}{\mathbb{K}}
\newcommand{\PLpow}{\PolynomialLogPowerSeriesRing{\K}}   % K[L][[t]]
\newcommand{\PLlau}{\PolynomialLogLaurentRing{\K}}       % K[L]((t))
\newcommand{\PLrat}{\RationalLogLaurentField{\K}}        % K(L)((t))
\newcommand{\PLit}{\IteratedLaurentLogField{\K}}         % K((L^{-1}))((t))
\newcommand{\Gnear}{\NearIdentityPolynomialLogGroup{\K}} % X + K[L][[t]]
\newcommand{\ordt}{\nu_{t}}                              % outer valuation
\DeclareMathOperator{\lm}{lm}
\DeclareMathOperator{\lc}{lc}
\DeclareMathOperator{\degL}{deg_{L}}
\newcommand{\Trunc}[2]{\operatorname{Trunc}_{<#2}#1}     % exclusive cutoff
\newcommand{\DX}{D_{X}}                                  % distinguished derivation
\newcommand{\EulerOp}{\vartheta}                         % Euler operator X D_X
\newcommand{\shiftop}[2]{\Delta_{#1,#2}}                 % repeated-derivative operator
\newcommand{\compose}{\mathbin{\circ}}
\newcommand{\rev}[1]{\CompositionalInverseOf{#1}}
\newcommand{\recip}[1]{\MultiplicativeInverseOf{#1}}
\newcommand{\stirlingone}[2]{\UnsignedStirlingFirstKind{#1}{#2}}
\newcommand{\stirlingsigned}[2]{\SignedStirlingFirstKind{#1}{#2}}
\newcommand{\LamP}[1]{\LambertPolynomial{#1}}
\newcommand{\LamPsc}[1]{\ScaledLambertPolynomial{#1}}
\newcommand{\RevErr}[2]{\operatorname{Err}_{\mathrm{rev}}\!\left(#1,#2\right)}
\newcommand{\dprofile}[1]{d_{#1}}                        % log-degree profile
\newcommand{\precdom}{\prec}
\newcommand{\succdom}{\succ}
\newcommand{\preceqdom}{\preccurlyeq}
\newcommand{\asymdom}{\asymp}

\begin{document}
\title{\textbf{Polynomial--Logarithmic Transseries}\\[0.35em]
  \Large Algebra, Composition, Series Reversal, and the Lambert
  \texorpdfstring{\(\LambertW\)}{W} Archetype}
\author{Vladimir Reshetnikov\\
  \small ProveIt formal-development synthesis}
\date{2 September 2026}

\maketitle

\begin{abstract}
This volume is a self-contained operational treatment of the smallest
transseries scale that is closed under the operations one actually performs
when solving an implicit equation with a logarithmic term: the
polynomial--logarithmic scale generated by \(t=X^{-1}\) and \(L=\log X\) as
\(X\to+\infty\).  Its elements are formal series
\(\sum_{n\ge n_{0}}p_{n}(L)t^{n}\) whose coefficient blocks are polynomials in
\(L\).

We develop the four ambient classes \(\PLpow\subset\PLlau\subset\PLrat\subset
\PLit\) together with the valuation, leading data, and truncation calculus that
makes every computation finite; the arithmetic (multiplication, the exact unit
criterion, reciprocals and division, powers, logarithms, exponentials, and the
closure boundary); the differential calculus of the distinguished derivation
\(\DX=t\partial_{L}-t^{2}\partial_{t}\), including the repeated-derivative
operator and the resonance that governs antiderivatives; composition, the exact
formal Taylor formula \(F\compose(X+H)=\sum_{k\ge0}H^{k}\DX^{k}F/k!\), and the
near-identity composition group; and series reversal at infinity by four
independent routes -- triangular correction, functional iteration, Newton's
method with its precision-doubling law, and a Lagrange--B\"urmann formula
proved directly in the formal setting.

The Lambert \(\LambertW\) function is treated as the archetype rather than as
an application.  The inverse of \(Y\mapsto Y+\log Y\) is Wright's
\(\WrightOmega\), whose expansion coefficients are the canonical Lambert
polynomials \(\LamP{n}\); we prove their operator definition, recurrence,
Stirling-number closed form, generating equation, and integral transform, and
derive the branch expansions of \(\LambertW_{k}\) from them.  Throughout, a
formal identity and an analytic estimate are kept rigorously apart: the
transfer from a formal residual to an error bound on an actual inverse branch
is stated as a theorem with its hypotheses printed, and no convergence claim is
made that is not proved.

Two further parts cover the passage from formal transseries to Poincar\'e
asymptotics -- including power--log series near a finite point and the Dulac
conventions -- and the algorithms, residual certificates, and diagnostics
needed to compute with these objects reliably.  Appendices collect the
Bell-polynomial bookkeeping, closed-form coefficient extraction, Lambert
polynomial tables, an operational formula sheet, exercises with solutions, a
concordance with the corpus notation catalogue, and the provenance of the six
absorbed sources.
\end{abstract}

\tableofcontents
\clearpage
\section{Reader's guide and the structure of this volume}
\label{plt:sec:readers-guide}

\subsection{The eight parts}
\label{plt:sec:mot-parts}

The volume is organized so that each part depends only on those before it, with
one deliberate exception: Part~V may be read directly after Part~IV by a
reader who wants only the Lambert expansion.

\begin{center}
\begin{tabular}{@{}llp{0.52\textwidth}@{}}
\hline
Part & Subject & What it establishes \\
\hline
I & The scale &
  The generators \(t=X^{-1}\), \(L=\log X\); dominance and the rank-two
  valuation; the rings \(\PLpow\), \(\PLlau\), \(\PLrat\), \(\PLit\);
  coefficient support, leading data, exclusive truncation
  \(\Trunc{F}{N}\), and \(t\)-adic completeness. \\
II & Arithmetic and differential calculus &
  Cauchy convolution block by block; the exact unit criterion and the
  reciprocal and quotient recurrences; the distinguished derivation
  \(\DX=t\partial_{L}-t^{2}\partial_{t}\) and the Euler operator
  \(\EulerOp\); repeated derivatives \(\shiftop{n}{k}\); antidifferentiation
  and the resonant block; powers, logarithms and exponentials of small
  series. \\
III & Composition &
  Admissible inner arguments; substitution on the generators; the exact formal
  Taylor formula and its formal summability; finite coefficient formulas;
  associativity and the identity; the near-identity group \(\Gnear\); the
  precise points at which composition leaves the class. \\
IV & Series reversal &
  Existence and uniqueness of \(\rev{F}\) on \(\Gnear\); four algorithms ---
  contractive fixed point, triangular coefficient matching, Newton iteration
  with precision doubling, and Lagrange--B\"urmann at infinity; residual
  certificates \(\RevErr{F}{H}\); reversion of a nontrivial leading
  monomial. \\
V & Wright omega, the Lambert polynomials, and Lambert \(\LambertW\) &
  The reversion of \(X+L\); the canonical zero-based family \(\LamP{n}\) and
  the scaled integer family \(\LamPsc{n}\); recurrence, generating equation,
  operator formula, closed form through unsigned Stirling cycle numbers,
  degree law, Laplace transform; all complex branches, the lower real branch
  \(\LambertW_{-1}\), and the square-root branch point. \\
VI & From formal transseries to analytic asymptotics &
  What a formal identity does and does not give; blockwise Poincaré
  expansions with the logarithmic degree printed; transport of asymptotics
  through the algebraic operations; from a formal residual to an analytic
  inverse error; convergence versus divergence; branch consistency. \\
VII & Algorithms, certificates, diagnostics &
  Sparse data structures for \(n\mapsto p_{n}(L)\); truncated primitives and
  their complexity; reversion strategies compared; verification invariants,
  residual certificates, and a decision procedure for the common failure
  modes. \\
VIII & Extensions and the wider landscape &
  Puiseux and Hahn exponent groups; rational-log and Laurent-log coefficients;
  iterated logarithms and multiseries; power-log and Dulac classes at zero;
  the relation to the logarithmic-exponential transseries field and to
  Hardy fields
  \cite{hardy-orders,edgar-beginners,edgar-composition,vanderhoeven-book,adh-book,vdDMM-1997,vdDMM-2001,costin-book,mantova-monotonicity,edgar-fractional,hahn}. \\
\hline
\end{tabular}
\end{center}

\subsection{Appendices}
\label{plt:sec:mot-appendices}

\begin{itemize}
\item \emph{Bell polynomials and Fa\`a di Bruno bookkeeping.}  The partial and
  complete exponential Bell polynomials
  \(\ExponentialPartialBellPolynomial{n}{k}\),
  \(\ExponentialCompleteBellPolynomial{n}\) in the form used by the
  composition and reversion formulas.
\item \emph{Closed-form coefficient extraction.}  Explicit formulas for
  \([t^{n}]F\) and \([t^{n}L^{j}]F\) after each operation, with the
  conventions of \cref{plt:rmk:mot-normalizations} in force.
\item \emph{Lambert coefficient tables.}  \(\LamP{n}\) and \(\LamPsc{n}\)
  through a stated order, with the Stirling diagonals and a reproducible
  residual check for each row.
\item \emph{Operational formula sheet.}  One page per operation: hypotheses,
  formula, cost, and the certificate that validates a computed truncation.
\item \emph{Exercises with solutions.}
\item \emph{Notation concordance.}  The mapping from each source draft's local
  letters to the normative catalogue commands.
\item \emph{Provenance.}  Which source contributed which result, and the
  editorial repairs applied to it.
\end{itemize}

\subsection{Reading paths}
\label{plt:sec:mot-paths}

\begin{enumerate}
\item \emph{Only the Lambert \(\LambertW\) expansion.}  Read
  Part~I \(\to\) Part~IV \(\to\) Part~V.  Part~I supplies the scale, the ring
  and the valuation; Part~IV supplies the one theorem needed --- that a
  near-identity element of \(\Gnear\) has a compositional inverse in
  \(\Gnear\), computable at every order by a finite calculation; Part~V
  performs the reversion of \(X+L\) and reads off every branch.  The only
  results imported from Parts~II and~III by this path are the block derivative
  \eqref{plt:eq:mot-block-derivative} and the formal Taylor formula, both
  quoted with precise statements where used.
\item \emph{The full calculus.}  Read Parts~I--IV in order.  This is the
  self-contained algebraic core.
\item \emph{Analytic use.}  After Part~I, read Part~VI.  A reader who intends
  to convert a formal truncation into a numerical bound should read Part~VI
  before computing anything, because none of the formal machinery by itself
  licenses an analytic claim.
\item \emph{Implementation.}  Read Part~I, the recurrences of Parts~II--IV,
  then Part~VII and the formula-sheet appendix.
\item \emph{Context.}  Read Part~VIII for the embedding of this class into
  Hahn series, Dulac classes, Hardy fields, and the full
  logarithmic-exponential transseries field.
\end{enumerate}

\begin{checkpointbox}
After Part~I a reader should be able to: write an element of \(\PLlau\) in the
normal form \eqref{plt:eq:mot-normal-form}; compute \(\ordt\), \(\lm\),
\(\lc\), \(\degL\) and the log-degree profile \(\dprofile{F}\) of a given
series; decide from the leading block alone whether the series is a unit; state
which of \(\PLpow\), \(\PLlau\), \(\PLrat\), \(\PLit\) a given expression
belongs to; and say, for a displayed truncation, whether the tail is a formal
valuation tail or an analytic remainder.
\end{checkpointbox}

\subsection{Standing conventions and normalizations}
\label{plt:sec:mot-conventions}

\begin{remark}[Normalizations adopted, and how they differ from the source
drafts]
\label{plt:rmk:mot-normalizations}
This volume merges six independently written treatments.  Where their
conventions collided, the notation catalogue of the corpus is the arbiter, and
the following choices are in force everywhere.
\begin{enumerate}
\item \emph{Regime and field.}  Unless a statement says otherwise, \(X\to
  +\infty\) along the positive real ray, \(\K\) is a field of characteristic
  zero, and \(t=X^{-1}\), \(L=\log X\).  Complex work names its sector and its
  branch of \(\log\).  Local letters used by individual sources ---
  \(\lambda\), \(\ell\), \(x\) --- are renamings of \(L\) or \(X\) and are not
  additional generators.
\item \emph{Truncation is exclusive.}  \(\Trunc{F}{N}=\sum_{n<N}p_{n}(L)t^{n}\)
  keeps every logarithmic monomial in every retained block; ``through outer
  order \(t^{N-1}\)'' means the exclusive cutoff \(<N\).  One source uses the
  inclusive \([F]_{\le N}\); the conversion is
  \([F]_{\le N}=\Trunc{F}{N+1}\), a one-block difference that must be applied
  whenever a formula is transported from that source.
\item \emph{Coefficient extraction.}  \([t^{n}]F=p_{n}(L)\) and
  \([t^{n}L^{j}]F=a_{n,j}\), with the series written to the right of the
  bracket.  The bracket glyph also denotes unsigned Stirling cycle numbers
  \(\stirlingone{n}{k}\) with two integer indices; the two roles are told
  apart by argument type and never abbreviated to the bare glyph in prose.
\item \emph{Lambert polynomials.}  The canonical zero-based family
  \(\LamP{n}\) of \eqref{plt:eq:mot-lambert-first} is used throughout, with
  \(\LamP{0}(u)=-u\).  The degree law \(\deg\LamP{n}=n\) is asserted only for
  \(n\ge1\).  The scaled family is \(\LamPsc{n}=n!\LamP{n}\in
  \IntegerNumbers[u]\) for \(n\ge1\); it is a storage normalization and not a
  different coefficient sequence.  The bare symbol \(P_{n}\) is not reserved
  and is not used for this family.
\item \emph{The three \(O\)-roles} are kept apart as in the warning box of
  \cref{plt:sec:mot-three-statuses}.
\item \emph{Two inverses.}  \(\recip{F}\) is the multiplicative reciprocal and
  \(\rev{F}\) the compositional inverse.  They are unrelated operations and the
  ambiguous notation \(F^{-1}\) is not used for either.
\item \emph{Filtration, not ideal.}  In the Laurent ring \(\PLlau\) the
  generator \(t\) is a unit, so \(t^{N}\PLpow\) is an additive
  valuation-filtration tail and \emph{not} an ideal; the principal ideal
  generated by \(t^{N}\) is all of \(\PLlau\).  Statements about agreement
  below outer order \(N\) are therefore phrased with \(\ordt(F-G)\ge N\), with
  \(\FormalBigOAt{t^{N}}{t}\), or with equality of exclusive truncations, never
  as a congruence modulo an ideal.
  % ed.: source draft 2 tabulates the output class of differentiation as
  % "t K[L]((t)) up to the leading power".  That set equals K[L]((t)) because t
  % is a unit there, so the row is vacuous.  The correct statement is the
  % valuation inequality nu_t(D_X F) >= nu_t(F) + 1, proved in Part II, which
  % can be strict only through annihilation of a constant leading block.
\item \emph{Coefficient blocks are polynomials.}  In \(\PLpow\) and
  \(\PLlau\) every \(p_{n}\) lies in \(\K[L]\), so logarithmic exponents are
  nonnegative and each block is finite; no uniform bound on \(\deg_{L}p_{n}\)
  is implied by membership, and any such bound is a hypothesis or a proved
  invariant.
  % ed.: source draft 4 prints the coefficient condition of its defining
  % display as "p_n in K[X]", with the large variable in place of the
  % logarithmic one.  The intended and correct condition is p_n in K[L].
\end{enumerate}
\end{remark}

\clearpage
\part{The polynomial--logarithmic scale}

% ---- fragment 01-motivation ----
% ---------------------------------------------------------------------------
% Cluster: mot -- Part I opening.  Motivation and reader's guide.
% ---------------------------------------------------------------------------

\section{Why powers and logarithms belong in one expansion}
\label{plt:sec:motivation}

This volume is about a single asymptotic scale and the exact algebra it
supports.  The scale is generated by two symbols,
\begin{equation}\label{plt:eq:mot-generators}
  t\DefinitionEquals X^{-1},
  \qquad
  L\DefinitionEquals\log X,
  \qquad X\to+\infty,
\end{equation}
and its monomials are the products \(t^nL^j\).  A reader who has met only
power series may reasonably ask why two generators are needed at all, and, if
two, why exactly these two.  This section answers both questions by exhibiting
the smallest problem that forces them --- inverting \(y\mapsto y+\log y\) ---
and then proving three separate minimality statements: a pure power scale
cannot continue past the first term, a pure logarithmic scale cannot even
start, and the ring generated by \(t\) and \(L\) is the smallest subring of
germs at \(+\infty\) that contains a logarithm and is closed under
differentiation.

Throughout this section \(X\) is a real variable tending to \(+\infty\), all
logarithms are the real logarithm on \((0,\infty)\), and \(f\sim g\) means
\(f/g\to1\).  Complex branches are introduced only in
Part~V, where every branch of \(\LambertW\) and every choice of
\(\log\) is printed with the statement it belongs to.

\subsection{The inversion problem that starts everything}
\label{plt:sec:mot-inversion}

The whole subject is visible in one implicit equation with one logarithmic
term.

\begin{proposition}[The basic inverse and its elementary envelope]
\label{plt:prop:mot-omega-basic}
The map \(\Phi(y)=y+\log y\) is a strictly increasing bijection from
\((0,\infty)\) onto \(\RealNumbers\).  Denote its inverse by \(\WrightOmega\),
so that
\begin{equation}\label{plt:eq:mot-omega-equation}
  \WrightOmega(X)+\log\WrightOmega(X)=X
  \qquad(X\in\RealNumbers).
\end{equation}
Then \(\WrightOmega\) is real analytic, strictly increasing,
\(\WrightOmega(X)\to+\infty\) as \(X\to+\infty\), and for every \(X\ge1\),
writing \(L=\log X\),
\begin{equation}\label{plt:eq:mot-envelope}
  X-L\;\le\;\WrightOmega(X)\;\le\;X .
\end{equation}
Moreover \(\LambertW_0(z)=\WrightOmega(\log z)\) for every \(z>0\), where
\(\LambertW_0\) is the principal real branch of Lambert's function, defined by
\(\LambertW_0(z)\EulerE^{\LambertW_0(z)}=z\).
\end{proposition}

\begin{proof}
On \((0,\infty)\) we have \(\Phi'(y)=1+1/y>0\), so \(\Phi\) is strictly
increasing and real analytic with nowhere vanishing derivative;
\(\Phi(y)\to-\infty\) as \(y\downarrow0\) and \(\Phi(y)\to+\infty\) as
\(y\to+\infty\), so \(\Phi\) is a bijection onto \(\RealNumbers\).  The inverse
function theorem gives that \(\WrightOmega\) is real analytic and strictly
increasing, and \(\WrightOmega(X)\to+\infty\) as \(X\to+\infty\).

Fix \(X\ge1\) and put \(y=\WrightOmega(X)\).  If \(y<1\) then \(\log y<0\) and
hence \(X=y+\log y<1\), a contradiction; so \(y\ge1\) and \(\log y\ge0\).  From
\(y=X-\log y\) we get \(y\le X\); and then \(\log y\le\log X=L\), whence
\(y=X-\log y\ge X-L\).  This is \eqref{plt:eq:mot-envelope}.

For the last claim let \(z>0\) and put \(w=\WrightOmega(\log z)\).  Then
\(w>0\) and \(w+\log w=\log z\), so \(\log(w\EulerE^{w})=\log z\) and therefore
\(w\EulerE^{w}=z\).  Since \(v\mapsto v\EulerE^{v}\) is strictly increasing on
\((-1,\infty)\) and \(w>0>-1\), the value \(w\) is the principal branch
\(\LambertW_0(z)\).
\end{proof}

The envelope \eqref{plt:eq:mot-envelope} already contains the whole message of
this volume in miniature.  It says that \(\WrightOmega(X)\) is \(X\) up to a
correction which is at most \(\log X\); so the answer to a question posed
purely in powers of \(X\) is not expressible purely in powers of \(X\).  The
next subsection makes the correction explicit to two further orders.

\begin{remark}[Names]
\label{plt:rmk:mot-names}
The function \(\WrightOmega\) is Wright's omega function
\cite{wright-omega,corless-et-al-1996}.  Working with \(\WrightOmega\) rather
than with \(\LambertW\) removes one layer of logarithms from every display:
the two nested logarithms \(\log z\) and \(\log\log z\) in the classical
Lambert expansion become the single large variable \(X\) and the single
logarithmic coefficient variable \(L\).  Every result about \(\LambertW\) in
Part~V is obtained from a result about \(\WrightOmega\) by
substituting a branch coordinate for \(X\).
\end{remark}

\subsection{Two orders by hand}
\label{plt:sec:mot-two-orders}

The following computation is the motivating one.  It is short enough to do on
paper and already displays every structural feature of the scale: the
corrections are inverse powers of \(X\); their coefficients are polynomials in
\(\log X\); and the logarithmic degree grows with the inverse power.

\begin{proposition}[Two orders of the basic inverse]
\label{plt:prop:mot-two-orders}
With \(L=\log X\) and \(\WrightOmega\) as in
\cref{plt:prop:mot-omega-basic}, as \(X\to+\infty\) along the real axis,
\begin{equation}\label{plt:eq:mot-two-orders}
  \WrightOmega(X)
  =X-L+\frac{L}{X}+\frac{1}{X^{2}}\left(\frac{L^{2}}{2}-L\right)
   +\BigOAt{\frac{L^{3}}{X^{3}}}{X\to+\infty}.
\end{equation}
In particular \(\WrightOmega(X)-X\sim-L\) and
\(\WrightOmega(X)-X+L\sim L/X\).
\end{proposition}

\begin{proof}
Assume \(X\ge X_{0}\), where \(X_{0}\ge4\) is chosen so large that
\(\varepsilon\DefinitionEquals L/X\le\tfrac12\) and \(L\ge1\) for \(X\ge X_{0}\).
Write
\[
  r\DefinitionEquals\WrightOmega(X)-X+L,
  \qquad
  u\DefinitionEquals\frac{L-r}{X},
\]
so that \(\WrightOmega(X)=X(1-u)\).  By \eqref{plt:eq:mot-envelope} we have
\(0\le r\le L\) and hence
\begin{equation}\label{plt:eq:mot-u-range}
  0\le u\le\varepsilon\le\tfrac12 .
\end{equation}
Because \(\WrightOmega(X)>0\), we may take logarithms in
\(\WrightOmega(X)=X(1-u)\):
\[
  \log\WrightOmega(X)=L+\log(1-u).
\]
Substituting this into \eqref{plt:eq:mot-omega-equation} and solving for
\(\WrightOmega(X)\) gives \(\WrightOmega(X)=X-L-\log(1-u)\), that is,
\begin{equation}\label{plt:eq:mot-r-fixed-point}
  r=-\log(1-u)=\sum_{k\ge1}\frac{u^{k}}{k}.
\end{equation}
The series converges because \(0\le u\le\tfrac12\).

\emph{First estimate.}  For \(0\le u\le\tfrac12\) one has
\(u\le-\log(1-u)\le2u\), the right inequality because
\(u\mapsto-\log(1-u)/u\) increases on \((0,\tfrac12]\) with value
\(2\log2<2\) at \(u=\tfrac12\).  Hence, by \eqref{plt:eq:mot-r-fixed-point}
and \eqref{plt:eq:mot-u-range},
\begin{equation}\label{plt:eq:mot-r-crude}
  u\le r\le 2u\le2\varepsilon .
\end{equation}
Consequently \(u=(L-r)/X\ge(L-2L/X)/X\ge\tfrac12\varepsilon\) for \(X\ge4\), so
\begin{equation}\label{plt:eq:mot-r-two-sided}
  \tfrac12\varepsilon\le u\le\varepsilon,
  \qquad
  \tfrac12\varepsilon\le r\le2\varepsilon .
\end{equation}

\emph{Second estimate.}  Insert \(r=O(\varepsilon)\) into the definition of
\(u\).  By \eqref{plt:eq:mot-r-fixed-point} and \eqref{plt:eq:mot-r-two-sided},
\(r=u+O(u^{2})=\varepsilon+O(\varepsilon/X)+O(\varepsilon^{2})\), and since
\(L\ge1\) we have \(\varepsilon/X=L/X^{2}\le L^{2}/X^{2}=\varepsilon^{2}\); thus
\(r=\varepsilon+O(\varepsilon^{2})\).  Therefore
\begin{equation}\label{plt:eq:mot-u-refined}
  u=\varepsilon-\frac{r}{X}
   =\varepsilon-\frac{\varepsilon}{X}+\BigOAt{\frac{\varepsilon^{2}}{X}}{X\to+\infty}.
\end{equation}
Squaring \eqref{plt:eq:mot-u-refined} and using \(r=O(\varepsilon)\),
\[
  u^{2}=\Bigl(\varepsilon-\frac{r}{X}\Bigr)^{2}
       =\varepsilon^{2}-\frac{2\varepsilon r}{X}+\frac{r^{2}}{X^{2}}
       =\varepsilon^{2}+\BigOAt{\frac{\varepsilon^{2}}{X}}{X\to+\infty}.
\]
Finally, for \(0\le u\le\tfrac12\),
\[
  \Bigl|\sum_{k\ge3}\frac{u^{k}}{k}\Bigr|
  \le\frac{u^{3}}{3}\sum_{k\ge0}u^{k}
  \le\frac{2}{3}u^{3}\le\frac{2}{3}\varepsilon^{3}.
\]
Adding the three contributions in \eqref{plt:eq:mot-r-fixed-point},
\[
  r=u+\frac{u^{2}}{2}+\BigOAt{\varepsilon^{3}}{X\to+\infty}
   =\varepsilon-\frac{\varepsilon}{X}+\frac{\varepsilon^{2}}{2}
    +\BigOAt{\frac{\varepsilon^{2}}{X}}{X\to+\infty}
    +\BigOAt{\varepsilon^{3}}{X\to+\infty}.
\]
Now substitute \(\varepsilon=L/X\).  The displayed main terms are
\(L/X\), \(-L/X^{2}\) and \(L^{2}/(2X^{2})\), and the two error terms are
\(O(L^{2}/X^{3})\) and \(O(L^{3}/X^{3})\), both \(O(L^{3}/X^{3})\) because
\(L\ge1\).  Since \(\WrightOmega(X)=X-L+r\), this is exactly
\eqref{plt:eq:mot-two-orders}.

The two asymptotic equivalences follow: \(\WrightOmega(X)-X=-L+r\) with
\(r\le2L/X=\LittleOAt{L}{X\to+\infty}\), and
\(\WrightOmega(X)-X+L=r\sim\varepsilon=L/X\) by
\eqref{plt:eq:mot-two-orders}.
\end{proof}

\begin{remark}[What the two coefficients are]
\label{plt:rmk:mot-lambert-preview}
The coefficient blocks produced by the computation are the first two members
of the canonical zero-based Lambert family
\begin{equation}\label{plt:eq:mot-lambert-first}
  \LamP{0}(u)=-u,\qquad
  \LamP{1}(u)=u,\qquad
  \LamP{2}(u)=\frac{u^{2}}{2}-u,\qquad
  \LamP{3}(u)=\frac{u^{3}}{3}-\frac{3}{2}u^{2}+u,
\end{equation}
defined and studied in Part~V.  In that normalization
\eqref{plt:eq:mot-two-orders} reads
\[
  \WrightOmega(X)=X+\sum_{n=0}^{2}\LamP{n}(L)\,t^{n}
   +\BigOAt{L^{3}t^{3}}{X\to+\infty},
\]
the block \(n=0\) being the term \(-L\).  Continuing the same expansion by one
more step gives
\(\LamP{4}(u)=u^{4}/4-\tfrac{11}{6}u^{3}+3u^{2}-u\).  The degree law
\(\deg\LamP{n}=n\) holds for \(n\ge1\) only: \(\LamP{0}(u)=-u\) has degree one.
\end{remark}

\begin{remark}[The same computation for \(\LambertW\)]
\label{plt:rmk:mot-lambert-transfer}
Substituting \(X=\log z\) into \eqref{plt:eq:mot-two-orders} and using
\(\LambertW_0(z)=\WrightOmega(\log z)\) from
\cref{plt:prop:mot-omega-basic} gives, with the iterated logarithms
\(L_{1}=\log z\) and \(L_{2}=\log L_{1}\),
\begin{equation}\label{plt:eq:mot-W-two-orders}
  \LambertW_0(z)=L_{1}-L_{2}+\frac{L_{2}}{L_{1}}
   +\frac{L_{2}^{2}/2-L_{2}}{L_{1}^{2}}
   +\BigOAt{\frac{L_{2}^{3}}{L_{1}^{3}}}{z\to+\infty},
\end{equation}
the classical large-argument expansion \cite{corless-et-al-1996,dlmf-W}.  This
is the sense in which the algebra of this volume is ``one level down'' from
Lambert's function: the two nested logarithms of \(z\) are the single pair
\((X,L)\).  Note that \(L_{1}\) and \(L_{2}\) here are iterated logarithms and
never Lambert polynomials; the Lambert family is always written
\(\LamP{n}\).
\end{remark}

\subsection{Why one generator is never enough}
\label{plt:sec:mot-one-generator}

We now prove that neither of the two obvious one-generator scales can carry
the expansion \eqref{plt:eq:mot-two-orders}.  The failures are of different
kinds and it is worth naming them separately: the power scale \emph{stalls},
the logarithmic scale is \emph{blind}.

\begin{definition}[Asymptotic scale and expansion; working definition]
\label{plt:def:mot-scale}
An \emph{asymptotic scale} at \(+\infty\) is a set \(\mathfrak{M}\) of
functions, each positive on some half-line, such that for all distinct
\(m,m'\in\mathfrak{M}\) either \(m/m'\to0\) or \(m'/m\to0\) as
\(X\to+\infty\).  A function \(f\) defined near \(+\infty\) has the
\emph{asymptotic expansion} \(\sum_{i\ge1}c_{i}m_{i}\) with respect to
\(\mathfrak{M}\) if the \(c_{i}\in\RealNumbers\setminus\SetOf{0}\), the
\(m_{i}\in\mathfrak{M}\) satisfy \(m_{i+1}/m_{i}\to0\), and for every index
\(N\) present in the (finite or infinite) list,
\[
  f-\sum_{i\le N}c_{i}m_{i}=\LittleOAt{m_{N}}{X\to+\infty}.
\]
A finite such expansion is \emph{maximal} when it admits no further term,
i.e.\ when the remainder is not asymptotically equivalent to a constant
multiple of any element of \(\mathfrak{M}\).
\end{definition}

This is the coarse blockwise notion that suffices for a motivating discussion.
Part~VI replaces it by the sharper statements needed for analytic work,
including the remainder estimates that are uniform in \(L\) and the exact
sense in which a truncated polynomial-logarithmic series approximates a
function.

\begin{proposition}[Both one-generator scales fail]
\label{plt:prop:mot-one-generator-fails}
Set \(\psi_{1}(X)=\WrightOmega(X)-X\) and
\(\psi_{2}(X)=\WrightOmega(X)-X+\log X\).  Let
\(\mathfrak{M}_{\mathrm{pow}}=\SetOf{X^{a}:a\in\RealNumbers}\) and
\(\mathfrak{M}_{\log}=\SetOf{(\log X)^{b}:b\in\RealNumbers}\); both are
asymptotic scales at \(+\infty\).
\begin{enumerate}
\item \emph{(The power scale stalls.)}  With respect to
  \(\mathfrak{M}_{\mathrm{pow}}\), the function \(\WrightOmega\) has the
  one-term expansion \(X\), and that expansion is maximal.  Indeed
  \(\psi_{1}\sim-\log X\), so that \(X^{a}=\LittleOAt{|\psi_{1}|}{X\to+\infty}\)
  for every \(a\le0\) while
  \(\psi_{1}=\LittleOAt{X^{a}}{X\to+\infty}\) for every \(a>0\); no element of
  \(\mathfrak{M}_{\mathrm{pow}}\) is asymptotically proportional to
  \(\psi_{1}\).
\item \emph{(The logarithmic scale is blind.)}  With respect to
  \(\mathfrak{M}_{\log}\), the function \(\WrightOmega\) has no expansion at
  all, since \(\WrightOmega(X)/(\log X)^{b}\to+\infty\) for every
  \(b\in\RealNumbers\).  Worse, the residual \(\psi_{2}\) is strictly positive
  for large \(X\) and satisfies
  \(\psi_{2}=\LittleOAt{(\log X)^{b}}{X\to+\infty}\) for \emph{every}
  \(b\in\RealNumbers\); hence \(\mathfrak{M}_{\log}\) assigns \(\psi_{2}\) the
  empty expansion and cannot distinguish \(\psi_{2}\) from \(0\).
\item \emph{(A single mixed monomial repairs both failures.)}
  \(\psi_{2}(X)\sim X^{-1}\log X\).
\end{enumerate}
\end{proposition}

\begin{proof}
That the two sets are asymptotic scales is the case \(b=b'=0\), respectively
\(a=a'=0\), of \cref{plt:lem:mot-dominance} below, whose proof is
independent of this proposition.

(i) By \cref{plt:prop:mot-two-orders}, \(\psi_{1}\sim-\log X\).  Since
\(\log X\to+\infty\), for \(a\le0\) we have
\(X^{a}\le1=\LittleOAt{\log X}{X\to+\infty}\), and for \(a>0\) we have
\(\log X=\LittleOAt{X^{a}}{X\to+\infty}\).  So there is no
\(a\in\RealNumbers\) and no \(c\ne0\) with \(\psi_{1}\sim cX^{a}\): for
\(a>0\) the ratio \(\psi_{1}/X^{a}\to0\), and for \(a\le0\) it is unbounded.
Finally \(\WrightOmega\sim X\) by \eqref{plt:eq:mot-envelope}, and
\(\WrightOmega-X=\psi_{1}=\LittleOAt{X}{X\to+\infty}\), so \(X\) is a legitimate
first (and, by the above, last) term.

(ii) \(\WrightOmega(X)\ge X-\log X\to+\infty\) faster than every fixed power of
\(\log X\): indeed \((X-\log X)/(\log X)^{b}\to+\infty\) by
\cref{plt:lem:mot-dominance}.  Hence no leading term exists.  For the second
claim, \eqref{plt:eq:mot-r-two-sided} in the proof of
\cref{plt:prop:mot-two-orders} gives
\(\tfrac12\log X/X\le\psi_{2}(X)\le2\log X/X\) for \(X\ge X_{0}\); in
particular \(\psi_{2}>0\).  For any \(b\),
\(\psi_{2}(X)(\log X)^{-b}\le2(\log X)^{1-b}/X\to0\) by
\cref{plt:lem:mot-dominance}.

(iii) This is the last assertion of \cref{plt:prop:mot-two-orders}.
\end{proof}

\begin{remark}[The two failures are not symmetric]
\label{plt:rmk:mot-asymmetry}
Part~(i) is a \emph{gap}: the exact size of the correction, \(\log X\), falls
strictly between \(X^{0}\) and \(X^{a}\) for every \(a>0\), so the power scale
runs out of monomials.  Part~(ii) is a \emph{loss of resolution}: the
logarithmic scale has monomials on both sides of \(\psi_{2}\) but they are all
too coarse, so the expansion exists and is uninformative.  A source draft that
reports only ``a pure power series is impossible'' has stated the weaker of
the two obstructions.
\end{remark}

\subsection{The smallest closed class}
\label{plt:sec:mot-smallest}

The negative results above say what does not work.  The positive statement is
an exact minimality theorem, and it is where the pair \((t,L)\) is singled
out.  Let \(\mathcal{G}\) denote the ring of germs at \(+\infty\) of real
functions that are smooth on some half-line, with the derivation
\(\DX=\Differential/\Differential X\).

\begin{theorem}[The two-generator ring is the smallest differentially closed
class containing a logarithm]
\label{plt:thm:mot-smallest-differential-algebra}
Inside \(\mathcal{G}\), with \(t=X^{-1}\) and \(L=\log X\):
\begin{enumerate}
\item The smallest \(\DX\)-stable subring of \(\mathcal{G}\) containing
  \(\RationalNumbers\) and \(X\) is \(\RationalNumbers[X]\).  Every element of
  it is a polynomial; in particular it contains no unbounded element of
  subpolynomial growth, and no logarithm.
\item The smallest \(\DX\)-stable subring of \(\mathcal{G}\) containing
  \(\RationalNumbers\) and \(L\) is \(\RationalNumbers[t,L]\).
\item The smallest \(\DX\)-stable subring of \(\mathcal{G}\) containing
  \(\RationalNumbers\), \(X\) and \(L\) is
  \(\RationalNumbers[X,t,L]=\RationalNumbers[t,t^{-1}][L]\), the ring of
  Laurent polynomials in \(t\) with coefficients polynomial in \(L\).
\item \(t\) and \(L\) are algebraically independent over \(\RealNumbers\).
  Hence \(\RealNumbers[t,L]\) is a genuine two-variable polynomial ring, and a
  germ in it has a \emph{unique} representation \(\sum_{n\ge0}p_{n}(L)t^{n}\)
  with \(p_{n}\in\RealNumbers[L]\) and only finitely many \(p_{n}\ne0\).
\item On each monomial block, for \(n\in\IntegerNumbers\) and
  \(p\in\RationalNumbers[L]\),
  \begin{equation}\label{plt:eq:mot-block-derivative}
    \DX\bigl(t^{n}p(L)\bigr)=t^{n+1}\bigl(p'(L)-np(L)\bigr).
  \end{equation}
\end{enumerate}
The same statements hold with \(\RationalNumbers\) replaced by any subfield
\(\K\subseteq\RealNumbers\); more generally, over an abstract field \(\K\) of
characteristic zero the ring in (iii) is the polynomial ring
\(\K[t,t^{-1}][L]\) equipped with the unique derivation satisfying
\(\DX t=-t^{2}\) and \(\DX L=t\).
\end{theorem}

\begin{proof}
We first verify \eqref{plt:eq:mot-block-derivative}, which is (v).  As germs,
\(\DX t=-X^{-2}=-t^{2}\) and \(\DX L=X^{-1}=t\).  By the product and chain
rules,
\[
  \DX\bigl(t^{n}p(L)\bigr)
  =nt^{n-1}(-t^{2})p(L)+t^{n}p'(L)\,t
  =t^{n+1}\bigl(p'(L)-np(L)\bigr).
\]

(i) \(\RationalNumbers[X]\) contains \(\RationalNumbers\) and \(X\), and is
\(\DX\)-stable because \(\DX X=1\).  Conversely any \(\DX\)-stable subring
containing \(\RationalNumbers\) and \(X\) contains all powers of \(X\) and all
rational linear combinations of them.  A nonzero element of
\(\RationalNumbers[X]\) is a polynomial, hence either bounded or of growth
\(\asymp X^{d}\) with \(d\ge1\); no element is asymptotic to a constant
multiple of \(\log X\), since \(\log X\to\infty\) while
\(\log X/X\to0\).

(ii) Put \(R=\RationalNumbers[t,L]\).  It contains \(\RationalNumbers\) and
\(L\).  It is \(\DX\)-stable: a derivation is determined on a generated ring
by its values on generators, and \(\DX t=-t^{2}\in R\), \(\DX L=t\in R\).
Conversely, let \(S\) be \(\DX\)-stable with \(\RationalNumbers\cup\SetOf{L}
\subseteq S\).  Then \(t=\DX L\in S\), so \(S\) contains the subring generated
by \(\RationalNumbers\), \(t\) and \(L\), which is \(R\).

(iii) Put \(R'=\RationalNumbers[X,t,L]\).  Since \(Xt=1\), \(R'\) is the ring
of Laurent polynomials in \(t\) with coefficients in \(\RationalNumbers[L]\).
It is \(\DX\)-stable by \(\DX X=1\), \(\DX t=-t^{2}\), \(\DX L=t\).  Any
\(\DX\)-stable subring containing \(\RationalNumbers\), \(X\), \(L\) contains
\(\DX L=t\) and hence contains \(R'\).

(iv) Suppose \(P\in\RealNumbers[T,\Lambda]\) is nonzero and
\(P(t,L)=0\) as a germ, i.e.\ \(P(X^{-1},\log X)=0\) for all large \(X\).
Write \(P=\sum_{n\ge0}T^{n}q_{n}(\Lambda)\) with \(q_{n}\in\RealNumbers[\Lambda]\),
and let \(n_{0}\) be minimal with \(q_{n_{0}}\ne0\).  Multiplying the relation
by \(X^{n_{0}}\) gives, for all large \(X\),
\[
  q_{n_{0}}(\log X)=-\sum_{n>n_{0}}X^{-(n-n_{0})}q_{n}(\log X).
\]
Each summand on the right is \(X^{-k}\) times a polynomial in \(\log X\) with
\(k\ge1\), hence tends to \(0\) by \cref{plt:lem:mot-dominance}; as the sum is
finite, the right side tends to \(0\).  The left side, however, is a nonzero
polynomial evaluated along \(\log X\to+\infty\), so it tends to \(\pm\infty\)
or to a nonzero constant.  This contradiction proves algebraic independence,
and uniqueness of the representation follows because two representations
differ by a polynomial relation.

The final sentence is immediate: all the arguments above used only that the
coefficient ring is a subring of a characteristic-zero field, and the abstract
version replaces the germ realization by the free polynomial ring, on which
\(\DX t=-t^{2}\), \(\DX L=t\) extends uniquely to a derivation because
\(t,t^{-1},L\) are free generators.
\end{proof}

\begin{remark}[Reading the theorem]
\label{plt:rmk:mot-reading-minimality}
Part~(i) says the pure power world is closed but expressively poor: a
logarithm can never be produced from \(X\) by differentiation and ring
operations, so if a problem's answer contains a logarithm, that logarithm must
enter through an operation --- inversion, integration, composition --- that
leaves the class.  Part~(ii) says the converse asymmetry does not exist:
once \(\log X\) is admitted, \(X^{-1}\) is forced immediately, and then all of
\(\RationalNumbers[t,L]\).  There is no ``logarithms only'' differential
algebra.  This is the precise sense in which powers and logarithms belong in
one expansion.  The \(t\)-adic completion of the Laurent ring in (iii) is the
ambient object \(\PLlau=\K[L]((t))\) of Part~I.
\end{remark}

\begin{corollary}[Both exponent directions are forced by \(\WrightOmega\)]
\label{plt:cor:mot-both-generators-needed}
Let \(S\subseteq\RealNumbers^{2}\) and let
\(\mathfrak{M}_{S}=\SetOf{X^{a}(\log X)^{b}:(a,b)\in S}\); by
\cref{plt:lem:mot-dominance} this is an asymptotic scale.  If
\(\WrightOmega\) has an asymptotic expansion with respect to
\(\mathfrak{M}_{S}\) with at least three terms, then
\[
  (1,0)\in S,\qquad (0,1)\in S,\qquad (-1,1)\in S,
\]
with respective coefficients \(1\), \(-1\), \(1\).  Consequently the subgroup
of \((\RealNumbers^{2},+)\) generated by \(S\) contains
\(\IntegerNumbers^{2}\); in particular \(S\) is contained neither in
\(\RealNumbers\times\SetOf{0}\) (a pure power scale) nor in
\(\SetOf{0}\times\RealNumbers\) (a pure logarithmic scale).
\end{corollary}

\begin{proof}
By \cref{plt:lem:mot-dominance}, for \((a,b)\ne(a',b')\) the ratio
\(X^{a}(\log X)^{b}/\bigl(X^{a'}(\log X)^{b'}\bigr)\) tends to \(0\) or to
\(+\infty\); in particular a monomial \(X^{a}(\log X)^{b}\) can be
asymptotically equivalent to \(c\,X^{a'}(\log X)^{b'}\) with
\(c\in(0,\infty)\) only if \((a,b)=(a',b')\) and \(c=1\).

The first term of an expansion satisfies \(\WrightOmega\sim c_{1}m_{1}\).  By
\eqref{plt:eq:mot-envelope}, \(\WrightOmega\sim X\), so
\(c_{1}m_{1}\sim X=X^{1}(\log X)^{0}\); by the previous paragraph
\(m_{1}=X\) and \(c_{1}=1\), so \((1,0)\in S\).  The second term satisfies
\(\WrightOmega-X\sim c_{2}m_{2}\), and
\(\WrightOmega-X\sim-\log X\) by \cref{plt:prop:mot-two-orders}; hence
\(c_{2}m_{2}\sim-\log X\), which forces \(c_{2}<0\),
\(m_{2}=(\log X)^{1}\), \(c_{2}=-1\), so \((0,1)\in S\).  The third term
satisfies \(\WrightOmega-X+\log X\sim c_{3}m_{3}\), and this residual is
\(\sim X^{-1}\log X\) by \cref{plt:prop:mot-two-orders}; hence
\((-1,1)\in S\) with \(c_{3}=1\).  Finally \((1,0)\) and \((0,1)\) generate
\(\IntegerNumbers^{2}\).
\end{proof}

\begin{remark}[Scope of the corollary]
\label{plt:rmk:mot-corollary-scope}
The corollary is stated for monomial scales in the two generators, and this
restriction is not cosmetic.  Nothing here rules out an exotic scale built
from other generators --- for instance one containing \(X+X^{1/2}\) and
\(X^{1/2}\) --- in which the first terms of \(\WrightOmega\) look different.
What \emph{is} ruled out, by
\cref{plt:prop:mot-one-generator-fails}, is any scale generated by powers
alone or by logarithms alone, whatever the exponents; and what is asserted by
\cref{plt:thm:mot-smallest-differential-algebra} is that among differentially
closed rings the two-generator ring is minimal.  Together these are the honest
content of the phrase ``smallest closed class''.  A stronger uniqueness
statement --- that \(\PLlau\) is the unique minimal differential ring in which
the reversion of \(X+L\) exists --- is not proved in this volume and should
not be inferred.
\end{remark}

\subsection{The dominance order and the block decomposition}
\label{plt:sec:mot-dominance}

We now record, with proof, the elementary comparison that has been used
several times above and that organizes every series in this volume.

\begin{lemma}[Lexicographic dominance of power-logarithmic monomials]
\label{plt:lem:mot-dominance}
Let \((a,b),(a',b')\in\RealNumbers^{2}\).  As \(X\to+\infty\),
\[
  \frac{X^{a}(\log X)^{b}}{X^{a'}(\log X)^{b'}}
  \longrightarrow
  \begin{cases}
    0,&\text{if }a<a',\text{ or }a=a'\text{ and }b<b',\\
    1,&\text{if }(a,b)=(a',b'),\\
    +\infty,&\text{if }a>a',\text{ or }a=a'\text{ and }b>b'.
  \end{cases}
\]
Equivalently, writing \(\succdom\) for strict dominance and using the
generators \eqref{plt:eq:mot-generators}, for \(n,m\in\IntegerNumbers\) and
\(j,k\in\IntegerNumbers\),
\begin{equation}\label{plt:eq:mot-dominance}
  t^{n}L^{j}\succdom t^{m}L^{k}
  \quad\Longleftrightarrow\quad
  n<m,\ \text{ or }\ (n=m\ \text{ and }\ j>k).
\end{equation}
Thus \(\SetOf{X^{a}(\log X)^{b}}\) is an asymptotic scale, totally ordered by
the lexicographic order on the exponent pair \((a,b)\), reversed in the
\(t\)-exponent.
\end{lemma}

\begin{proof}
Put \(s=\log X\), so \(s\to+\infty\), and write the ratio as
\[
  X^{a-a'}(\log X)^{b-b'}=\EulerE^{(a-a')s}s^{\,b-b'}.
\]
If \(a<a'\), set \(c=a-a'<0\) and \(\delta=b-b'\).  For \(s\ge1\) we have
\(\EulerE^{cs}s^{\delta}=\exp(cs+\delta\log s)\), and
\(cs+\delta\log s\to-\infty\) because \(\log s=\LittleOAt{s}{s\to+\infty}\);
hence the ratio tends to \(0\).  The case \(a>a'\) is the reciprocal
statement.  If \(a=a'\) the ratio is \(s^{\,b-b'}\), which tends to \(0\),
equals \(1\), or tends to \(+\infty\) according as \(b<b'\), \(b=b'\), or
\(b>b'\).

For \eqref{plt:eq:mot-dominance}, apply the above with \((a,b)=(-n,j)\) and
\((a',b')=(-m,k)\): the monomial \(t^{n}L^{j}\) dominates \(t^{m}L^{k}\) iff
\(-n>-m\), i.e.\ \(n<m\), or \(n=m\) and \(j>k\).  Since the trichotomy is
exhaustive, distinct exponent pairs give ratios tending to \(0\) or
\(+\infty\), which is the definition of an asymptotic scale
(\cref{plt:def:mot-scale}).
\end{proof}

The dominance order is not merely a convention chosen to make bookkeeping
convenient: the grouping of monomials into inverse-power blocks is intrinsic.

\begin{proposition}[Blocks are the logarithmic-comparability classes]
\label{plt:prop:mot-blocks}
Let \(\mu=t^{n}L^{j}\) and \(\nu=t^{m}L^{k}\) with
\(n,m,j,k\in\IntegerNumbers\).  The following are equivalent.
\begin{enumerate}
\item \(n=m\).
\item There is an integer \(N\ge0\) with
  \((\log X)^{-N}\le\mu/\nu\le(\log X)^{N}\) for all sufficiently large \(X\).
\end{enumerate}
Hence the relation ``differ by at most a fixed power of \(\log X\)'' is an
equivalence relation on power-logarithmic monomials whose classes are exactly
the sets \(\SetOf{t^{n}L^{j}:j\in\IntegerNumbers}\) for fixed \(n\), and the
order \(\succdom\) descends to a total order on the classes given by the
natural order of \(-n\).
\end{proposition}

\begin{proof}
(i)\(\Rightarrow\)(ii): if \(n=m\) then \(\mu/\nu=L^{\,j-k}\), so
\(N=\AbsoluteValue{j-k}\) works with equality on one side.

(ii)\(\Rightarrow\)(i): suppose \(n\ne m\).  If \(n>m\) then, with
\(s=\log X\), \(\mu/\nu=\EulerE^{(m-n)s}s^{\,j-k}\) and
\[
  \frac{\mu/\nu}{s^{-N}}=\EulerE^{(m-n)s}s^{\,j-k+N}\longrightarrow0
\]
for every \(N\), by \cref{plt:lem:mot-dominance}; so the lower bound in (ii)
fails for every \(N\).  If \(n<m\) the same computation applied to
\(\nu/\mu\) shows the upper bound fails.

The last sentence follows because (ii) is visibly reflexive and symmetric, and
transitive with \(N\) replaced by the sum of the two constants.
\end{proof}

\Cref{plt:prop:mot-blocks} is the structural justification for the normal
form used throughout this volume: a polynomial-logarithmic transseries is
written
\begin{equation}\label{plt:eq:mot-normal-form}
  F=\sum_{n\ge n_{0}}p_{n}(L)\,t^{n},
  \qquad p_{n}\in\K[L],\quad n_{0}\in\IntegerNumbers,
\end{equation}
one finite polynomial per archimedean class, with the classes indexed by the
integer \(n\) and ordered by \(-n\).  The formal ring housing
\eqref{plt:eq:mot-normal-form} is \(\PLlau=\K[L]((t))\), with the
power-series subring \(\PLpow=\K[L][[t]]\); their construction, valuation and
truncation topology occupy the rest of Part~I.

\begin{remark}[Where the class ends]
\label{plt:rmk:mot-boundary}
Three failures of closure mark the exact boundary of
\eqref{plt:eq:mot-normal-form}, and each names its own minimal enlargement.
\begin{itemize}
\item \(\MultiplicativeInverseOf{(L+1)}\) is not a polynomial in \(L\).
  Unrestricted division needs the coefficient field \(\K(L)\), giving
  \(\PLrat=\K(L)((t))\); expanding those rational coefficients at
  \(L=\infty\) gives instead \(\PLit=\K((L^{-1}))((t))\).  These are different
  objects: \(1/(L+1)\) is exact in the first and becomes
  \(L^{-1}-L^{-2}+L^{-3}-\cdots\) in the second.
\item \(\log L\) is produced by composing with an inner argument whose
  leading behaviour is logarithmic; it requires a second logarithmic level.
\item \(X^{1/2}\) is produced by inverting a non-integer leading power and
  requires a Puiseux or Hahn exponent group.
\end{itemize}
% ed.: source draft 6 lists the block scale as
% "..., X^{-n}L^2, X^{-n}L, X^{-n}, X^{-n}L^{-1}, ..." , which silently places
% negative powers of L inside the polynomial-logarithmic class.  They are not
% there: K[L]((t)) has nonnegative L-exponents only, and L^{-1} belongs to the
% coefficient-field enlargements named above.
Negative powers of \(L\) do not belong to \(\PLlau\).  A display that mixes
\(L^{-1}\) into a coefficient block has already moved to \(\PLrat\) or
\(\PLit\), and must say so.  These enlargements are the subject of
Part~VIII.
\end{remark}

\subsection{Where a logarithm is created}
\label{plt:sec:mot-creation}

\Cref{plt:thm:mot-smallest-differential-algebra}(i) says that a logarithm
cannot be manufactured from powers by differentiation.  It can be
manufactured by \emph{anti}differentiation, and by its discrete analogue,
summation.  These two mechanisms account for almost every occurrence of the
scale in practice, so we isolate them.

\begin{lemma}[Block antidifferentiation and the resonant block]
\label{plt:lem:mot-block-antiderivative}
Let \(\K\) have characteristic zero, let \(c\in\K\) and consider the
\(\K\)-linear operator \(\partial_{L}-c\) on \(\K[L]\).
\begin{enumerate}
\item If \(c\ne0\), then \(\partial_{L}-c\) is a bijection of \(\K[L]\) onto
  itself and preserves degree.  Consequently, for \(n\ne1\) and
  \(p\in\K[L]\), the germ \(t^{n}p(L)\) has a unique antiderivative of the
  form \(t^{n-1}q(L)\) with \(q\in\K[L]\), and \(\deg q=\deg p\).
\item If \(c=0\), then \(\partial_{L}\) is surjective with kernel \(\K\), and
  \(\deg q=\deg(\partial_{L}q)+1\) for nonconstant \(q\).  Consequently the
  block \(t\,p(L)\) has antiderivative \(q(L)\) with \(q'=p\), unique up to an
  additive constant, and \(\deg q=\deg p+1\).
\end{enumerate}
In particular \(t\cdot1=X^{-1}\) integrates to \(L=\log X\): the outer level
\(n=1\) is the unique block at which antidifferentiation raises logarithmic
degree, and it is the only source of new logarithms in the calculus.
\end{lemma}

\begin{proof}
(i) Write \(q=\sum_{i=0}^{d}q_{i}L^{i}\) with \(q_{d}\ne0\).  Then
\((\partial_{L}-c)q=-cq_{d}L^{d}+(\text{terms of degree}<d)\), so the operator
preserves degree; it is therefore injective (a nonzero \(q\) has nonzero
image) and, being degree-preserving and linear on each finite-dimensional
space \(\K[L]_{\le d}\) of polynomials of degree at most \(d\), it maps
\(\K[L]_{\le d}\) bijectively onto itself.  Taking the union over \(d\) gives
bijectivity on \(\K[L]\).  For the antiderivative statement, apply
\eqref{plt:eq:mot-block-derivative} with \(n-1\) in place of \(n\):
\(\DX(t^{n-1}q)=t^{n}\bigl(q'-(n-1)q\bigr)=t^{n}(\partial_{L}-c)q\) with
\(c=n-1\ne0\); solve \((\partial_{L}-c)q=p\).

(ii) \(\partial_{L}\) on \(\K[L]\) is surjective because
\(\partial_{L}\bigl(L^{i+1}/(i+1)\bigr)=L^{i}\), which uses
\(\operatorname{char}\K=0\); its kernel is \(\K\); and it lowers degree by
exactly one on nonconstant polynomials.  Taking \(n=1\) in
\eqref{plt:eq:mot-block-derivative} with \(n-1=0\) gives
\(\DX q(L)=t\,q'(L)\), so the antiderivative of \(t\,p(L)\) is any \(q\) with
\(q'=p\).  Finally \(\int X^{-1}\Differential X=\log X=L\) is the case
\(p=1\), \(q=L\).
\end{proof}

\begin{lemma}[The discrete counterpart: a harmonic increment makes a logarithm]
\label{plt:lem:mot-harmonic}
Let \((w_{n})_{n\ge n_{1}}\) be real, \(w_{n}\to+\infty\), and suppose that for
some constants \(c_{0}>0\) and \(c_{1}\in\RealNumbers\),
\begin{equation}\label{plt:eq:mot-harmonic-hypothesis}
  w_{n+1}=w_{n}+c_{0}+\frac{c_{1}}{w_{n}}
   +\BigOAt{w_{n}^{-2}}{n\to\infty}.
\end{equation}
Then there is a constant \(C\) with
\begin{equation}\label{plt:eq:mot-harmonic-conclusion}
  w_{n}=c_{0}n+\frac{c_{1}}{c_{0}}\log n+C+\LittleOAt{1}{n\to\infty}.
\end{equation}
\end{lemma}

\begin{proof}
Since \(w_{n}\to+\infty\), the last two terms in
\eqref{plt:eq:mot-harmonic-hypothesis} tend to \(0\), so
\(w_{n+1}-w_{n}\to c_{0}\).  By the Cesàro theorem applied to the increments,
\(w_{n}/n\to c_{0}\); in particular \(w_{n}\sim c_{0}n\) and
\(1/w_{n}=\BigOAt{n^{-1}}{n\to\infty}\).

Put \(v_{n}=w_{n}-c_{0}n\).  Then
\(v_{n+1}-v_{n}=c_{1}/w_{n}+\BigOAt{w_{n}^{-2}}{n\to\infty}
=\BigOAt{n^{-1}}{n\to\infty}\), so summing,
\(v_{n}=\BigOAt{\log n}{n\to\infty}\).

Now bootstrap.  Since \(w_{n}=c_{0}n+v_{n}\) with \(v_{n}=O(\log n)=o(n)\),
\[
  \frac{1}{w_{n}}
  =\frac{1}{c_{0}n}\cdot\frac{1}{1+v_{n}/(c_{0}n)}
  =\frac{1}{c_{0}n}-\frac{v_{n}}{c_{0}^{2}n^{2}}
   +\BigOAt{\frac{v_{n}^{2}}{n^{3}}}{n\to\infty}
  =\frac{1}{c_{0}n}+\BigOAt{\frac{\log n}{n^{2}}}{n\to\infty}.
\]
Hence
\[
  v_{n+1}-v_{n}=\frac{c_{1}}{c_{0}n}
   +\BigOAt{\frac{\log n}{n^{2}}}{n\to\infty}.
\]
Summing from \(n_{1}\) to \(n-1\) and using
\(\sum_{k=n_{1}}^{n-1}k^{-1}=\log n+\gamma-\sum_{k<n_{1}}k^{-1}
+\LittleOAt{1}{n\to\infty}\) together with the absolute convergence of
\(\sum_{k}\log k\,/\,k^{2}\), we obtain
\(v_{n}=(c_{1}/c_{0})\log n+C+\LittleOAt{1}{n\to\infty}\) for a constant
\(C\), which is \eqref{plt:eq:mot-harmonic-conclusion}.
\end{proof}

\begin{example}[Iterating the sine]
\label{plt:ex:mot-sine}
Let \(x_{0}\in(0,\pi)\) and \(x_{n+1}=\sin x_{n}\).  Since
\(0<\sin x\le\min(x,1)\) for \(x\in(0,\pi)\), the sequence stays in
\((0,1]\subset(0,\pi)\), is decreasing, and its limit \(\ell\) satisfies
\(\ell=\sin\ell\), so \(x_{n}\downarrow0\).  Put \(u_{n}=x_{n}^{-2}\).  From
\(\sin x=x\bigl(1-x^{2}/6+x^{4}/120+O(x^{6})\bigr)\) we get
\[
  \Bigl(\frac{\sin x}{x}\Bigr)^{-2}
  =1+\frac{x^{2}}{3}+\frac{x^{4}}{15}+\BigOAt{x^{6}}{x\to0},
\]
because \((1+a)^{-2}=1-2a+3a^{2}+O(a^{3})\) with \(a=-x^{2}/6+x^{4}/120\)
gives \(1+x^{2}/3+(3/36-1/60)x^{4}+O(x^{6})\) and
\(3/36-1/60=1/15\).  Multiplying by \(u_{n}\) and using
\(u_{n}x_{n}^{2}=1\),
\[
  u_{n+1}=u_{n}+\frac13+\frac{1}{15u_{n}}
   +\BigOAt{u_{n}^{-2}}{n\to\infty}.
\]
\Cref{plt:lem:mot-harmonic} with \(c_{0}=1/3\), \(c_{1}=1/15\) gives
\(u_{n}=n/3+\tfrac15\log n+C+o(1)\), and therefore
\begin{equation}\label{plt:eq:mot-sine-asymptotic}
  x_{n}=\sqrt{\frac{3}{n}}
   \left(1-\frac{3}{10}\cdot\frac{\log n}{n}
    +\BigOAt{\frac1n}{n\to\infty}\right).
\end{equation}
The coefficient of \(n^{-1}\) is a polynomial of degree one in \(\log n\).
Continuing the same bootstrap produces, at each further order \(n^{-r}\), a
polynomial in \(\log n\) of degree \(r\): the expansion lives in the Puiseux
enlargement \(\K[L]\bigl((t^{1/2})\bigr)\) with \(X=n\).  This is de~Bruijn's
example \cite[Ch.~8]{debruijn}.
\end{example}

\begin{example}[Iterating a parabolic germ]
\label{plt:ex:mot-parabolic}
Let \(f(z)=z-z^{2}+az^{3}+O(z^{4})\) be an analytic germ at \(0\) tangent to
the identity, and let \(z_{n}=f^{\circ n}(z_{0})\) for \(z_{0}>0\) small, so
that \(z_{n}\downarrow0\).  With \(w_{n}=1/z_{n}\),
\[
  w_{n+1}=\frac{1}{z_{n}(1-z_{n}+az_{n}^{2}+O(z_{n}^{3}))}
   =w_{n}\bigl(1+z_{n}+(1-a)z_{n}^{2}+O(z_{n}^{3})\bigr)
   =w_{n}+1+\frac{1-a}{w_{n}}+\BigOAt{w_{n}^{-2}}{n\to\infty},
\]
so \cref{plt:lem:mot-harmonic} gives
\(w_{n}=n+(1-a)\log n+C+o(1)\) and hence
\[
  z_{n}=\frac1n-\frac{(1-a)\log n}{n^{2}}
   +\BigOAt{\frac{1}{n^{2}}}{n\to\infty}.
\]
The same logarithm appears in the Fatou coordinate of \(f\), which has the
shape \(-1/z+(1-a)\log(1/z)+\text{const}+o(1)\); the corresponding
normal-form and conjugacy theory is one of the classical homes of this scale.
\end{example}

\begin{summarybox}
Two mechanisms create the logarithm, and they are the same mechanism in
continuous and discrete clothing.  Antidifferentiation is degree-preserving on
every outer block except \(t^{1}\), where it raises logarithmic degree by one
(\cref{plt:lem:mot-block-antiderivative}); summation is bounded on every
increment except the harmonic one, where it produces \(\log n\)
(\cref{plt:lem:mot-harmonic}).  Once one logarithm is present,
\cref{plt:thm:mot-smallest-differential-algebra}(ii) forces the whole ring
\(\K[t,L]\).
\end{summarybox}

\subsection{Provenance: where these expansions actually arise}
\label{plt:sec:mot-provenance}

\subsubsection{Implicit equations carrying a logarithmic term}

The generic source is an implicit relation in which the unknown appears both
linearly and inside a logarithm,
\begin{equation}\label{plt:eq:mot-generic-implicit}
  y+\alpha\log y+\frac{p_{1}(\log y)}{y}+\frac{p_{2}(\log y)}{y^{2}}+\cdots=X,
  \qquad \alpha\in\K,\ p_{i}\in\K[\cdot],
\end{equation}
or its exponentiated form \(y^{\alpha}\EulerE^{y}=\EulerE^{X}\).  Dominant
balance gives \(y\sim X\); substituting that into the logarithm creates
\(\log X\); iterating creates powers of \(X^{-1}\) carrying polynomials in
\(\log X\).  The reversion theory of Part~IV shows that this heuristic is a
theorem: the left side of \eqref{plt:eq:mot-generic-implicit} lies in the
near-identity class \(\Gnear\), and \(\Gnear\) is a group under composition, so
the solution \(y\) is again of the same shape.

A classical instance outside the Lambert family is the asymptotic size of the
\(n\)-th prime.  Inverting the relation \(\pi(p_{n})=n\) against the
logarithmic-integral approximation gives Cipolla's expansion
\begin{equation}\label{plt:eq:mot-cipolla}
  \frac{p_{n}}{n}=\log n+\log\log n-1
   +\frac{\log\log n-2}{\log n}
   +\BigOAt{\frac{(\log\log n)^{2}}{(\log n)^{2}}}{n\to\infty}.
\end{equation}
Read with \(X=\log n\) and \(L=\log\log n\), the right side of
\eqref{plt:eq:mot-cipolla} is an element of
\(X+\K[L][[t]]=\Gnear\) truncated at outer order \(t^{2}\): the very same
near-identity class that houses \(\WrightOmega\).  Quantile asymptotics for
distributions with exponential-type tails have the same shape for the same
reason.

\subsubsection{Dulac maps and first-return maps of hyperbolic sectors}

Let \(v\) be a real analytic planar vector field with a hyperbolic saddle at
the origin, and let \(D\) be the transition (Dulac, or corner) map from a
transversal on one separatrix to a transversal on the other, expressed in
coordinates \(x\downarrow0\) on the transversals.  Then \(D\) has an
asymptotic expansion in the power-logarithmic scale at zero,
\begin{equation}\label{plt:eq:mot-dulac-series}
  D(x)\;\sim\;\sum_{i\ge0}x^{\lambda_{i}}P_{i}\bigl(\log(1/x)\bigr),
  \qquad
  0<\lambda_{0}<\lambda_{1}<\cdots\to+\infty,\quad P_{i}\in\RealNumbers[\cdot],
\end{equation}
with \(\lambda_{0}\) the hyperbolicity ratio of the saddle.  The logarithms
are forced exactly by resonances among the exponents \(\lambda_{i}\); for a
nonresonant saddle the \(P_{i}\) are constants.  Composing finitely many such
maps with regular transition maps produces the first-return map of a
hyperbolic polycycle, whose asymptotic expansion is again of the form
\eqref{plt:eq:mot-dulac-series}.  The finiteness of the number of limit cycles
accumulating on such a polycycle --- Dulac's problem, settled by Écalle and
by Ilyashenko --- turns on how much of this formal expansion can be promoted
to an analytic statement \cite{ecalle,mardesic-normalforms,resman-dulac}; the
\(\varepsilon\)-neighbourhood length of an orbit of a Dulac map is computed in
the same scale \cite{mardesic-length}, and the closure properties of the
associated logarithmic transseries algebras are studied in
\cite{peran-logtransseries}.

The change of variable \(X=1/x\), \(L=\log X=\log(1/x)\) carries
\eqref{plt:eq:mot-dulac-series} to \(\sum_{i}t^{\lambda_{i}}P_{i}(L)\), which
is the normal form \eqref{plt:eq:mot-normal-form} with a real exponent set;
Part~VIII treats that Hahn enlargement.  The core of this volume keeps the
integer lattice, where the algorithms are exact and finite.

\begin{warningbox}
Part of the dynamical literature uses \(\ell(x)=-1/\log x\) as its small
logarithmic generator, so that \(\ell\downarrow0\) with \(x\), rather than
\(\log(1/x)\), which tends to \(+\infty\).  The two conventions differ by an
inversion of the coefficient variable, and a formula transported between them
without that inversion will be wrong.  In this volume \(L=\log X\) always
tends to \(+\infty\).
\end{warningbox}

\subsubsection{Iteration, normal forms, and recursions}

\Cref{plt:ex:mot-sine,plt:ex:mot-parabolic} are the model cases: whenever the
increment of a slowly diverging quantity carries a term of size \(1/w\), the
partial sums manufacture a logarithm, and the whole subsequent expansion is
organized in inverse powers of \(n\) with polynomial coefficients in
\(\log n\).  In the analytic classification of parabolic germs the same
logarithm appears as the \(\log\) term of the Fatou coordinate, and in the
formal classification of maps tangent to the identity at infinity the
conjugating series is an element of \(\Gnear\).  Symbolic algorithms for
computing such expansions --- and for the more general exp-log asymptotic
inversion problem --- are due to Salvy and Shackell
\cite{salvy-shackell-1998,salvy-shackell-1999}; the algorithmic material of
Part~VII specializes their framework to the present class.

\subsubsection{The Fabius--Rvachev endpoint, where this corpus meets the scale}

The corpus in which this volume sits studies the Fabius function
\(\FabiusBounded\) and the Rvachev function \(\RvachevUp\).  Their contact
with the polynomial-logarithmic scale is at the endpoints, and it is genuine
rather than decorative: the endpoint problem is an inversion problem of
exactly the type analysed above, and its solution requires two generators.

The corpus proves --- and has kernel-verified in Lean --- the effective
flatness bounds
\begin{equation}\label{plt:eq:mot-flatness}
  \FabiusBounded(x)\le2^{\binom{n+1}{2}}x^{n},
  \qquad
  \FabiusBounded(x)\le\frac{2^{\binom{n+1}{2}}}{n!}\,x^{n},
  \qquad
  n\in\NaturalNumbers,\ x\ge0,\ 2^{n}x\le1 .
\end{equation}
Optimizing the free index \(n\) in the sharper of the two bounds turns them
into a quadratic-in-\(\log\) decay law.

\begin{proposition}[A two-term endpoint upper bound from the flatness family]
\label{plt:prop:mot-fabius-endpoint}
Write \(\ell=\log(1/x)\) and \(a=\log2\).  There is an absolute constant
\(C\) such that for all \(0<x\le\tfrac14\),
\begin{equation}\label{plt:eq:mot-fabius-endpoint}
  \log\FabiusBounded(x)
  \le-\frac{\ell^{2}}{2a}-\frac{\ell\log\ell}{a}+C\ell .
\end{equation}
\end{proposition}

\begin{proof}
Put \(m=\ell/a\); since \(x\le\tfrac14\) we have \(\ell\ge2a\) and hence
\(m\ge2\).  Let \(n=\Floor{m}\ge2\) and \(\theta=m-n\in[0,1)\).  The
constraint \(2^{n}x\le1\) is equivalent to \(na\le\ell\), i.e.\ to \(n\le m\),
which holds.  Taking logarithms in the second bound of
\eqref{plt:eq:mot-flatness} and using \(\log x=-\ell\),
\[
  \log\FabiusBounded(x)\le\frac{n(n+1)}{2}a-n\ell-\log n!
  \;=\;A+B,
  \qquad
  A\DefinitionEquals\frac{an(n+1)}{2}-n\ell,\quad
  B\DefinitionEquals-\log n! .
\]

\emph{The quadratic part.}  Using \(\ell=am\) and \(n=m-\theta\),
\[
  A=\frac{an}{2}\bigl(n+1-2m\bigr)
   =-\frac{a}{2}(m-\theta)(m+\theta-1)
   =-\frac{a}{2}\bigl(m^{2}-m+\theta-\theta^{2}\bigr).
\]
Since \(\theta-\theta^{2}\in[0,\tfrac14]\) and \(am^{2}/2=\ell^{2}/(2a)\),
\(am/2=\ell/2\), we get
\[
  A\le-\frac{\ell^{2}}{2a}+\frac{\ell}{2}.
\]

\emph{The factorial part.}  From \(\EulerE^{n}\ge n^{n}/n!\) we get
\(n!\ge n^{n}\EulerE^{-n}\), hence \(B\le-n\log n+n\).  Since \(m\ge2\) and
\(0\le\theta<1\) we have \(0\le\theta/m\le\tfrac12\), so
\(\log(1-\theta/m)\ge-2\theta/m\ge-2/m\) and therefore
\[
  n\log n=(m-\theta)\bigl(\log m+\log(1-\theta/m)\bigr)
   \ge(m-\theta)\log m-2 .
\]
Consequently, using \(0\le\theta<1\) and \(n\le m\),
\[
  B\le-m\log m+\log m+2+m .
\]
Now \(m\log m=(\ell/a)(\log\ell-\log a)=\dfrac{\ell\log\ell}{a}
-\dfrac{\ell\log a}{a}\), and \(\log a=\log\log2<0\), so
\(-\ell\log a/a\le\AbsoluteValue{\log a}\ell/a\).  Also
\(\log m=\log(\ell/a)\le\ell\) and \(m=\ell/a\).  Hence
\[
  B\le-\frac{\ell\log\ell}{a}
   +\Bigl(\frac{\AbsoluteValue{\log a}+1}{a}+1+\frac{2}{\ell}\Bigr)\ell
  \le-\frac{\ell\log\ell}{a}+C_{1}\ell
\]
for an absolute constant \(C_{1}\), using \(\ell\ge2a>1\).  Adding the two
estimates gives \eqref{plt:eq:mot-fabius-endpoint} with
\(C=C_{1}+\tfrac12\).
\end{proof}

\begin{corollary}[The inverse endpoint lives on a Puiseux-logarithmic scale]
\label{plt:cor:mot-fabius-inverse}
Let \(\FabiusQuantile\) denote the inverse of the strictly increasing
restriction \(\FabiusBounded|_{[0,1]}\), and put \(B=\log(1/y)\).  Then, as
\(y\downarrow0\),
\begin{equation}\label{plt:eq:mot-fabius-inverse}
  \log\frac{1}{\FabiusQuantile(y)}
  \le\sqrt{2B\log2}+\BigOAt{1}{y\downarrow0}.
\end{equation}
\end{corollary}

\begin{proof}
Let \(x=\FabiusQuantile(y)\), so \(y=\FabiusBounded(x)\) and \(x\downarrow0\)
as \(y\downarrow0\); for \(y\) small we have \(x\le\tfrac14\).  With
\(\ell=\log(1/x)\) and \(a=\log2\), \cref{plt:prop:mot-fabius-endpoint} and
\(\log\ell>0\) give
\[
  B=\log\frac{1}{y}=-\log\FabiusBounded(x)\ge\frac{\ell^{2}}{2a}-C\ell .
\]
Hence \(\ell^{2}-2aC\ell-2aB\le0\), so
\(\ell\le aC+\sqrt{a^{2}C^{2}+2aB}\le\sqrt{2aB}+2aC\).  This is
\eqref{plt:eq:mot-fabius-inverse}.
\end{proof}

\begin{remark}[What the corpus adds, and what must not be inferred]
\label{plt:rmk:mot-fabius-status}
\Cref{plt:prop:mot-fabius-endpoint,plt:cor:mot-fabius-inverse} are one-sided:
they follow from a family of upper bounds, and an upper bound for
\(\FabiusBounded\) gives only an upper bound for
\(\log(1/\FabiusQuantile(y))\).  The corpus establishes the matching lower
bounds by a lattice computation on the dyadic scale ladder, obtaining the
two-sided laws
\[
  \log\FabiusBounded(x)=-\frac{\ell^{2}}{2\log2}
   -\frac{\ell\log\ell}{\log2}+\BigOAt{\ell}{x\downarrow0},
  \qquad
  \log\frac{1}{\FabiusQuantile(y)}=\sqrt{2B\log2}-\tfrac12\log B
   +\BigOAt{1}{y\downarrow0},
\]
with \(\ell=\log(1/x)\) and \(B=\log(1/y)\).  Those results are quoted here,
not proved; the proof belongs to the endpoint volumes of the corpus.

Two structural observations connect them to the present volume.  First, the
inverse law is not a polynomial-logarithmic series in the core sense: read
with \(X=B\), \(t=X^{-1}\) and \(L=\log X\), its displayed terms are
\(\sqrt{2\log2}\,X^{1/2}-\tfrac12L+\BigOAt{1}{X\to+\infty}\), an element of the
Puiseux enlargement \(\K[L]\bigl((t^{1/2})\bigr)\) of Part~VIII, not of
\(\PLlau\).  The half-integer exponent is exactly the signature of inverting a
leading behaviour that is quadratic rather than linear, and it is treated in
Part~IV under reversion beyond the near-identity class.  Second, the corpus
records that the refined dyadic endpoint expansions carry \emph{log-periodic}
corrections --- bounded one-periodic functions of \(\ell/\log2\) --- and that
the natural coordinate for the associated saddle equations is the lower real
branch \(\LambertW_{-1}\).  A bounded oscillating function of \(\ell/\log2\) is
not a monomial of any power-logarithmic scale: it is not asymptotic to a
constant multiple of \(t^{n}L^{j}\), and it therefore cannot appear as a
coefficient in \eqref{plt:eq:mot-normal-form}.  Its presence is a genuine
departure from the class studied here, and the algebra of this volume applies
to such an expansion only after the periodic factor has been separated out.
The appearance of \(\LambertW_{-1}\), on the other hand, is squarely inside
the volume: it is the branch expanded in Part~V.
\end{remark}

\subsection{Formal, asymptotic, convergent}
\label{plt:sec:mot-three-statuses}

The material above has already used three logically independent kinds of
statement, and the rest of the volume depends on never confusing them.

\begin{warningbox}
A displayed infinite expansion can have any of three statuses.
\begin{enumerate}
\item A \emph{formal identity}: an equality of coefficient arrays in a series
  ring.  Written with the formal tail \(\FormalBigOAt{t^{N}}{t}\), which means
  \(\ordt(F-G)\ge N\) and asserts no numerical bound, no convergence, and no
  uniformity in \(L\).
\item An \emph{asymptotic expansion}: a hierarchy of remainder estimates for
  an actual function on a stated domain.  Written with
  \(\BigOAt{A(X)}{X\to+\infty}\) or \(\BigOAt{A(X)}{X\to\infty,\,X\in S}\),
  with the limit and the sector printed.
\item A \emph{convergent expansion}: an actual infinite sum with a stated
  domain of convergence.
\end{enumerate}
A fourth use of the same glyph, algorithmic complexity \(\BigO(N^{2})\), is
neither of the three.  No one of these four meanings may be inferred from
another; in particular, a bare \(\BigO\) is never acceptable for an asymptotic
claim in this volume.
\end{warningbox}

For the Lambert expansion all three statuses are in fact available on suitable
domains, and this is exceptional rather than typical: the large-argument
Lambert series is not only asymptotic but convergent in appropriate real and
complex regions, a fact requiring separate analytic work
\cite{corless-et-al-1996,dlmf-W}.
% ed.: source draft 1 asserts convergence "on the principal real branch for
% z >= e" with only a general citation.  The threshold is not established
% anywhere in the six drafts, and the statement mixes a convergence claim with
% the (separate, easy) observation that the series is exact at z = e because
% every Lambert polynomial vanishes at 0.  We therefore record the exactness at
% z = e as the verifiable part and leave the convergence domain to Part VI.
What is elementary, and worth isolating, is that the expansion
\eqref{plt:eq:mot-W-two-orders} is \emph{exact} at \(z=\EulerE\): there
\(L_{1}=1\) and \(L_{2}=0\), and \(\LamP{n}(0)=0\) for every \(n\), so every
correction block vanishes and the display returns
\(\LambertW_0(\EulerE)=1\).  A convergence threshold is a different assertion
and is not proved in this volume; Part~VI states what is proved and cites what
is not.

\begin{summarybox}
The inverse of \(y\mapsto y+\log y\) cannot be expanded in powers of \(X\)
alone --- the correction \(-\log X\) falls in a gap of that scale --- nor in
powers of \(\log X\) alone, which cannot even see the correction
\(X^{-1}\log X\).  It can be expanded in the two-generator scale
\(t^{n}L^{j}\), and that scale is forced: the smallest differentially closed
ring containing a logarithm is \(\K[t,L]\).  Dominance is lexicographic, one
inverse power of \(X\) beating every fixed power of \(\log X\), so a series is
organized into inverse-power blocks, each a polynomial in \(L\).  The same
scale organizes implicit equations with a logarithmic term, Dulac and
first-return maps, parabolic normal forms, recursions with a harmonic
increment, and the endpoint asymptotics of the Fabius--Rvachev system.
\end{summarybox}

\clearpage
\part{Arithmetic and differential calculus}

% ---- fragment 03-multiplication ----
\section{Multiplication: convolution of logarithmic blocks}
\label{plt:sec:multiplication}

Multiplication is the operation on which every later construction rests:
reciprocals, powers, composition, and reversion are all built from it by
triangular recurrences.  Its whole content is that the Cauchy rule in the
outer variable~$t$ is \emph{coefficientwise finite}, so that each output
block is an exact finite expression in finitely many input blocks.  Nothing
here is an analytic rearrangement, and nothing here requires convergence.

Throughout this section $\K$ is a field of characteristic zero, $X$ is the
large variable, the regime is $X\to+\infty$ on the positive ray unless a
different domain is announced, and
\begin{equation}\label{plt:eq:mul-generators}
 t\DefinitionEquals X^{-1},
 \qquad
 L\DefinitionEquals\log X .
\end{equation}
In the formal algebra $t$ and $L$ are independent generators; their analytic
coupling is recorded only by the distinguished derivation, which plays no
role in this section.  The four ambient classes are

\begin{center}
\begin{tabular}{@{}lll@{}}
\hline
class & coefficient ring & elements \\
\hline
$\PLpow$ & $\K[L]$ & $\sum_{n\ge0}p_n(L)t^n$ \\
$\PLlau$ & $\K[L]$ & $\sum_{n\ge n_0}p_n(L)t^n$, $n_0\in\IntegerNumbers$ \\
$\PLrat$ & $\K(L)$ & $\sum_{n\ge n_0}p_n(L)t^n$, $p_n$ rational in $L$ \\
$\PLit$  & $\K((L^{-1}))$ & $\sum_{n\ge n_0}t^n\sum_{j\ge j_n}a_{n,j}L^{-j}$ \\
\hline
\end{tabular}
\end{center}

All truncations below are \emph{exclusive}: $\Trunc{F}{N}$ keeps every
logarithmic monomial in every block $t^n$ with $n<N$, and discards the rest.

\begin{remark}[Reconciliation of the source conventions]
\label{plt:rmk:mul-conventions}
Two of the merged reports use an inclusive outer cutoff $[F]_{\le N}$.  The
translation is $[F]_{\le N}=\Trunc{F}{N+1}$, and every dependency range,
algorithm bound and error exponent quoted from those reports has been shifted
by one block accordingly.  Likewise, one report writes the coefficient
bracket as $[F]_M$; here the series always stands to the right of the
bracket, $[t^n]F=p_n(L)$ and $[t^nL^j]F=a_{n,j}$.
\end{remark}

\subsection{The block Cauchy product}
\label{plt:sec:mul-cauchy}

It is economical to prove the algebra of all four classes at once.  What the
three coefficient rings $\K[L]$, $\K(L)$, $\K((L^{-1}))$ have in common is a
degree function in $L$ together with a coefficient-extraction map obeying the
usual convolution rule.

\begin{definition}[Logarithmically graded domain]
\label{plt:def:mul-graded-domain}
A \emph{logarithmically graded domain over $\K$} is a commutative
$\K$-algebra $R$ together with $\K$-linear maps
$\kappa_d\colon R\to\K$, one for each $d\in\IntegerNumbers$, such that
for every $f\in R$ the set $\SetOf{d:\kappa_d(f)\ne0}$ is nonempty and
bounded above whenever $f\ne0$, and:
\begin{enumerate}
\item[(G1)] $f=0$ if and only if $\kappa_d(f)=0$ for every $d$;
\item[(G2)] $\kappa_d(fg)=\sum_{i+j=d}\kappa_i(f)\kappa_j(g)$ for all
$f,g\in R$ and all $d\in\IntegerNumbers$, the sum having only finitely many
nonzero terms;
\item[(G3)] $\kappa_0(1)=1$ and $\kappa_d(1)=0$ for $d\ne0$.
\end{enumerate}
For $f\ne0$ put $\degL f\DefinitionEquals\max\SetOf{d:\kappa_d(f)\ne0}$ and
$\lc_L(f)\DefinitionEquals\kappa_{\degL f}(f)\in\K^\times$; set
$\degL 0\DefinitionEquals-\infty$ and $\lc_L(0)\DefinitionEquals0$.
\end{definition}

Here $\lc_L$ is the leading coefficient \emph{inside one block}; it is not
the same object as the leading scalar $\lc$ of a series, defined in
\cref{plt:def:mul-leading-data} below, although the two agree on the leading
block.

\begin{lemma}[Elementary consequences]
\label{plt:lem:mul-graded-basics}
Let $R$ be a logarithmically graded domain over $\K$.  For all $f,g\in R$:
\begin{equation}\label{plt:eq:mul-degL-laws}
 \degL(fg)=\degL f+\degL g,
 \qquad
 \lc_L(fg)=\lc_L(f)\lc_L(g),
 \qquad
 \degL(f+g)\le\max(\degL f,\degL g).
\end{equation}
In particular $R$ is an integral domain.  Moreover, if $\degL f\le d$ and
$\degL g\le e$, then $\kappa_{d+e}(fg)=\kappa_d(f)\kappa_e(g)$.
\end{lemma}

\begin{proof}
The third assertion in \eqref{plt:eq:mul-degL-laws} is immediate from
$\K$-linearity of each $\kappa_d$.  For the first two, assume $f,g\ne0$ and
set $d=\degL f$, $e=\degL g$.  If $c>d+e$ and $i+j=c$, then $i>d$ or $j>e$, so
every term of (G2) vanishes and $\kappa_c(fg)=0$.  If $c=d+e$, the only
possibly nonvanishing term is $i=d$, $j=e$, whence
\[
 \kappa_{d+e}(fg)=\kappa_d(f)\kappa_e(g)=\lc_L(f)\lc_L(g)\ne0,
\]
because $\K$ is a field and both factors are nonzero.  This proves
$\degL(fg)=d+e$, $\lc_L(fg)=\lc_L(f)\lc_L(g)$, and, by (G1), $fg\ne0$; so $R$
is a domain.  The last sentence is the same computation with $d,e$ merely
upper bounds, all other terms of (G2) again vanishing.
\end{proof}

\begin{lemma}[The three coefficient rings]
\label{plt:lem:mul-three-rings}
Each of
\[
 \K[L],\qquad \K(L),\qquad \K((L^{-1}))
\]
is a logarithmically graded domain over $\K$, with $\degL$ equal
respectively to the polynomial degree, to $\degL(p/q)=\deg p-\deg q$, and to
the largest exponent occurring in the Laurent expansion.  There are injective
$\K$-algebra homomorphisms
\begin{equation}\label{plt:eq:mul-coeff-chain}
 \K[L]\hookrightarrow\K(L)\hookrightarrow\K((L^{-1}))
\end{equation}
that preserve $\degL$ and $\lc_L$.
\end{lemma}

\begin{proof}
For $\K[L]$ take $\kappa_d(f)$ to be the coefficient of $L^d$ in $f$ for
$d\ge0$ and $\kappa_d=0$ for $d<0$; (G1)--(G3) are the definition of
polynomial multiplication.  For $\K((L^{-1}))$, an element is
$f=\sum_{j\le d}a_jL^j$ with only finitely many $j>0$ allowed but arbitrarily
negative $j$ permitted; put $\kappa_j(f)=a_j$.  Then (G1) and (G3) are clear,
and in (G2) the constraint $i+j=c$ with $i\le\degL f$, $j\le\degL g$ forces
$i\ge c-\degL g$, so the sum is finite; it is the definition of the Laurent
product.  The support of $f$ is bounded above by construction.

The inclusion $\K[L]\subset\K((L^{-1}))$ is a $\K$-algebra map, and every
nonzero element of $\K((L^{-1}))$ is invertible because $\K((L^{-1}))$ is a
Laurent-series field in the variable $L^{-1}$ over the field $\K$.  Hence
every nonzero polynomial becomes invertible, and the universal property of the
field of fractions supplies a unique $\K$-algebra homomorphism
$\iota\colon\K(L)\to\K((L^{-1}))$ extending the inclusion; it is injective
because $\K(L)$ is a field and $\iota(1)=1$.  Pulling the maps $\kappa_d$ back
along $\iota$ makes $\K(L)$ a logarithmically graded domain, and
\cref{plt:lem:mul-graded-basics} applied to $\iota(p)=\iota(q)\iota(p/q)$
gives $\degL(p/q)=\deg p-\deg q$ and
$\lc_L(p/q)=\lc_L(p)\lc_L(q)^{-1}$.  Both maps in
\eqref{plt:eq:mul-coeff-chain} preserve $\kappa_d$ by construction, hence
preserve $\degL$ and $\lc_L$.
\end{proof}

\begin{definition}[The block Cauchy product]
\label{plt:def:mul-cauchy}
Let $R$ be a logarithmically graded domain over $\K$ and let
\begin{equation}\label{plt:eq:mul-FG-blocks}
 F=\sum_{n\ge a}p_n\,t^n,
 \qquad
 G=\sum_{m\ge b}q_m\,t^m,
 \qquad p_n,q_m\in R,
\end{equation}
be elements of $R((t))$, the ring of formal Laurent series in $t$ over $R$.
Their product is
\begin{equation}\label{plt:eq:mul-cauchy}
 FG\DefinitionEquals\sum_{N\ge a+b}C_N\,t^N,
 \qquad
 C_N\DefinitionEquals\sum_{\substack{n+m=N\\ n\ge a,\ m\ge b}}p_nq_m
 =\sum_{n=a}^{N-b}p_n\,q_{N-n}.
\end{equation}
\end{definition}

\begin{lemma}[Finiteness of every fiber]
\label{plt:lem:mul-finite-fiber}
For fixed $N$ the index set in \eqref{plt:eq:mul-cauchy} is
$\SetOf{n\in\IntegerNumbers:a\le n\le N-b}$, of cardinality
$\max(N-a-b+1,\,0)$.  In particular it is finite, and it is empty exactly
when $N<a+b$.
\end{lemma}

\begin{proof}
The constraints $n\ge a$ and $m=N-n\ge b$ are equivalent to
$a\le n\le N-b$.  An interval of integers $[a,N-b]$ has $N-b-a+1$ elements
when $N-b\ge a$ and is empty otherwise.
\end{proof}

This is the entire mechanism.  There is no limit, no rearrangement, and no
hypothesis of numerical convergence: each $C_N$ is a finite sum of products
in~$R$.

\begin{theorem}[Closure of the four ambient classes]
\label{plt:thm:mul-closure}
Let $R$ be a logarithmically graded domain over $\K$.  Then
\eqref{plt:eq:mul-cauchy} makes $R((t))$ a commutative $\K$-algebra with unit
$1$, containing $R[[t]]$ as a subalgebra.  Consequently each of
\[
 \PLpow=\K[L][[t]],\quad
 \PLlau=\K[L]((t)),\quad
 \PLrat=\K(L)((t)),\quad
 \PLit=\K((L^{-1}))((t))
\]
is closed under multiplication, and the chain of inclusions
\begin{equation}\label{plt:eq:mul-class-chain}
 \PLpow\subset\PLlau\subset\PLrat\subset\PLit
\end{equation}
consists of injective $\K$-algebra homomorphisms.
\end{theorem}

\begin{proof}
By \cref{plt:lem:mul-finite-fiber} each $C_N$ is a well-defined element
of~$R$, and $C_N=0$ for $N<a+b$; so $FG$ again has support bounded below and
lies in $R((t))$.  Commutativity is the symmetry of the index set under
$(n,m)\mapsto(m,n)$, and bilinearity over $\K$ is clear, as is
$1\cdot F=F$ from (G3).  For associativity let $H=\sum_{k\ge c}h_kt^k$.  Then
\[
 [t^N]\bigl((FG)H\bigr)
 =\sum_{r+k=N}\Bigl(\sum_{n+m=r}p_nq_m\Bigr)h_k
 =\sum_{n+m+k=N}p_nq_mh_k,
\]
the last sum ranging over the finite set of triples with $n\ge a$, $m\ge b$,
$k\ge c$: indeed $n+m+k=N$ with the three lower bounds confines each index to
a finite interval, so the regrouping is a finite rearrangement.  The same
computation with the other bracketing gives the same finite sum, whence
$(FG)H=F(GH)$.  Distributivity is likewise a finite regrouping.

If $\ordt F\ge0$ and $\ordt G\ge0$, then $C_N=0$ for $N<0$, so $R[[t]]$ is
closed under the product; it is a subalgebra.  Taking $R=\K[L]$ gives the
first two classes, $R=\K(L)$ the third, $R=\K((L^{-1}))$ the fourth, using
\cref{plt:lem:mul-three-rings}.  Finally, a $\K$-algebra homomorphism
$\varphi\colon R\to R'$ induces a $\K$-algebra homomorphism
$R((t))\to R'((t))$ by $\sum p_nt^n\mapsto\sum\varphi(p_n)t^n$, because
$\varphi$ commutes with the finite sums \eqref{plt:eq:mul-cauchy}; injectivity
is inherited blockwise.  Applying this to
\eqref{plt:eq:mul-coeff-chain} yields \eqref{plt:eq:mul-class-chain}.
\end{proof}

\begin{proposition}[The monomial-level convolution and the support bound]
\label{plt:prop:mul-monomial}
Write $F=\sum_{n,j}a_{n,j}t^nL^j$ and $G=\sum_{m,k}b_{m,k}t^mL^k$ in one of
the four classes, the inner index running over $\NaturalNumbers$ in the first
three and over $\IntegerNumbers$ (bounded above at each $n$) in $\PLit$.
Then
\begin{equation}\label{plt:eq:mul-2d-convolution}
 [t^NL^d](FG)
 =\sum_{n+m=N}\ \sum_{i+j=d}a_{n,i}b_{m,j},
\end{equation}
a finite sum, and
\begin{equation}\label{plt:eq:mul-support}
 \CoefficientSupportOf{FG}
 \subseteq
 \CoefficientSupportOf{F}+\CoefficientSupportOf{G},
\end{equation}
the Minkowski sum in $\IntegerNumbers^2$.
\end{proposition}

\begin{proof}
Formula \eqref{plt:eq:mul-2d-convolution} is (G2) applied to
\eqref{plt:eq:mul-cauchy}.  Finiteness holds because the outer sum is finite
by \cref{plt:lem:mul-finite-fiber} and, for each of its terms, the inner sum
is finite by (G2).  For \eqref{plt:eq:mul-support}: if
$(N,d)\notin\CoefficientSupportOf{F}+\CoefficientSupportOf{G}$, then every
summand of \eqref{plt:eq:mul-2d-convolution} has $a_{n,i}=0$ or $b_{m,j}=0$,
so $[t^NL^d](FG)=0$.
\end{proof}

\begin{remark}[Why multiplication is a two-level, not a flat, operation]
\label{plt:rmk:mul-two-level}
Although \eqref{plt:eq:mul-2d-convolution} is a genuine two-dimensional
convolution, it should not be implemented as one.  The outer index is the
asymptotic hierarchy and the inner index is not: at fixed $N$ the inner sum
is bounded, whereas a cutoff imposed jointly on $|N|+|d|$ does not respect
dominance.  See \cref{plt:rmk:mul-rectangular-cutoff}.
\end{remark}

\begin{remark}[Hahn supports]
\label{plt:rmk:mul-hahn}
Nothing in \cref{plt:lem:mul-finite-fiber} uses that the exponents are
integers, only that the two supports are well ordered in the direction of
increasing smallness.  If $S,T\subseteq\Gamma$ are reverse-well-ordered
subsets of an ordered abelian group, then $S+T$ is reverse-well-ordered and
each $\gamma\in S+T$ has only finitely many decompositions $\gamma=s+t$ with
$s\in S$, $t\in T$; this is the Hahn support theorem
\cite{hahn,edgar-beginners}.  Everything in this section that does not
mention the specific value group $\IntegerNumbers$ transfers verbatim to the
Hahn power-log extension $\sum_{\alpha\in S}p_\alpha(L)t^\alpha$ and, in
particular, to the Puiseux-log extensions $\K[L]((t^{1/s}))$ used later for
fractional powers.  We do not repeat the statements.
\end{remark}

\subsection{Finite determination and precision bookkeeping}
\label{plt:sec:mul-finite-det}

The content of \cref{plt:lem:mul-finite-fiber} deserves to be recorded as a
dependency statement, because that is the form in which it is used: it says
exactly how much of the input must be known in order to know a given part of
the output, and it says that no more is needed and no less will do.

\begin{theorem}[Finite determination, with exact dependency range]
\label{plt:thm:mul-finite-determination}
Let $R$ be a logarithmically graded domain over $\K$, let $a,b,N\in
\IntegerNumbers$, and let $F,G\in R((t))$ satisfy $\ordt F\ge a$ and
$\ordt G\ge b$.  Then
\begin{equation}\label{plt:eq:mul-determination}
 [t^N](FG)=\sum_{n=a}^{N-b}p_n\,q_{N-n},
\end{equation}
so that $[t^N](FG)$ is a $\K$-bilinear function of the two finite tuples
\[
 (p_a,p_{a+1},\ldots,p_{N-b})
 \qquad\text{and}\qquad
 (q_b,q_{b+1},\ldots,q_{N-a}),
\]
equivalently of $\Trunc{F}{N-b+1}$ and $\Trunc{G}{N-a+1}$.  Both ranges are
exact: for every $n_0$ with $a\le n_0\le N-b$ there are $F,F'\in R((t))$ with
$\ordt\ge a$ that differ only in the block of index $n_0$, and a $G$ with
$\ordt G\ge b$, such that $[t^N](FG)\ne[t^N](F'G)$; and symmetrically in the
second argument.
\end{theorem}

\begin{proof}
Equation \eqref{plt:eq:mul-determination} and the bilinearity are
\cref{plt:def:mul-cauchy} together with \cref{plt:lem:mul-finite-fiber}; the
blocks listed are precisely those occurring in the sum.  For exactness fix
$n_0\in[a,N-b]$ and put $m_0=N-n_0$, so that $m_0\ge b$.  Take
$G=t^{m_0}$, which has $\ordt G=m_0\ge b$.  If $n_0>a$ set $F=t^a+t^{n_0}$
and $F'=t^a$; if $n_0=a$ set $F=t^a$ and $F'=2t^a$ (here $2\ne0$ because
$\operatorname{char}\K=0$).  In both cases $\ordt F=\ordt F'=a$, the two
series agree outside index $n_0$, and $[t^N](FG)-[t^N](F'G)$ equals $1$
respectively $-1$.  The symmetric statement follows by commutativity.
\end{proof}

\begin{checkpointbox}
\textbf{The precision rule.}  If $F$ is known through $\Trunc{F}{P}$ and $G$
through $\Trunc{G}{Q}$, and $\ordt F=a$, $\ordt G=b$, then $FG$ is known
through
\[
 \Trunc{FG}{\min(P+b,\;Q+a)}
\]
and no further.  Precision is lost from the factor whose \emph{partner} has
the larger valuation.  In particular, multiplying by a series with a pole of
order $r$ costs $r$ blocks of precision.
\end{checkpointbox}

\begin{corollary}[Truncated product]
\label{plt:cor:mul-truncated-product}
Let $F,G\in R((t))$ be nonzero with $a=\ordt F$, $b=\ordt G$, and let
$N\in\IntegerNumbers$.  Then
\begin{equation}\label{plt:eq:mul-truncated-product}
 \Trunc{FG}{N}
 =\Trunc{\bigl(\Trunc{F}{N-b}\bigr)\bigl(\Trunc{G}{N-a}\bigr)}{N}.
\end{equation}
If moreover $N>a+b$, neither cutoff can be lowered by one: there exist $F,G$
with these valuations for which replacing $N-b$ by $N-b-1$, or $N-a$ by
$N-a-1$, falsifies \eqref{plt:eq:mul-truncated-product}.  If
$F,G\in\PLpow$, the cruder law
\begin{equation}\label{plt:eq:mul-truncated-product-crude}
 \Trunc{FG}{N}=\Trunc{\bigl(\Trunc{F}{N}\bigr)\bigl(\Trunc{G}{N}\bigr)}{N}
 \qquad(N\ge0)
\end{equation}
also holds.
\end{corollary}

\begin{proof}
By \cref{plt:thm:mul-finite-determination} every block $[t^M](FG)$ with $M<N$
depends only on $p_n$ with $n\le M-b\le N-1-b$ and on $q_m$ with
$m\le M-a\le N-1-a$, that is, only on $\Trunc{F}{N-b}$ and $\Trunc{G}{N-a}$;
and the discarded blocks of $F$ and $G$ contribute nothing below $t^N$
because $n\ge N-b$ and $m\ge b$ force $n+m\ge N$, and symmetrically.  This is
\eqref{plt:eq:mul-truncated-product}.

For the sharpness, assume $N>a+b$ and put $F=t^a+t^{N-b-1}$ (or $F=t^a$ if
$N-b-1=a$, in which case take $F=t^a$ and note $\Trunc{F}{N-b-1}=0$) and
$G=t^b$.  Then $[t^{N-1}](FG)=1\ne0$, while
$\Trunc{(\Trunc{F}{N-b-1})G}{N}$ has zero block at $t^{N-1}$.  The second
cutoff is symmetric.

For \eqref{plt:eq:mul-truncated-product-crude} note that $a,b\ge0$ gives
$N-b\le N$ and $N-a\le N$, so $\Trunc{F}{N}$ and $\Trunc{G}{N}$ retain at
least the blocks required by \eqref{plt:eq:mul-truncated-product}; the extra
retained blocks pair only into orders $\ge N$ and are removed by the outer
truncation.
\end{proof}

\begin{warningbox}
The crude law \eqref{plt:eq:mul-truncated-product-crude} is \emph{false} on
the Laurent ring.  Take $F=t^{-1}$, $G=1+t$, $N=1$.  Then $FG=t^{-1}+1$, so
$\Trunc{FG}{1}=t^{-1}+1$, whereas
$\Trunc{F}{1}\Trunc{G}{1}=t^{-1}\cdot1=t^{-1}$.  The correct cutoff for $G$
is $N-\ordt F=2$, and indeed $t^{-1}(1+t)=t^{-1}+1$.  Implementations that
``normalize every series to start at $t^0$'' silently commit this error
whenever a genuine principal part is present.
\end{warningbox}

\begin{corollary}[Error rule for approximate factors]
\label{plt:cor:mul-error-rule}
Let $F,G\in\PLlau$ with $\ordt F=a$, $\ordt G=b$, and suppose
$\widetilde F,\widetilde G\in\PLlau$ satisfy
\[
 F=\widetilde F+\FormalBigOAt{t^{p}}{t},
 \qquad
 G=\widetilde G+\FormalBigOAt{t^{q}}{t}.
\]
Then
\begin{equation}\label{plt:eq:mul-error-rule}
 FG=\widetilde F\widetilde G
 +\FormalBigOAt{t^{\min(p+b,\;q+a,\;p+q)}}{t}.
\end{equation}
If in addition $p\ge a$ and $q\ge b$ --- that is, if neither approximation is
coarser than the leading order of the quantity it approximates --- then the
third entry is redundant and
\begin{equation}\label{plt:eq:mul-error-rule-normalized}
 FG=\widetilde F\widetilde G
 +\FormalBigOAt{t^{\min(p+b,\;q+a)}}{t}.
\end{equation}
\end{corollary}

\begin{proof}
Write $E=F-\widetilde F$ and $H=G-\widetilde G$, so $\ordt E\ge p$ and
$\ordt H\ge q$.  Expanding,
\begin{equation}\label{plt:eq:mul-error-expansion}
 FG-\widetilde F\widetilde G
 =E\widetilde G+\widetilde FH+EH .
\end{equation}
Now $\widetilde G=G-H$ gives
$\ordt\widetilde G\ge\min(\ordt G,\ordt H)\ge\min(b,q)$, and likewise
$\ordt\widetilde F\ge\min(a,p)$.  Since $\ordt$ is additive on products
(\cref{plt:thm:mul-leading-law} below, or directly
\cref{plt:lem:mul-finite-fiber}) and superadditive on sums,
\[
 \ordt\bigl(FG-\widetilde F\widetilde G\bigr)
 \ \ge\
 \min\bigl(p+\min(b,q),\ \min(a,p)+q,\ p+q\bigr)
 \ =\ \min(p+b,\ q+a,\ p+q),
\]
which is \eqref{plt:eq:mul-error-rule}.  If $p\ge a$ and $q\ge b$, then
$p+q\ge a+q$ and $p+q\ge p+b$, so $p+q\ge\min(p+b,q+a)$ and the third entry
may be dropped.
\end{proof}

% ed.: S6, Theorem "Closure and finite determination", eq. (mult-error-rule).
% ed.: The source states the conclusion as O(t^{min(p+m_0, q+n_0)}) without
% ed.: any hypothesis relating p,q to the valuations n_0,m_0, and its proof
% ed.: silently bounds ord(E*Gtilde) by p+m_0 and ord(Ftilde*H) by q+n_0,
% ed.: which requires ord(Gtilde) >= m_0 and ord(Ftilde) >= n_0.  Neither
% ed.: follows from the stated hypotheses.  The corrected statement carries
% ed.: the third entry p+q, and the hypotheses p >= n_0, q >= m_0 under
% ed.: which the source's form is recovered.
\begin{remark}[What \eqref{plt:eq:mul-error-rule} does and does not say]
\label{plt:rmk:mul-not-an-ideal}
The symbol $\FormalBigOAt{t^{N}}{t}$ asserts coefficient agreement below
outer order $N$ and nothing else: no numerical bound, no convergence, no
uniformity in $L$.  It is also not a statement about a quotient ring.  In
$\PLlau$ the element $t$ is a unit, so the ideal generated by $t^N$ is the
whole ring; the subset $t^N\PLpow$ is an additive piece of the
valuation filtration, not an ideal.  Congruences of the form
\eqref{plt:eq:mul-error-rule} must therefore be read as valuation
inequalities $\ordt(F-G)\ge N$, or equivalently as equalities of exclusive
truncations.  Quotient-algebra language becomes legitimate only after the
finite principal part has been removed and one works inside $\PLpow$, where
$(t^N)$ is a genuine ideal.
\end{remark}

The corresponding algorithm is the schoolbook convolution restricted to the
window supplied by \cref{plt:thm:mul-finite-determination}.

\begin{lstlisting}[style=pseudo,caption={Truncated block product},label={plt:alg:mul-truncated}]
# F, G sparse maps  exponent -> coefficient in R;  a = ord(F), b = ord(G)
# returns the blocks of Trunc_{<N}(F*G)
function block_product(F, G, a, b, N):
    C := empty map
    for (n, Pn) in F with a <= n <= N-b-1:
        for (m, Qm) in G with b <= m <= N-n-1:
            C[n+m] := C.get(n+m, 0) + Pn*Qm
    normalize every coefficient of C          # collect in L before zero-testing
    delete the blocks of C that are exactly 0
    return C
\end{lstlisting}

\begin{remark}[Cost, and what actually dominates it]
\label{plt:rmk:mul-cost}
With dense consecutive support the loop performs
$\BigO\bigl((N-a-b)^2\bigr)$ coefficient products; with sparse supports of
sizes $s_F,s_G$ it performs at most $s_Fs_G$.  If the retained blocks have
logarithmic degree $\BigO(d)$, each coefficient product costs
$\BigO(d^2)$ scalar operations by schoolbook polynomial arithmetic, so the
total is $\BigO\bigl((N-a-b)^2d^2\bigr)$; fast convolution in $t$ and fast
polynomial multiplication in $L$ reduce both factors.  These are complexity
$\BigO$'s and carry no analytic meaning.  In practice, at the moderate orders
typical of hand asymptotics, exact coefficient normalization --- expansion,
collection, and gcd reduction if the coefficients are rational functions of
$L$ --- costs more than the convolution itself, and the decisive
implementation choice is to never form a product whose outer order exceeds
the requested cutoff.
\end{remark}

\begin{remark}[A rectangular cutoff is not a truncation]
\label{plt:rmk:mul-rectangular-cutoff}
Discarding all monomials $t^nL^j$ with $|n|+|j|$ or $|j|$ above a threshold
does not respect dominance.  With a cap $|j|\le10$ one deletes $t^2L^{100}$
while retaining $t^3L^{-100}$, although
$t^3L^{-100}\precdom t^2L^{100}$ by an enormous margin.  Bounding
$\degL p_n$ is a second, independent approximation; it is legitimate only
when accompanied by a proved degree bound such as
\cref{plt:prop:mul-affine-profile}.
\end{remark}

\subsection{Leading data and the leading-term law}
\label{plt:sec:mul-leading}

\begin{definition}[Leading data]
\label{plt:def:mul-leading-data}
Let $R$ be a logarithmically graded domain over $\K$ and let
$F=\sum_{n\ge n_0}p_nt^n\in R((t))$ be nonzero with
$\ordt F=\min\SetOf{n:p_n\ne0}$; set $\ordt0=+\infty$.  The
\emph{leading block} of $F$ is the whole coefficient $p_{\ordt F}\in R$, the
\emph{leading monomial} is
\begin{equation}\label{plt:eq:mul-leading-monomial}
 \lm(F)\DefinitionEquals t^{\ordt F}L^{\degL p_{\ordt F}},
\end{equation}
and the \emph{leading scalar} is
$\lc(F)\DefinitionEquals\lc_L\bigl(p_{\ordt F}\bigr)\in\K^\times$.  The
\emph{rank-two valuation} of $F$ is
\begin{equation}\label{plt:eq:mul-rank-two}
 v(F)\DefinitionEquals\bigl(\ordt F,\ -\degL p_{\ordt F}\bigr)
 \in\IntegerNumbers^2,
\end{equation}
with $\IntegerNumbers^2$ lexicographically ordered and $v(0)=+\infty$.  These
are three different objects and none is determined by the coarse integer
$\ordt F$ alone.
\end{definition}

\begin{lemma}[Dominance is growth]
\label{plt:lem:mul-dominance}
On the positive ray, for integers $n,m$ and $j,k$,
\[
 t^nL^j\succdom t^mL^k
 \quad\Longleftrightarrow\quad
 \lim_{X\to+\infty}\frac{X^{-n}(\log X)^j}{X^{-m}(\log X)^k}=+\infty
 \quad\Longleftrightarrow\quad
 n<m,\ \text{or}\ (n=m\ \text{and}\ j>k).
\]
Consequently $\lm(F)$ is the $\succdom$-largest monomial of
$\CoefficientSupportOf{F}$.
\end{lemma}

\begin{proof}
The ratio equals $X^{m-n}(\log X)^{j-k}$.  If $m>n$ it tends to $+\infty$
because $\log X=\LittleOAt{X^{\varepsilon}}{X\to+\infty}$ for every
$\varepsilon>0$; if $m=n$ it is $(\log X)^{j-k}$, which tends to $+\infty$
exactly when $j>k$; if $m<n$ it tends to $0$ for the same reason as the first
case.  For the last sentence: a monomial $t^nL^j$ in the support has
$n\ge\ordt F$, with equality attained; among those with $n=\ordt F$ the
largest admissible $j$ is $\degL p_{\ordt F}$ by
\cref{plt:def:mul-graded-domain}.
\end{proof}

\begin{theorem}[Leading-term law]
\label{plt:thm:mul-leading-law}
Let $R$ be a logarithmically graded domain over $\K$ and let $F,G\in R((t))$
be nonzero.  Then $FG\ne0$ and
\begin{align}
 \ordt(FG)&=\ordt F+\ordt G,
 \label{plt:eq:mul-ord-additive}\\
 \lm(FG)&=\lm(F)\,\lm(G),
 \label{plt:eq:mul-lm-law}\\
 \lc(FG)&=\lc(F)\,\lc(G),
 \label{plt:eq:mul-lc-law}\\
 v(FG)&=v(F)+v(G).
 \label{plt:eq:mul-v-additive}
\end{align}
Moreover $\ordt(F+G)\ge\min(\ordt F,\ordt G)$ and
$v(F+G)\ge\min\bigl(v(F),v(G)\bigr)$ in the lexicographic order, both
inequalities being equalities unless the relevant leading data cancel.
\end{theorem}

\begin{proof}
Put $a=\ordt F$, $b=\ordt G$.  By \cref{plt:lem:mul-finite-fiber} the
coefficient $C_N$ of $FG$ vanishes for $N<a+b$, and for $N=a+b$ the index set
$[a,N-b]=\SetOf{a}$ is a single point, so $C_{a+b}=p_aq_b$.  Since $R$ is a
domain (\cref{plt:lem:mul-graded-basics}) and $p_a,q_b\ne0$, we get
$C_{a+b}\ne0$; hence $FG\ne0$ and \eqref{plt:eq:mul-ord-additive} holds.
Applying \eqref{plt:eq:mul-degL-laws} to $C_{a+b}=p_aq_b$ gives
\[
 \degL C_{a+b}=\degL p_a+\degL q_b,
 \qquad
 \lc_L(C_{a+b})=\lc_L(p_a)\lc_L(q_b),
\]
which are exactly \eqref{plt:eq:mul-lm-law} and \eqref{plt:eq:mul-lc-law};
\eqref{plt:eq:mul-v-additive} is the same pair of identities rewritten in the
notation \eqref{plt:eq:mul-rank-two}.

For the sum, let $c=\min(a,b)$.  Every block of $F+G$ of index $<c$ vanishes,
so $\ordt(F+G)\ge c$.  If $a<b$ then $[t^c](F+G)=p_a\ne0$ and the leading
data of $F+G$ and $F$ coincide, so $v(F+G)=v(F)=\min(v(F),v(G))$ --- note
that $a<b$ makes $v(F)<v(G)$ lexicographically.  If $a=b=c$ and
$p_c+q_c\ne0$, then $\ordt(F+G)=c$ and
$\degL(p_c+q_c)\le\max(\degL p_c,\degL q_c)$ by
\eqref{plt:eq:mul-degL-laws}, i.e.
$-\degL(p_c+q_c)\ge\min(-\degL p_c,-\degL q_c)$, which is
$v(F+G)\ge\min(v(F),v(G))$.  If $p_c+q_c=0$, then $\ordt(F+G)>c$ and
$v(F+G)>\min(v(F),v(G))$ outright.
\end{proof}

\begin{corollary}[Integrality and fields]
\label{plt:cor:mul-domain}
Each of $\PLpow$, $\PLlau$, $\PLrat$, $\PLit$ is an integral domain.  The map
$v$ of \eqref{plt:eq:mul-rank-two} is a rank-two valuation on $\PLrat$ and on
$\PLit$, with value group $\IntegerNumbers^2$ lexicographically ordered; on
$\PLlau$ it takes values in the sub-semigroup
$\IntegerNumbers\times\IntegerNumbers_{\le0}$.  The classes $\PLrat$ and
$\PLit$ are fields; $\PLpow$ and $\PLlau$ are not.
\end{corollary}

\begin{proof}
Integrality is \cref{plt:thm:mul-leading-law}, and the valuation axioms are
its last two displays together with $v(F)=+\infty$ iff $F=0$.  Surjectivity
onto $\IntegerNumbers^2$ for $\PLrat$ and $\PLit$ follows from
$v(t^nL^{-j})=(n,j)$, legitimate there because $L^{-j}$ lies in $\K(L)$ and
in $\K((L^{-1}))$; the group $\IntegerNumbers^2$ with the lexicographic order
has exactly one proper nontrivial convex subgroup,
$\SetOf{0}\times\IntegerNumbers$, so the valuation has rank two.  On $\PLlau$
the second coordinate is $-\degL p_{\ordt F}\le0$ because polynomial degrees
are nonnegative.  A Laurent-series ring $R((t))$ over a field $R$ is a field,
which covers $\PLrat$ and $\PLit$.  Finally $L\in\PLlau$ is not invertible in
$\PLlau$: if $LH=1$ then \eqref{plt:eq:mul-lm-law} gives
$L\cdot\lm(H)=1$, forcing $\degL$ of the leading block of $H$ to be $-1$,
impossible for a polynomial.
\end{proof}

\begin{corollary}[Necessity half of the unit criterion]
\label{plt:cor:mul-unit-necessary}
If $F\in\PLlau$ is invertible in $\PLlau$, then its leading block is a
nonzero scalar: $p_{\ordt F}\in\K^\times$.  The same statement holds in
$\PLpow$ with the additional conclusion $\ordt F=0$.
\end{corollary}

\begin{proof}
If $FH=1$ then \eqref{plt:eq:mul-ord-additive} gives
$\ordt F+\ordt H=0$ and \eqref{plt:eq:mul-lm-law} gives
$\degL p_{\ordt F}+\degL h_{\ordt H}=0$.  Both degrees are nonnegative
integers, hence both vanish, so $p_{\ordt F}\in\K^\times$.  In $\PLpow$ both
orders are $\ge0$ and sum to $0$, so both are $0$.
\end{proof}

The converse --- that a scalar leading block suffices --- is a division
statement and is proved with the reciprocal recurrence; it is not a
multiplication fact and is not repeated here.

\subsection{Degree and support bounds}
\label{plt:sec:mul-degree}

\begin{definition}[Log-degree profile]
\label{plt:def:mul-profile}
For $F=\sum_np_nt^n$ in one of the four classes, the \emph{log-degree
profile} is the function $\dprofile{F}\colon\IntegerNumbers\to
\IntegerNumbers\cup\SetOf{-\infty}$,
\begin{equation}\label{plt:eq:mul-profile}
 \dprofile{F}(n)\DefinitionEquals
 \begin{cases}
  \degL p_n,&p_n\ne0,\\
  -\infty,&p_n=0.
 \end{cases}
\end{equation}
The profile is data additional to $\ordt F$; it is not recoverable from it.
For two profiles define the \emph{max-plus convolution}
\begin{equation}\label{plt:eq:mul-maxplus}
 \bigl(\dprofile{F}\boxplus\dprofile{G}\bigr)(N)
 \DefinitionEquals
 \max_{n+m=N}\bigl(\dprofile{F}(n)+\dprofile{G}(m)\bigr),
\end{equation}
with the conventions $-\infty+c=-\infty$ and $\max\emptyset=-\infty$; by
\cref{plt:lem:mul-finite-fiber} the maximum is over a finite set.
\end{definition}

\begin{theorem}[Degree bound and exact criterion for attainment]
\label{plt:thm:mul-degree-bound}
Let $F,G\in R((t))$ with $R$ a logarithmically graded domain over $\K$, let
$N\in\IntegerNumbers$, and set
$M\DefinitionEquals(\dprofile{F}\boxplus\dprofile{G})(N)$.  Let
\[
 P_N\DefinitionEquals
 \SetOf{(n,m):n+m=N,\ p_n\ne0,\ q_m\ne0,\
        \dprofile{F}(n)+\dprofile{G}(m)=M}
\]
be the (finite, possibly empty) set of maximizing pairs.  Then
\begin{equation}\label{plt:eq:mul-degree-bound}
 \dprofile{FG}(N)\le M,
\end{equation}
and the coefficient of $L^{M}$ in the block $[t^N](FG)$ is exactly
\begin{equation}\label{plt:eq:mul-top-coefficient}
 \kappa_{M}\bigl([t^N](FG)\bigr)
 =\sum_{(n,m)\in P_N}\lc_L(p_n)\,\lc_L(q_m).
\end{equation}
Hence equality holds in \eqref{plt:eq:mul-degree-bound} if and only if the
sum \eqref{plt:eq:mul-top-coefficient} is nonzero.
\end{theorem}

\begin{proof}
Each summand of $C_N=\sum_{n+m=N}p_nq_m$ has
$\degL(p_nq_m)=\dprofile{F}(n)+\dprofile{G}(m)\le M$ by
\eqref{plt:eq:mul-degL-laws}, with the convention that a vanishing factor
gives $-\infty$.  A finite sum of elements of $\degL\le M$ has
$\degL\le M$, which is \eqref{plt:eq:mul-degree-bound}.  If $M=-\infty$ the
index set is empty or all summands vanish, and both sides of
\eqref{plt:eq:mul-top-coefficient} are read as $0$; assume $M\in
\IntegerNumbers$.  By $\K$-linearity of $\kappa_M$,
\[
 \kappa_M(C_N)=\sum_{n+m=N}\kappa_M(p_nq_m).
\]
If $(n,m)\notin P_N$ then either a factor vanishes or
$\degL(p_nq_m)<M$, and in both cases $\kappa_M(p_nq_m)=0$.  If
$(n,m)\in P_N$ then $\degL(p_nq_m)=M$ and
$\kappa_M(p_nq_m)=\lc_L(p_nq_m)=\lc_L(p_n)\lc_L(q_m)$ by
\cref{plt:lem:mul-graded-basics}.  This proves
\eqref{plt:eq:mul-top-coefficient}, and equality in
\eqref{plt:eq:mul-degree-bound} means precisely $\kappa_M(C_N)\ne0$.
\end{proof}

% ed.: S6, Proposition "Degree and support bounds", asserts equality "whenever
% ed.: the maximizing blocks are nonzero" under the hypothesis that all
% ed.: coefficients are nonnegative over an ordered field.  That hypothesis is
% ed.: far stronger than necessary and the quoted sufficient condition is
% ed.: stated ambiguously (a zero block has degree -infinity and cannot
% ed.: maximize at all).  The criterion below is exact and the positivity
% ed.: hypothesis has been reduced to the leading coefficients of the
% ed.: maximizing blocks only.
\begin{corollary}[Three sufficient conditions for equality]
\label{plt:cor:mul-degree-equality}
In the situation of \cref{plt:thm:mul-degree-bound}, equality holds in
\eqref{plt:eq:mul-degree-bound} in each of the following cases.
\begin{enumerate}
\item $P_N$ is a singleton.
\item $N=\ordt F+\ordt G$ (both factors nonzero); then $P_N$ is the
singleton $\SetOf{(\ordt F,\ordt G)}$ and one recovers
\eqref{plt:eq:mul-lm-law}.
\item $\K$ is an ordered field and $\lc_L(p_n)>0$, $\lc_L(q_m)>0$ for every
$(n,m)\in P_N$.
\end{enumerate}
\end{corollary}

\begin{proof}
(1) The sum \eqref{plt:eq:mul-top-coefficient} is a single product of two
elements of $\K^\times$, hence nonzero.  (2) By
\cref{plt:lem:mul-finite-fiber} the only admissible pair at
$N=\ordt F+\ordt G$ is $(\ordt F,\ordt G)$, and both blocks are nonzero, so
$P_N$ is that singleton; apply (1).  (3) The sum is a nonempty sum of
strictly positive elements of an ordered field.
\end{proof}

\begin{example}[The bound is attained, at every order and by every pair]
\label{plt:ex:mul-degree-sharp}
Over $\K=\RationalNumbers$ let
\[
 F=G=\sum_{n\ge0}L^nt^n\in\PLpow,
 \qquad\text{so}\qquad
 \dprofile{F}(n)=\dprofile{G}(n)=n .
\]
Then $(\dprofile{F}\boxplus\dprofile{G})(N)=\max_{n+m=N}(n+m)=N$, and
\emph{every} one of the $N+1$ admissible pairs is maximizing, so $P_N$ is as
large as it can be and \cref{plt:cor:mul-degree-equality}(1) does not apply.
Nevertheless
\[
 [t^N](FG)=\sum_{n=0}^{N}L^n\cdot L^{N-n}=(N+1)L^N,
\]
so $\dprofile{FG}(N)=N$ for every $N$: the bound
\eqref{plt:eq:mul-degree-bound} is attained at every order.  Formula
\eqref{plt:eq:mul-top-coefficient} predicts exactly the top coefficient
$N+1=\sum_{(n,m)\in P_N}1\cdot1$.
\end{example}

\begin{warningbox}
\cref{plt:ex:mul-degree-sharp} is the place where characteristic zero is
used.  Over a field of characteristic $p>0$ the same computation gives
$[t^{p-1}](FG)=p\,L^{p-1}=0$, so the block vanishes entirely and
$\dprofile{FG}(p-1)=-\infty$, a drop of $p-1$ below the bound.  Every
statement in this section is asserted over a field of characteristic zero.
\end{warningbox}

\begin{example}[The bound can fail by an arbitrary amount]
\label{plt:ex:mul-degree-drop}
Let $F=1+L^{d}t$ and $G=1-L^{d}t$ in $\PLpow$, $d\ge1$.  Then
$(\dprofile{F}\boxplus\dprofile{G})(1)=\max(0+d,\ d+0)=d$, while
$[t^1](FG)=-L^{d}+L^{d}=0$, so $\dprofile{FG}(1)=-\infty$.  Here
$P_1=\SetOf{(0,1),(1,0)}$ and the sum
\eqref{plt:eq:mul-top-coefficient} is $1\cdot(-1)+1\cdot1=0$, as it must be.
A drop by exactly one occurs for $F=1+(L^2+L)t$, $G=1-(L^2-L)t$, where
$[t^1](FG)=2L$ has degree $1=d-1$.
\end{example}

\begin{proposition}[Affine profiles and the profile subalgebra]
\label{plt:prop:mul-affine-profile}
Let $\alpha,\alpha'\ge0$ and $\beta,\gamma\in\RealNumbers$.  Let
$F,G\in\PLlau$ be nonzero with $a=\ordt F$, $b=\ordt G$ and
\[
 \dprofile{F}(n)\le\alpha n+\beta\ \ (n\ge a),
 \qquad
 \dprofile{G}(m)\le\alpha'm+\gamma\ \ (m\ge b).
\]
Put $\mu=\max(\alpha,\alpha')$.  Then for all $N$,
\begin{equation}\label{plt:eq:mul-affine-bound}
 \dprofile{FG}(N)\le\mu N+\beta+\gamma-(\mu-\alpha)a-(\mu-\alpha')b .
\end{equation}
In particular, if $a,b\ge0$ then $\dprofile{FG}(N)\le\mu N+\beta+\gamma$.
Consequently, for each $\alpha\ge0$ the set
\[
 \mathfrak A_\alpha
 \DefinitionEquals
 \SetOf{F\in\PLlau:\ \exists\beta\in\RealNumbers,\
        \dprofile{F}(n)\le\alpha n+\beta\ \text{for all }n}
\]
is a $\K$-subspace with
$\mathfrak A_\alpha\cdot\mathfrak A_{\alpha'}\subseteq
\mathfrak A_{\max(\alpha,\alpha')}$, and
$\mathfrak A\DefinitionEquals\bigcup_{\alpha\ge0}\mathfrak A_\alpha$ is a
$\K$-subalgebra of $\PLlau$ containing $\K[L][t,t^{-1}]$.
\end{proposition}

\begin{proof}
Fix $N$ and $n+m=N$ with $n\ge a$, $m\ge b$.  Then
\[
 \alpha n+\alpha'm
 =\mu(n+m)-(\mu-\alpha)n-(\mu-\alpha')m
 \le\mu N-(\mu-\alpha)a-(\mu-\alpha')b,
\]
because $\mu-\alpha\ge0$, $\mu-\alpha'\ge0$, $n\ge a$ and $m\ge b$.  Adding
$\beta+\gamma$ and taking the maximum over the finite index set gives
\eqref{plt:eq:mul-affine-bound} via \eqref{plt:eq:mul-degree-bound}.  If
$a,b\ge0$ then the two subtracted terms are $\ge0$, so they may be dropped.

For the algebra statements: $\mathfrak A_\alpha$ is a subspace because
$\dprofile{F+G}(n)\le\max(\dprofile{F}(n),\dprofile{G}(n))
\le\alpha n+\max(\beta,\gamma)$ and scalars do not change profiles.  The
product statement is \eqref{plt:eq:mul-affine-bound}, whose right-hand side is
$\mu N+\text{const}$.  Since $\mathfrak A_\alpha\subseteq\mathfrak
A_{\alpha''}$ for $\alpha\le\alpha''$, the union $\mathfrak A$ is closed under
addition and multiplication and contains $1$; every Laurent polynomial
$\sum_{n=n_1}^{n_2}p_n(L)t^n$ lies in $\mathfrak A_0$ with
$\beta=\max_n\degL p_n$.
\end{proof}

\begin{remark}[$\mathfrak A$ is not $t$-adically closed]
\label{plt:rmk:mul-profile-not-closed}
The affine-profile condition is preserved by finite sums and products but not
by $t$-adic limits.  Take $F_r=t^rL^{r^2}$ for $r\ge1$; then
$\ordt F_r=r\to\infty$, so $(F_r)$ is summable, yet the sum
$S=\sum_{r\ge1}t^rL^{r^2}$ has $\dprofile{S}(r)=r^2$, which is bounded by no
affine function.  Hence $S\notin\mathfrak A$.  A degree bound must therefore
be re-derived, not inherited, whenever an infinite construction is performed.
\end{remark}

\begin{proposition}[Profile of a power]
\label{plt:prop:mul-power-profile}
Let $F\in\PLpow$ with $\dprofile{F}(n)\le\alpha n+\beta$ for all $n\ge0$,
where $\alpha\ge0$ and $\beta\ge0$.  Then for every integer $k\ge0$,
\begin{equation}\label{plt:eq:mul-power-profile}
 \dprofile{F^k}(N)\le\alpha N+k\beta
 \qquad(N\ge0).
\end{equation}
\end{proposition}

\begin{proof}
Induction on $k$.  For $k=0$ we have $F^0=1$, so $\dprofile{F^0}(0)=0\le0$
and $\dprofile{F^0}(N)=-\infty$ for $N>0$.  Assume the bound for $k$.  By
\eqref{plt:eq:mul-degree-bound} applied to $F^{k+1}=F^k\cdot F$,
\[
 \dprofile{F^{k+1}}(N)
 \le\max_{n+m=N,\ n,m\ge0}\bigl((\alpha n+k\beta)+(\alpha m+\beta)\bigr)
 =\alpha N+(k+1)\beta .
\]
\end{proof}

\begin{remark}[The Lambert profile, corrected]
\label{plt:rmk:mul-lambert-profile}
Two of the sources record the degree law of the canonical Lambert family as
$\deg\LamP{n}=n$ without restriction, and use it to assert that products of
Lambert-type expansions have profile $\le N$.  The law
$\deg\LamP{n}=n$ holds only for $n\ge1$: the normalization
$\LamP{0}(u)=-u$ has degree $1$, not $0$.  The correct hypothesis for a
series $F=\sum_{n\ge0}\LamP{n}(L)t^n$ is therefore
$\dprofile{F}(n)\le n+1$ for all $n\ge0$, i.e. $\alpha=1$, $\beta=1$ in
\cref{plt:prop:mul-affine-profile}; a product of two such series has
$\dprofile{FG}(N)\le N+2$, and a $k$-th power has
$\dprofile{F^k}(N)\le N+k$ by \cref{plt:prop:mul-power-profile}.  The
unrestricted claim $\dprofile{FG}(N)\le N$ is false already at $N=0$, where
$[t^0](FG)=\LamP{0}(L)^2=L^2$ has degree $2$.
\end{remark}

\subsection{Powers and binomial powers}
\label{plt:sec:mul-powers}

Infinite sums enter through powers, so we first record the exact sense in
which they are legitimate.

\begin{definition}[Summable family]
\label{plt:def:mul-summable}
A family $(S_i)_{i\in I}$ in $\PLlau$ is \emph{summable} if for every
$N\in\IntegerNumbers$ the set $\SetOf{i\in I:\ordt S_i<N}$ is finite.  Its
sum is defined blockwise:
$[t^n]\sum_{i}S_i\DefinitionEquals\sum_i[t^n]S_i$, a finite sum in the
coefficient ring.
\end{definition}

\begin{lemma}[Summable families are well behaved]
\label{plt:lem:mul-summable}
Let $(S_i)_{i\in I}$ and $(T_j)_{j\in J}$ be summable families in $\PLlau$.
\begin{enumerate}
\item $\inf_i\ordt S_i>-\infty$, so $\sum_iS_i$ is a genuine element of
$\PLlau$ (its support is bounded below).
\item The family $(S_iT_j)_{(i,j)\in I\times J}$ is summable and
$\bigl(\sum_iS_i\bigr)\bigl(\sum_jT_j\bigr)=\sum_{(i,j)}S_iT_j$.
\end{enumerate}
\end{lemma}

\begin{proof}
(1) Taking $N=0$ in \cref{plt:def:mul-summable}, only finitely many $i$ have
$\ordt S_i<0$, and each of those has $\ordt S_i\in\IntegerNumbers$ or
$+\infty$; the minimum of a finite set of integers is finite.  Hence
$\ordt S_i\ge\min(0,\text{that minimum})$ for all $i$, and every block of
$\sum_iS_i$ below that bound vanishes.

(2) Let $\alpha=\inf_i\ordt S_i$ and $\beta=\inf_j\ordt T_j$, both finite by
(1).  Fix $N$.  If $\ordt(S_iT_j)=\ordt S_i+\ordt T_j<N$ then
$\ordt S_i<N-\beta$ and $\ordt T_j<N-\alpha$, and each of those conditions
holds for only finitely many indices; so the family is summable.  For the
identity, fix $n$ and compute in the coefficient ring, all sums having only
finitely many nonzero terms:
\[
 [t^n]\Bigl(\sum_iS_i\sum_jT_j\Bigr)
 =\sum_{p+q=n}\Bigl(\sum_i[t^p]S_i\Bigr)\Bigl(\sum_j[t^q]T_j\Bigr)
 =\sum_{i,j}\sum_{p+q=n}[t^p]S_i\,[t^q]T_j
 =\sum_{i,j}[t^n](S_iT_j).
\]
\end{proof}

\begin{remark}
\label{plt:rmk:mul-summable-vs-cauchy}
Part (1) is worth isolating.  For a \emph{sequence} in a Laurent ring,
$t$-adic Cauchyness alone does not guarantee a limit with a finite principal
part; a common eventual lower support bound must be imposed.  For a
\emph{summable family} no such extra hypothesis is needed, because the
finiteness condition at $N=0$ already forces a common lower bound.
\end{remark}

\begin{proposition}[Coefficients of an integer power]
\label{plt:prop:mul-integer-power}
Let $F=\sum_{n\ge0}p_nt^n\in\PLpow$ and $k\ge0$.  Then
\begin{equation}\label{plt:eq:mul-power-coefficient}
 [t^N]F^k
 =\sum_{\substack{n_1,\ldots,n_k\ge0\\ n_1+\cdots+n_k=N}}
   p_{n_1}\cdots p_{n_k}
 \DefinitionEquals C_{N,k}(p_0,p_1,\ldots),
\end{equation}
with $C_{0,0}=1$ and $C_{N,0}=0$ for $N>0$.  These \emph{convolution
polynomials} satisfy
\begin{equation}\label{plt:eq:mul-Cmr-recurrence}
 C_{N,k+1}=\sum_{n=0}^{N}p_n\,C_{N-n,k},
\end{equation}
which computes the whole triangular table in $\BigO(N^2k)$ coefficient
products.  If $p_0=0$, then $C_{N,k}=0$ for $N<k$, and, writing
$F=\sum_{n\ge1}p_nt^n$,
\begin{equation}\label{plt:eq:mul-bell}
 [t^N]F^k
 =\frac{k!}{N!}\,
 \ExponentialPartialBellPolynomial{N}{k}
   \bigl(1!\,p_1,\,2!\,p_2,\ldots,(N-k+1)!\,p_{N-k+1}\bigr).
\end{equation}
\end{proposition}

\begin{proof}
Formula \eqref{plt:eq:mul-power-coefficient} follows by induction on $k$ from
\cref{plt:def:mul-cauchy}: the case $k=0$ is the convention, and
\[
 [t^N]F^{k+1}=\sum_{n+m=N}p_n\,[t^m]F^k
 =\sum_{n=0}^{N}p_n\!\!\sum_{\substack{n_1+\cdots+n_k=N-n}}\!\!
   p_{n_1}\cdots p_{n_k},
\]
which is both \eqref{plt:eq:mul-Cmr-recurrence} and the composition sum for
$k+1$.  All sums are finite because every $n_i\ge0$ and they sum to $N$.  If
$p_0=0$, every $n_i$ in a nonvanishing term is $\ge1$, so $N\ge k$.

For \eqref{plt:eq:mul-bell}, recall the defining generating identity of the
exponential partial Bell polynomials \cite{comtet-book,comtet1970}: in
$R[[z]]$ over a commutative $\RationalNumbers$-algebra $R$,
\[
 \frac1{k!}\Bigl(\sum_{j\ge1}x_j\frac{z^j}{j!}\Bigr)^{k}
 =\sum_{N\ge k}\ExponentialPartialBellPolynomial{N}{k}(x_1,x_2,\ldots)
   \frac{z^N}{N!}.
\]
Apply it with $R=\K[L]$ (a $\RationalNumbers$-algebra, as
$\operatorname{char}\K=0$), $z=t$ and $x_j=j!\,p_j$, so that
$\sum_{j\ge1}x_jz^j/j!=F$.  Comparing coefficients of $t^N$ gives
$\frac1{k!}[t^N]F^k=\frac1{N!}
\ExponentialPartialBellPolynomial{N}{k}(1!p_1,2!p_2,\ldots)$, which is
\eqref{plt:eq:mul-bell}; only the arguments $x_1,\ldots,x_{N-k+1}$ occur
because $\ExponentialPartialBellPolynomial{N}{k}$ is isobaric of weight $N$
in $k$ variables.
\end{proof}

\begin{theorem}[Binomial powers and their finite determination]
\label{plt:thm:mul-binomial}
Let $U\in\PLpow$ with $\ordt U\ge e$ for an integer $e\ge1$, write
$u_n=[t^n]U$, and let $\sigma\in\K$.  Put
$\binom{\sigma}{r}=\sigma(\sigma-1)\cdots(\sigma-r+1)/r!\in\K$, with
$\binom{\sigma}{0}=1$.  Then:
\begin{enumerate}
\item The family $\bigl(\binom{\sigma}{r}U^r\bigr)_{r\ge0}$ is summable, so
\begin{equation}\label{plt:eq:mul-binomial-def}
 (1+U)^\sigma\DefinitionEquals\sum_{r\ge0}\binom{\sigma}{r}U^r
\end{equation}
is a well-defined element of $\PLpow$ with constant block $1$.
\item \emph{(Finite determination.)}  For $N\ge0$,
\begin{equation}\label{plt:eq:mul-binomial-finite}
 [t^N](1+U)^\sigma
 =\sum_{r=0}^{\Floor{N/e}}\binom{\sigma}{r}\,[t^N]U^r,
\end{equation}
and this depends only on the blocks $u_e,u_{e+1},\ldots,u_N$, that is, only
on $\Trunc{U}{N+1}$.  The range of $r$ is exact: for $U=t^e$ with $e\mid N$
the term $r=N/e$ contributes $\binom{\sigma}{N/e}$, nonzero whenever
$\sigma\notin\SetOf{0,1,\ldots,N/e-1}$.  The top block is exact as well:
\begin{equation}\label{plt:eq:mul-binomial-top}
 [t^N](1+U)^\sigma=\sigma u_N+\Psi_N(u_e,\ldots,u_{N-e})
\end{equation}
for a universal polynomial $\Psi_N$ with coefficients in
$\RationalNumbers[\sigma]$, so for $\sigma\ne0$ the block $u_N$ cannot be
omitted.
\item \emph{(Exponent law.)}  For $\sigma,\tau\in\K$,
\begin{equation}\label{plt:eq:mul-binomial-group}
 (1+U)^\sigma(1+U)^\tau=(1+U)^{\sigma+\tau},
 \qquad (1+U)^0=1,
\end{equation}
and for $k\in\NaturalNumbers$ the series $(1+U)^k$ agrees with the $k$-fold
product.  Consequently $(1+U)^\sigma$ is a unit of $\PLpow$ with inverse
$(1+U)^{-\sigma}$, and $\sigma\mapsto(1+U)^\sigma$ is a homomorphism from
$(\K,+)$ to the units of $\PLpow$.
\item \emph{(Profile.)}  If $\dprofile{U}(n)\le\alpha n+\beta$ for all
$n\ge e$, with $\alpha\ge0$, then, writing $\beta^+=\max(\beta,0)$,
\begin{equation}\label{plt:eq:mul-binomial-profile}
 \dprofile{(1+U)^\sigma}(N)\le\Bigl(\alpha+\frac{\beta^+}{e}\Bigr)N
 \qquad(N\ge0).
\end{equation}
In particular $\mathfrak A\cap\PLpow$ is stable under
$U\mapsto(1+U)^\sigma$ whenever $\ordt U\ge1$.
\end{enumerate}
\end{theorem}

\begin{proof}
(1) By \eqref{plt:eq:mul-ord-additive}, $\ordt(U^r)\ge re$, which tends to
$+\infty$; so for each $N$ only the finitely many $r\le N/e$ can have
$\ordt(\binom{\sigma}{r}U^r)<N$.  \cref{plt:lem:mul-summable}(1) gives a
well-defined sum, and $\ordt U\ge e\ge1$ makes the $r=0$ term the only
contributor to $t^0$.

(2) Blockwise, $[t^N](1+U)^\sigma=\sum_{r\ge0}\binom{\sigma}{r}[t^N]U^r$, and
$[t^N]U^r=0$ once $re>N$; this is \eqref{plt:eq:mul-binomial-finite}.  By
\cref{plt:thm:mul-finite-determination} applied $r-1$ times, $[t^N]U^r$
depends only on the blocks $u_n$ with $e\le n\le N-(r-1)e$, and for $r\ge1$
this range is contained in $[e,N]$.  For the sharpness of the $r$-range take
$U=t^e$: then $U^r=t^{re}$ and the only contribution to $t^N$ comes from
$r=N/e$, with value $\binom{\sigma}{N/e}$, which vanishes exactly when
$\sigma$ is one of $0,1,\ldots,N/e-1$.  For
\eqref{plt:eq:mul-binomial-top}, the term $r=0$ contributes nothing to
$t^N$ for $N\ge1$, the term $r=1$ contributes $\sigma u_N$, and each term
with $r\ge2$ involves only $u_n$ with $n\le N-(r-1)e\le N-e$.

(3) By \cref{plt:lem:mul-summable}(2),
\[
 (1+U)^\sigma(1+U)^\tau
 =\sum_{r,s\ge0}\binom{\sigma}{r}\binom{\tau}{s}U^{r+s}
 =\sum_{m\ge0}\Bigl(\sum_{r+s=m}\binom{\sigma}{r}\binom{\tau}{s}\Bigr)U^m,
\]
the regrouping being legitimate because for each fixed $t$-block only
finitely many pairs $(r,s)$ contribute.  Vandermonde's convolution
\cite{graham-knuth-patashnik}
\[
 \sum_{r+s=m}\binom{\sigma}{r}\binom{\tau}{s}=\binom{\sigma+\tau}{m}
\]
holds for all $\sigma,\tau\in\K$: both sides are polynomial expressions in
$(\sigma,\tau)$ with coefficients in $\RationalNumbers$; they agree for all
pairs of nonnegative integers by comparing coefficients in
$(1+z)^{\sigma}(1+z)^{\tau}=(1+z)^{\sigma+\tau}$ in
$\RationalNumbers[z]$; a polynomial in two variables over the infinite domain
$\RationalNumbers$ vanishing on $\NaturalNumbers^2$ is the zero polynomial;
and $\RationalNumbers\subseteq\K$ because $\operatorname{char}\K=0$, so the
identity persists in $\K$.  This proves
\eqref{plt:eq:mul-binomial-group}; $(1+U)^0=1$ is immediate.  For
$k\in\NaturalNumbers$ we have $\binom{k}{r}=0$ for $r>k$, so
\eqref{plt:eq:mul-binomial-def} is the finite binomial expansion of the
$k$-fold product in the commutative ring $\PLpow$.  Invertibility follows by
taking $\tau=-\sigma$.

(4) By \eqref{plt:eq:mul-binomial-finite} it suffices to bound
$\degL[t^N]U^r$ for $1\le r\le\Floor{N/e}$ (the term $r=0$ contributes only
at $N=0$, where the bound reads $0\le0$).  By
\eqref{plt:eq:mul-power-coefficient} and \eqref{plt:eq:mul-degL-laws},
\[
 \degL[t^N]U^r
 \le\max_{\substack{n_1+\cdots+n_r=N\\ n_i\ge e}}
    \sum_{i=1}^{r}(\alpha n_i+\beta)
 =\alpha N+r\beta
 \le\alpha N+\Floor{N/e}\beta^+
 \le\Bigl(\alpha+\frac{\beta^+}{e}\Bigr)N,
\]
where the second inequality uses $r\le\Floor{N/e}$ when $\beta\ge0$ and
$r\beta\le0\le\Floor{N/e}\beta^+$ when $\beta<0$.  Taking the maximum over
$r$ gives \eqref{plt:eq:mul-binomial-profile}.
\end{proof}

\begin{corollary}[Formal roots]
\label{plt:cor:mul-roots}
With $U$ as above and $p,q\in\IntegerNumbers$, $q\ge1$,
\[
 \bigl((1+U)^{p/q}\bigr)^{q}=(1+U)^{p},
\]
the right-hand side being the ordinary $p$-th power (or the reciprocal of the
$|p|$-th power if $p<0$).  In particular $(1+U)^{1/q}$ is a $q$-th root of
$1+U$ in $\PLpow$, and it is the unique one with constant block $1$.
\end{corollary}

\begin{proof}
Iterating \eqref{plt:eq:mul-binomial-group} $q-1$ times gives
$((1+U)^{p/q})^q=(1+U)^{qp/q}=(1+U)^p$, and
\cref{plt:thm:mul-binomial}(3) identifies the right-hand side with the
ordinary power.  For uniqueness, suppose $V,W\in\PLpow$ have constant block
$1$ and $V^q=W^q=1+U$.  Then $(V-W)\sum_{i=0}^{q-1}V^iW^{q-1-i}=0$; the
second factor has constant block $q\ne0$ (characteristic zero), hence is
nonzero, and $\PLpow$ is a domain by \cref{plt:cor:mul-domain}, so $V=W$.
\end{proof}

\begin{remark}[Infinite products]
\label{plt:rmk:mul-infinite-products}
If $(U_r)_{r\ge1}$ lies in $\PLpow$ with $\ordt U_r\to+\infty$, the partial
products $P_R=\prod_{r\le R}(1+U_r)$ converge $t$-adically.  Indeed, given
$N$ choose $R_N$ with $\ordt U_r\ge N$ for $r>R_N$; for $R>R_N$,
\[
 P_R-P_{R_N}=P_{R_N}\Bigl(\prod_{R_N<r\le R}(1+U_r)-1\Bigr),
\]
and expanding the finite product shows that the bracket is a sum of products
of at least one $U_r$ with $r>R_N$, hence of $\ordt\ge N$; since
$\ordt P_{R_N}\ge0$ we get $\ordt(P_R-P_{R_N})\ge N$.  Thus $(P_R)$ is
$t$-adically Cauchy in the complete ring $\PLpow$ and
$\prod_{r\ge1}(1+U_r)$ is well defined.  On $\PLlau$ the same argument
applies once one notes that $\ordt U_r\to+\infty$ forces all but finitely
many factors to lie in $\PLpow$, so the partial products have a common lower
support bound.
\end{remark}

\subsection{Three worked computations}
\label{plt:sec:mul-worked}

\begin{example}[A Laurent product through four blocks, with degree
bookkeeping]
\label{plt:ex:mul-worked-laurent}
Over $\K=\RationalNumbers$ take
\begin{align*}
 F&=3t^{-2}+(L-1)t^{-1}+(L^2+2)+2L\,t+(L^2-L)t^2
   +\FormalBigOAt{t^{3}}{t},\\
 G&=t^{-1}+(2-L)+(L+1)t+L^2t^2+\FormalBigOAt{t^{3}}{t}.
\end{align*}
Here $\ordt F=-2$ and $\ordt G=-1$; $F$ is known through $\Trunc{F}{3}$ and
$G$ through $\Trunc{G}{3}$.  By the precision rule the product is known
through
\[
 \Trunc{FG}{\min(3+(-1),\,3+(-2))}=\Trunc{FG}{1},
\]
that is, through the block $t^{0}$ inclusive --- four blocks in all, starting
at $t^{-3}$.  The extra block $[t^2]F$ is wasted: it can influence only
$[t^{1}](FG)$ and beyond, for which $[t^{3}]G$ would also be needed.

Writing $A_n=[t^n]F$ and $B_m=[t^m]G$, the convolution
\eqref{plt:eq:mul-cauchy} gives
\begin{align*}
 C_{-3}&=A_{-2}B_{-1}=3,\\
 C_{-2}&=A_{-2}B_{0}+A_{-1}B_{-1}
        =3(2-L)+(L-1)=5-2L,\\
 C_{-1}&=A_{-2}B_{1}+A_{-1}B_{0}+A_{0}B_{-1}\\
       &=3(L+1)+(-L^2+3L-2)+(L^2+2)=6L+3,\\
 C_{0}&=A_{-2}B_{2}+A_{-1}B_{1}+A_{0}B_{0}+A_{1}B_{-1}\\
      &=3L^2+(L^2-1)+(-L^3+2L^2-2L+4)+2L
      =-L^3+6L^2+3 .
\end{align*}
Hence
\begin{equation}\label{plt:eq:mul-worked-laurent}
 FG=3t^{-3}+(5-2L)t^{-2}+(6L+3)t^{-1}+(-L^3+6L^2+3)
   +\FormalBigOAt{t^{1}}{t}.
\end{equation}

Leading data: $\lm(F)=t^{-2}$ with $\lc(F)=3$, $\lm(G)=t^{-1}$ with
$\lc(G)=1$; \cref{plt:thm:mul-leading-law} predicts $\lm(FG)=t^{-3}$ and
$\lc(FG)=3$, confirmed by $C_{-3}=3$.

The degree bookkeeping is instructive.  The profiles are
$\dprofile{F}(-2)=0$, $\dprofile{F}(-1)=1$, $\dprofile{F}(0)=2$,
$\dprofile{F}(1)=1$ and $\dprofile{G}(-1)=0$, $\dprofile{G}(0)=1$,
$\dprofile{G}(1)=1$, $\dprofile{G}(2)=2$.  Then:

\begin{center}
\begin{tabular}{@{}ccccc@{}}
\hline
$N$ & maximizing pairs $P_N$ & bound $M$ & top coefficient
 \eqref{plt:eq:mul-top-coefficient} & $\dprofile{FG}(N)$ \\
\hline
$-3$ & $(-2,-1)$ & $0$ & $3\cdot1=3$ & $0$ \\
$-2$ & $(-2,0),\ (-1,-1)$ & $1$ & $3(-1)+1\cdot1=-2$ & $1$ \\
$-1$ & $(-1,0),\ (0,-1)$ & $2$ & $1\cdot(-1)+1\cdot1=0$ & $1<2$ \\
$0$  & $(0,0)$ & $3$ & $1\cdot(-1)=-1$ & $3$ \\
\hline
\end{tabular}
\end{center}

At $N=-2$ two pairs compete and their contributions $-3$ and $+1$ do not
cancel, so the bound is attained, with top coefficient $-2$: indeed
$C_{-2}=-2L+5$.  At $N=-1$ two pairs compete and cancel exactly, so the
degree drops from the bound $2$ to the actual value $1$: indeed
$C_{-1}=6L+3$.  At $N=0$ the maximizing pair is unique, so
\cref{plt:cor:mul-degree-equality}(1) forces attainment, with top
coefficient $-1$: indeed $C_0=-L^3+\cdots$.  This is exactly the phenomenon
\cref{plt:thm:mul-degree-bound} isolates; no inspection of the finished
polynomial is needed to predict it.
\end{example}

\begin{example}[A square root through four blocks, verified by squaring]
\label{plt:ex:mul-worked-sqrt}
Let
\[
 U=2Lt+(L^2+2)t^2+4L\,t^3+\FormalBigOAt{t^{4}}{t}\in t\PLpow,
\]
so $e=1$ in \cref{plt:thm:mul-binomial}, and take $\sigma=\tfrac12$, with
\[
 \binom{1/2}{0}=1,\quad
 \binom{1/2}{1}=\tfrac12,\quad
 \binom{1/2}{2}=-\tfrac18,\quad
 \binom{1/2}{3}=\tfrac1{16}.
\]
By \eqref{plt:eq:mul-binomial-finite} only $r\le N$ contributes to $t^N$, so
through $t^3$ we need
\[
 U^2=4L^2t^2+(4L^3+8L)t^3+\FormalBigOAt{t^{4}}{t},
 \qquad
 U^3=8L^3t^3+\FormalBigOAt{t^{4}}{t}.
\]
Block by block:
\begin{align*}
 [t^1](1+U)^{1/2}&=\tfrac12(2L)=L,\\
 [t^2](1+U)^{1/2}&=\tfrac12(L^2+2)-\tfrac18(4L^2)
   =\tfrac{L^2}2+1-\tfrac{L^2}2=1,\\
 [t^3](1+U)^{1/2}&=\tfrac12(4L)-\tfrac18(4L^3+8L)+\tfrac1{16}(8L^3)
   =2L-\tfrac{L^3}2-L+\tfrac{L^3}2=L .
\end{align*}
Hence
\begin{equation}\label{plt:eq:mul-worked-sqrt}
 (1+U)^{1/2}=1+L\,t+t^2+L\,t^3+\FormalBigOAt{t^{4}}{t}.
\end{equation}
Two logarithmic cancellations occur, one at each of $t^2$ and $t^3$; the
degree bound \eqref{plt:eq:mul-binomial-profile} with $\alpha=1$,
$\beta^{+}=0$, $e=1$ predicts only $\dprofile{}(N)\le N$, and the actual
profile $0,1,0,1$ is strictly smaller at $N=2,3$.  A degree bound is an upper
bound; it is not a prediction.

\emph{Verification.}  \cref{plt:cor:mul-roots} says the square of
\eqref{plt:eq:mul-worked-sqrt} must reproduce $1+U$ through $t^3$.  Squaring
by \eqref{plt:eq:mul-cauchy}:
\[
 [t^0]=1,\quad
 [t^1]=2L,\quad
 [t^2]=2\cdot1+L^2=L^2+2,\quad
 [t^3]=2L+2L=4L,
\]
which are exactly the blocks $1,u_1,u_2,u_3$ of $1+U$.  Note the precision
accounting: squaring a series known through $\Trunc{\cdot}{4}$ with
$\ordt=0$ returns a result known through $\Trunc{\cdot}{4}$, so the check is
available at every block computed and at none beyond.
\end{example}

\begin{example}[A product in the iterated field]
\label{plt:ex:mul-worked-iterated}
In $\PLit$ the coefficient convolution is finite at the inner level too.
Take the two blocks
\[
 f=\frac1{L+1}=L^{-1}-L^{-2}+L^{-3}-\cdots,
 \qquad
 g=\frac1{L-1}=L^{-1}+L^{-2}+L^{-3}+\cdots
\]
in $\K((L^{-1}))$.  For $r\ge2$, the coefficient of $L^{-r}$ in $fg$ is a sum
over the finitely many decompositions $r=i+j$ with $i,j\ge1$:
\[
 \kappa_{-r}(fg)=\sum_{i=1}^{r-1}(-1)^{i+1}\cdot1
 =\begin{cases}1,&r\ \text{even},\\0,&r\ \text{odd},\end{cases}
\]
so $fg=L^{-2}+L^{-4}+L^{-6}+\cdots=(L^2-1)^{-1}$, as it must be.  The
degree law is visible: $\degL f=\degL g=-1$, $\degL(fg)=-2$, and
$\lc_L(f)=\lc_L(g)=\lc_L(fg)=1$.  Consequently, for
$F=f\,t^{-1}+\cdots$ and $G=g+\cdots$ in $\PLit$ one gets
$\lm(FG)=t^{-1}L^{-2}$ and $\lc(FG)=1$ by
\cref{plt:thm:mul-leading-law}, with no need to compute any further block.
Observe that $\lm$ here carries a \emph{negative} power of $L$: this is
legitimate in $\PLrat$ and $\PLit$ and impossible in $\PLlau$, which is
exactly the content of \cref{plt:cor:mul-domain}.
\end{example}

\subsection{The analytic product rule}
\label{plt:sec:mul-analytic}

Everything above is formal algebra.  We now record the exact circumstances
under which the block convolution computes the asymptotic expansion of an
actual product of functions.  Nothing in this subsection asserts
convergence: an asymptotic expansion is a statement about remainders, and the
formal product theorem is used as an input, not replaced.

\begin{definition}[Polynomial--logarithmic asymptotic expansion]
\label{plt:def:mul-asymptotic}
Let $\K\subseteq\ComplexNumbers$, let $X_0>1$, and let $S=[X_0,+\infty)$ with
the real branch $L=\log X$.  A function $f\colon S\to\ComplexNumbers$ is
\emph{PL-asymptotic} to $F\in\PLlau$, written $f\approx F$, if for every
$N\in\IntegerNumbers$ there exist $M_N\in\NaturalNumbers$, $C_N>0$ and
$X_N\ge X_0$ such that
\begin{equation}\label{plt:eq:mul-asymptotic}
 \AbsoluteValue{f(X)-\bigl(\Trunc{F}{N}\bigr)(X)}
 \le C_N\,X^{-N}(\log X)^{M_N}
 \qquad(X\ge X_N),
\end{equation}
where $\bigl(\Trunc{F}{N}\bigr)(X)=\sum_{\ordt F\le n<N}p_n(\log X)X^{-n}$ is
the actual finite sum obtained by substituting the generators.  Equivalently,
$f-\Trunc{F}{N}=\BigOAt{X^{-N}(\log X)^{M_N}}{X\to+\infty,\ X\in S}$.
\end{definition}

\begin{lemma}[Basic properties of $\approx$]
\label{plt:lem:mul-asymptotic-basic}
Let $f,g\colon S\to\ComplexNumbers$ and $F,G\in\PLlau$.
\begin{enumerate}
\item If \eqref{plt:eq:mul-asymptotic} holds for arbitrarily large $N$, it
holds for every $N$; so it suffices to verify it along a cofinal set of
cutoffs.
\item $f\approx F$ and $f\approx F'$ imply $F=F'$.
\item If $f\approx F$ and $g\approx G$, then $f+g\approx F+G$ and
$\lambda f\approx\lambda F$ for $\lambda\in\K$.
\item $f\approx0$ if and only if $f=\BigOAt{X^{-N}}{X\to+\infty}$ for every
$N$.
\end{enumerate}
\end{lemma}

\begin{proof}
(1) Let $N'<N$ and suppose \eqref{plt:eq:mul-asymptotic} holds at $N$.  Then
\[
 f-\Trunc{F}{N'}=\bigl(f-\Trunc{F}{N}\bigr)
   +\sum_{N'\le n<N}p_n(\log X)X^{-n},
\]
a finite sum; the first term is $\BigOAt{X^{-N}(\log
X)^{M_N}}{X\to+\infty}\subseteq\BigOAt{X^{-N'}(\log X)^{M_N}}{X\to+\infty}$
since $N\ge N'$, and each term of the sum is
$\BigOAt{X^{-N'}(\log X)^{\degL p_n}}{X\to+\infty}$.

(2) Suppose $D=F-F'\ne0$, with $n_1=\ordt D$ and leading block $p$ of degree
$d$ and leading coefficient $c\ne0$.  Take $N=n_1+1$ and subtract the two
instances of \eqref{plt:eq:mul-asymptotic}:
\[
 \bigl(\Trunc{D}{n_1+1}\bigr)(X)=p(\log X)X^{-n_1}
 =\BigOAt{X^{-n_1-1}(\log X)^{M}}{X\to+\infty}
\]
for some $M$.  But $|p(\log X)|\sim|c|(\log X)^{d}$, so the left side
divided by $X^{-n_1-1}(\log X)^M$ is $\sim|c|X(\log X)^{d-M}\to+\infty$, a
contradiction.

(3) Immediate from
$\Trunc{F+G}{N}=\Trunc{F}{N}+\Trunc{G}{N}$ and the triangle inequality.

(4) If $f\approx0$ then $\Trunc{0}{N}=0$ and $f=\BigOAt{X^{-N}(\log
X)^{M_N}}{X\to+\infty}$ for every $N$; since $(\log
X)^{M_N}=\LittleOAt{X}{X\to+\infty}$, this gives
$f=\BigOAt{X^{-N+1}}{X\to+\infty}$ for every $N$, hence
$f=\BigOAt{X^{-N}}{X\to+\infty}$ for every $N$.  The converse is trivial.
\end{proof}

\begin{theorem}[Analytic product rule]
\label{plt:thm:mul-analytic-product}
Let $\K\subseteq\ComplexNumbers$, $S=[X_0,+\infty)$ with the real logarithm,
and let $f,g\colon S\to\ComplexNumbers$ with $f\approx F$ and $g\approx G$
for $F,G\in\PLlau$.  Then
\[
 fg\approx FG,
\]
where $FG$ is the formal Cauchy product \eqref{plt:eq:mul-cauchy}.
Consequently the set
$\mathcal A(S)=\SetOf{f:\exists F\in\PLlau,\ f\approx F}$ is a
$\K$-algebra, the map $f\mapsto F$ is a well-defined $\K$-algebra
homomorphism $\mathcal A(S)\to\PLlau$, and its kernel is the ideal of
functions that are $\BigOAt{X^{-N}}{X\to+\infty}$ for every $N$.
\end{theorem}

\begin{proof}
If $F=0$, then by \cref{plt:lem:mul-asymptotic-basic}(4) $f$ is
$\BigOAt{X^{-P}}{X\to+\infty}$ for every $P$; taking $N=\ordt G+1$ in
\eqref{plt:eq:mul-asymptotic} for $g$ shows $g=\BigOAt{X^{-\ordt G}(\log
X)^{M}}{X\to+\infty}$ for some $M$, so $fg=\BigOAt{X^{-P}}{X\to+\infty}$ for
every $P$ and $fg\approx0=FG$.  Assume now $F,G\ne0$ and set $a=\ordt F$,
$b=\ordt G$.

Fix $N\in\IntegerNumbers$ and put
\[
 P=\max(N-b,\;a),
 \qquad
 Q=\max(N-a,\;b).
\]
Write $f_P=(\Trunc{F}{P})(X)$, $r=f-f_P$, $g_Q=(\Trunc{G}{Q})(X)$,
$s=g-g_Q$.  By hypothesis
$r=\BigOAt{X^{-P}(\log X)^{m_1}}{X\to+\infty}$ and
$s=\BigOAt{X^{-Q}(\log X)^{m_2}}{X\to+\infty}$ for suitable exponents.  Since
$\Trunc{F}{P}$ is a finite sum of blocks with outer exponents $\ge a$,
\[
 f_P=\BigOAt{X^{-a}(\log X)^{m_3}}{X\to+\infty},
 \qquad
 g_Q=\BigOAt{X^{-b}(\log X)^{m_4}}{X\to+\infty}
\]
with $m_3=\max_{a\le n<P}\degL p_n$ and similarly for $m_4$ (the estimates
being vacuously true if the sums are empty).  Therefore
\[
 fg-f_Pg_Q=f_Ps+rg_Q+rs
\]
is $\BigOAt{X^{-(a+Q)}(\log X)^{\bullet}}{X\to+\infty}
+\BigOAt{X^{-(P+b)}(\log X)^{\bullet}}{X\to+\infty}
+\BigOAt{X^{-(P+Q)}(\log X)^{\bullet}}{X\to+\infty}$.
Now $a+Q\ge a+(N-a)=N$, $P+b\ge(N-b)+b=N$, and $P+Q\ge a+(N-a)=N$; hence
\begin{equation}\label{plt:eq:mul-analytic-step1}
 fg-f_Pg_Q=\BigOAt{X^{-N}(\log X)^{M'}}{X\to+\infty}
\end{equation}
for some $M'$.

Next, $f_Pg_Q$ is the value at $X$ of the finite polynomial--logarithmic
expression $\bigl(\Trunc{F}{P}\bigr)\bigl(\Trunc{G}{Q}\bigr)\in\PLlau$,
because substitution of the generators $t=X^{-1}$, $L=\log X$ into a finite
sum of monomials is a ring homomorphism into functions on $S$.  Since
$P\ge N-b$ and $Q\ge N-a$, \cref{plt:cor:mul-truncated-product} gives
\[
 \Trunc{\bigl(\Trunc{F}{P}\bigr)\bigl(\Trunc{G}{Q}\bigr)}{N}
 =\Trunc{FG}{N}.
\]
Hence $\bigl(\Trunc{F}{P}\bigr)\bigl(\Trunc{G}{Q}\bigr)-\Trunc{FG}{N}$ is a
\emph{finite} sum of monomials $c\,t^nL^j$ with $n\ge N$, so its value at
$X$ is $\BigOAt{X^{-N}(\log X)^{M''}}{X\to+\infty}$ with $M''$ the largest
$j$ occurring.  Combining with \eqref{plt:eq:mul-analytic-step1},
\[
 fg-\bigl(\Trunc{FG}{N}\bigr)(X)
 =\BigOAt{X^{-N}(\log X)^{\max(M',M'')}}{X\to+\infty},
\]
which is \eqref{plt:eq:mul-asymptotic} for $fg$ and $FG$.  As $N$ was
arbitrary, $fg\approx FG$.

The algebra statements now follow: $\mathcal A(S)$ is closed under $+$,
scalars and $\cdot$ by \cref{plt:lem:mul-asymptotic-basic}(3) and the above;
the assignment $f\mapsto F$ is well defined by
\cref{plt:lem:mul-asymptotic-basic}(2) and is a homomorphism by the same two
statements; its kernel is described by
\cref{plt:lem:mul-asymptotic-basic}(4).
\end{proof}

\begin{remark}[Which hypotheses are actually used, and the sectorial version]
\label{plt:rmk:mul-analytic-hypotheses}
The proof uses only three things: that $f$ and $g$ have remainders bounded by
$X^{-N}$ times \emph{some} power of $\log X$; that finitely many blocks are
involved at each stage; and the formal finite-determination theorem
\cref{plt:thm:mul-finite-determination}, which is what allows the two
truncation cutoffs $P$ and $Q$ to be chosen independently of the shape of
$F$ and $G$.  No convergence and no uniformity in $L$ is used or obtained.

The same proof works verbatim on an unbounded set
$S\subseteq\SetOf{z:\AbsoluteValue{\arg z}\le\theta}$ with
$0<\theta<\pi$ and the principal logarithm, with $X^{-N}$ replaced by
$\AbsoluteValue{X}^{-N}$ and $(\log X)^{M}$ by
$(\log\AbsoluteValue{X})^{M}$; the only additional input is
$\AbsoluteValue{\log X}\le\log\AbsoluteValue{X}+\theta
\le2\log\AbsoluteValue{X}$ for $\AbsoluteValue{X}\ge\EulerE^{\theta}$, so
that a bound in $\AbsoluteValue{\log X}$ and a bound in
$\log\AbsoluteValue{X}$ differ only by a constant factor.  Uniformity claims
on a sector must be stated for a \emph{closed} subsector; the branch of
$\log$ is part of the statement and cannot be recovered from the formal
data.
\end{remark}

\begin{warningbox}
Three cautions, each corresponding to an overstatement found in the sources.
\begin{enumerate}
\item \emph{A formal identity is not an analytic theorem.}
\cref{plt:thm:mul-analytic-product} presupposes that $f$ and $g$
\emph{already have} expansions; it does not produce one.  The
homomorphism $f\mapsto F$ has a large kernel --- $\EulerE^{-X}\approx0$ ---
so a formal computation determines an actual function at best up to a flat
error.
\item \emph{The formal $\FormalBigOAt{t^N}{t}$ carries no uniformity in $L$.}
Asserting $F=G+\FormalBigOAt{t^N}{t}$ says only that the blocks agree below
outer order $N$.  In particular the constant $C_N$ and the exponent $M_N$ of
\eqref{plt:eq:mul-asymptotic} are not supplied by any formal statement and
must be produced separately.
\item \emph{The three $\BigO$-roles are distinct.}  The formal tail
$\FormalBigOAt{t^N}{t}$, the analytic estimate
$\BigOAt{X^{-N}(\log X)^{M}}{X\to+\infty}$ of
\eqref{plt:eq:mul-asymptotic}, and the complexity bound
$\BigO\bigl((N-a-b)^2d^2\bigr)$ of \cref{plt:rmk:mul-cost} are three
different assertions and none follows from the others.  A bare $\BigO$ for an
asymptotic claim, with the regime left to context, is never acceptable.
\end{enumerate}
\end{warningbox}

\begin{summarybox}
Multiplication is Cauchy convolution in $t=X^{-1}$ with exact coefficient
arithmetic in $L=\log X$.  Every output block below a cutoff is a finite
bilinear expression in finitely many input blocks, with the exact range
$a\le n\le N-b$, $b\le m\le N-a$; the resulting precision rule is
$\Trunc{FG}{\min(P+b,\,Q+a)}$ and it cannot be improved.  All four ambient
classes are closed and are integral domains, the leading monomial and leading
scalar are multiplicative, and the rank-two valuation
$v(F)=(\ordt F,-\degL p_{\ordt F})$ is additive.  The log-degree profile
obeys the max-plus bound $\dprofile{FG}\le\dprofile{F}\boxplus\dprofile{G}$,
with equality exactly when the maximizing pairs do not cancel; affine
profiles are stable under products, powers, and binomial powers, but not
under $t$-adic limits.  Binomial powers $(1+U)^\sigma$ are defined for every
$\sigma\in\K$ by a summable family, are finitely determined at every order,
and satisfy the exponent law.  Under printed remainder hypotheses on a stated
domain and branch, the formal product computes the asymptotic expansion of
the analytic product --- and nothing more.
\end{summarybox}

% ---- fragment 04-division ----
\section{Units, reciprocals, and division}
\label{plt:sec:division}

Addition and multiplication never leave \(\PLlau=\K[L]((t))\).  Division is the
first operation that does.  This section settles exactly when it does not --
the unit criterion -- and then follows the three escape routes that the
criterion forces: an accidental exact quotient inside \(\PLlau\), the
rational-log field \(\PLrat=\K(L)((t))\), and the iterated Laurent-log field
\(\PLit=\K((L^{-1}))((t))\).  Each of the three is a genuinely different
object, and conflating them is the most frequent error in informal
polynomial--logarithmic computation.

Throughout, \(\K\) is a field of characteristic zero, \(t\DefinitionEquals
X^{-1}\) and \(L\DefinitionEquals\log X\) are independent formal generators, and
a nonzero element of \(\PLlau\) is written
\begin{equation}\label{plt:eq:div-normal-form}
  F=\sum_{n\ge n_{0}}p_{n}(L)t^{n},
  \qquad p_{n}\in\K[L],\quad p_{n_{0}}\ne0,
\end{equation}
so that \(\ordt F=n_{0}\).  We call \(p_{\ordt F}\) the \emph{leading block} of
\(F\); it is a polynomial, not a scalar, and it must be distinguished from the
leading monomial \(\lm F=t^{n_{0}}L^{d}\) and from the scalar leading
coefficient \(\lc F\in\K^{\times}\), where \(d=\degL p_{n_{0}}\).  Truncation is
always the exclusive operator \(\Trunc{F}{N}=\sum_{n<N}p_{n}(L)t^{n}\), and
\(F=G+\FormalBigOAt{t^{N}}{t}\) abbreviates \(\ordt(F-G)\ge N\); it carries no
analytic content whatever.

\subsection{The exact unit criterion}
\label{plt:sec:div-units}

Two elementary facts do all the work.  Both are standard, but the second is the
step that three of the six sources compress into a single clause, and it is
precisely the step that decides the nonunit direction of the criterion.

\begin{lemma}[Units of the coefficient ring]\label{plt:lem:div-poly-units}
Let \(\K\) be a field.  Then \(\K[L]^{\times}=\K^{\times}\): a polynomial is
invertible in \(\K[L]\) if and only if it is a nonzero constant.
\end{lemma}

\begin{proof}
A nonzero constant \(c\) has inverse \(c^{-1}\).  Conversely, suppose
\(pq=1\) with \(p,q\in\K[L]\).  Since \(\K\) is a field, \(\K[L]\) is an
integral domain and \(\degL(pq)=\degL p+\degL q\).  Hence
\(\degL p+\degL q=\degL 1=0\), and as both degrees are nonnegative integers,
\(\degL p=\degL q=0\).  So \(p\in\K\), and \(p\ne0\) because \(pq=1\).
\end{proof}

\begin{lemma}[The Laurent ring is a domain and \(\ordt\) is additive]
\label{plt:lem:div-valuation}
Let \(F,G\in\PLlau\) be nonzero, with leading blocks \(p_{\ordt F}\) and
\(q_{\ordt G}\).  Then \(FG\ne0\),
\begin{equation}\label{plt:eq:div-valuation-additive}
  \ordt(FG)=\ordt F+\ordt G,
\end{equation}
and the leading block of \(FG\) is \(p_{\ordt F}\,q_{\ordt G}\).  In particular
\(\PLlau\) is an integral domain, and the same statements hold in \(\PLpow\),
in \(\K(L)((t))\), and in \(E((t))\) for any integral domain \(E\).
\end{lemma}

\begin{proof}
Write \(m=\ordt F\), \(r=\ordt G\).  The Cauchy convolution gives
\([t^{k}](FG)=\sum_{i+j=k}p_{i}q_{j}\), a finite sum because \(i\ge m\) and
\(j\ge r\) force \(m\le i\le k-r\).  For \(k<m+r\) every term has \(i<m\) or
\(j<r\) and therefore vanishes.  For \(k=m+r\) the only surviving pair is
\((i,j)=(m,r)\), so \([t^{m+r}](FG)=p_{m}q_{r}\), which is nonzero because
\(\K[L]\) is an integral domain.  Hence \(FG\ne0\) and both assertions hold.
\end{proof}

\begin{theorem}[Unit criterion]\label{plt:thm:div-unit-criterion}
Let \(F\in\PLlau\) be nonzero, written as in \cref{plt:eq:div-normal-form}.
The following are equivalent.
\begin{enumerate}
\item \(F\) is a unit of \(\PLlau\).
\item The leading block \(p_{n_{0}}\) is a nonzero constant, that is
      \(p_{n_{0}}\in\K^{\times}\).
\item \(F=c\,t^{n_{0}}(1+U)\) for some \(c\in\K^{\times}\) and some
      \(U\in t\PLpow\).
\end{enumerate}
For \(F\in\PLpow\) with \(n_{0}=0\), these are equivalent to \(p_{0}\in
\K^{\times}\), and \(F^{-1}\) then lies in \(\PLpow\).
\end{theorem}

\begin{proof}
\((2)\Leftrightarrow(3)\).  If \(p_{n_{0}}=c\in\K^{\times}\), put
\(U\DefinitionEquals c^{-1}t^{-n_{0}}F-1\); every block of \(c^{-1}t^{-n_{0}}F\)
of index \(\ge1\) is a polynomial and the index-\(0\) block is \(1\), so
\(U\in t\PLpow\).  Conversely, in \(F=ct^{n_{0}}(1+U)\) with \(U\in t\PLpow\)
the block of index \(n_{0}\) is \(c\).

\((3)\Rightarrow(1)\).  Set
\begin{equation}\label{plt:eq:div-geometric}
  G\DefinitionEquals c^{-1}t^{-n_{0}}\sum_{k\ge0}(-U)^{k}.
\end{equation}
Because \(\ordt(U^{k})\ge k\), for each \(N\) only the terms with \(k<N\)
contribute to blocks of index \(<N\); the family is therefore formally summable
and \(G\in\PLlau\), with polynomial blocks because \(\PLpow\) is closed under
multiplication.  For every \(N\ge0\),
\[
  \Trunc{\Bigl((1+U)\sum_{k\ge0}(-U)^{k}\Bigr)}{N}
  =\Trunc{\Bigl((1+U)\sum_{k<N}(-U)^{k}\Bigr)}{N}
  =\Trunc{\bigl(1-(-U)^{N}\bigr)}{N}
  =\Trunc{1}{N},
\]
using the telescoping identity \((1+U)\sum_{k<N}(-U)^{k}=1-(-U)^{N}\) and
\(\ordt((-U)^{N})\ge N\).  Since two elements of \(\PLlau\) with equal exclusive
truncations at every order are equal, \((1+U)\sum_{k\ge0}(-U)^{k}=1\), whence
\(FG=1\).

% ed.: S5 and S6 assert the nonunit direction by citing the general
% power-series unit criterion; the additivity step that makes it valid on the
% Laurent ring is supplied here explicitly.
\((1)\Rightarrow(2)\).  Suppose \(FG=1\) with \(G\in\PLlau\).  Then \(G\ne0\),
and \cref{plt:lem:div-valuation} gives \(\ordt F+\ordt G=\ordt 1=0\) together
with \(p_{\ordt F}\,q_{\ordt G}=1\) in \(\K[L]\), where \(q_{\ordt G}\) is the
leading block of \(G\).  By \cref{plt:lem:div-poly-units},
\(p_{\ordt F}\in\K^{\times}\).

The statement for \(\PLpow\) is the case \(n_{0}=0\) of the above, together
with the observation that the \(G\) constructed in \cref{plt:eq:div-geometric}
lies in \(\PLpow\) when \(n_{0}=0\).
\end{proof}

\begin{corollary}[The unit group]\label{plt:cor:div-unit-group}
The group of units of \(\PLlau\) is
\begin{equation}\label{plt:eq:div-unit-group}
  \PLlau^{\times}
  =\SetOf{c\,t^{m}(1+U):m\in\IntegerNumbers,\ c\in\K^{\times},\
          U\in t\PLpow},
\end{equation}
and the factorisation is unique, so that
\(\PLlau^{\times}\cong\IntegerNumbers\times\K^{\times}\times(1+t\PLpow)\) as
abelian groups.  Moreover \(\PLpow^{\times}=\K^{\times}\cdot(1+t\PLpow)\).
\end{corollary}

\begin{proof}
The description of the set is \cref{plt:thm:div-unit-criterion}.  Uniqueness:
in \(F=ct^{m}(1+U)\) one has \(m=\ordt F\) by
\cref{plt:lem:div-valuation}, then \(c\) is the leading block, then
\(U=c^{-1}t^{-m}F-1\).  The three subgroups \(t^{\IntegerNumbers}\),
\(\K^{\times}\) and \(1+t\PLpow\) generate \(\PLlau^{\times}\) and pairwise
intersect trivially by the uniqueness just proved, so the multiplication map
from their product is a group isomorphism.
\end{proof}

\begin{example}[A nonzero nonunit]\label{plt:ex:div-nonunit}
The element \(L+1\in\PLpow\) is nonzero, but it is not a unit of \(\PLlau\).
Indeed \(\ordt(L+1)=0\) and its leading block is \(L+1\), which is not
constant; \cref{plt:thm:div-unit-criterion} applies.  Unwinding the proof gives
the direct argument: if \((L+1)G=1\) with \(G=\sum_{n\ge n_{1}}q_{n}(L)t^{n}\)
nonzero, then \(\ordt G=0\) and \((L+1)q_{0}=1\) in \(\K[L]\), which is
impossible because \(\degL\bigl((L+1)q_{0}\bigr)=1+\degL q_{0}\ge1\).  The
same argument shows that \(L\), \(L+t\), \(L^{2}+t^{3}\) and
\(t^{-5}(L-7)\) are nonzero nonunits.  Their reciprocals exist in \(\PLrat\):
\begin{equation}\label{plt:eq:div-Lplust}
  \frac{1}{L+t}
  =\frac1L-\frac{t}{L^{2}}+\frac{t^{2}}{L^{3}}-\cdots
  \in\PLrat\setminus\PLlau .
\end{equation}
\end{example}

\begin{remark}[The dominant-monomial fallacy]\label{plt:rmk:div-dominant}
Factoring out the dominant monomial is \emph{not} enough to invert.  Writing
\(L+t=L(1+t/L)\) and expanding the second factor geometrically already commits
one to the coefficient ring \(\K(L)\) or \(\K((L^{-1}))\): the step that leaves
\(\PLlau\) is the inversion of the scalar-free factor \(L\), not the infinite
\(t\)-expansion.  \Cref{plt:eq:div-Lplust} makes the point in one line -- the
failure is visible already in the block of index \(0\).
\end{remark}

\subsection{Existence, uniqueness, and the reciprocal recurrence}
\label{plt:sec:div-recurrence}

\Cref{plt:eq:div-geometric} proves existence but is a poor algorithm.  The
practical construction is triangular.

\begin{theorem}[Reciprocal of a unit]\label{plt:thm:div-reciprocal}
Let \(B=\sum_{n\ge0}B_{n}(L)t^{n}\in\PLpow\) with \(B_{0}=c\in\K^{\times}\).
Then there is exactly one \(C\in\PLpow\) with \(BC=1\), namely
\(C=\recip{B}=\sum_{n\ge0}C_{n}(L)t^{n}\) with
\begin{align}
  C_{0}&=c^{-1},\label{plt:eq:div-recip-c0}\\
  C_{n}&=-c^{-1}\sum_{j=1}^{n}B_{j}\,C_{n-j}\qquad(n\ge1),
        \label{plt:eq:div-recip-cn}
\end{align}
and every \(C_{n}\) lies in \(\K[L]\).  Moreover \(C\) is the unique reciprocal
of \(B\) in \(\PLlau\), and for a Laurent unit \(F=t^{m}B\) one has
\(F^{-1}=t^{-m}C\).
\end{theorem}

\begin{proof}
\emph{Existence.}  Define \(C_{n}\in\K[L]\) by
\cref{plt:eq:div-recip-c0,plt:eq:div-recip-cn}; the right-hand side of
\cref{plt:eq:div-recip-cn} involves only \(C_{0},\dots,C_{n-1}\), so the
definition is legitimate by strong induction, and each \(C_{n}\) is a
\(\K\)-linear combination of products of polynomials, hence a polynomial.  Put
\(C\DefinitionEquals\sum_{n\ge0}C_{n}t^{n}\in\PLpow\).  For \(n\ge1\),
\[
  [t^{n}](BC)=\sum_{j=0}^{n}B_{j}C_{n-j}
  =cC_{n}+\sum_{j=1}^{n}B_{j}C_{n-j}=0
\]
by \cref{plt:eq:div-recip-cn}, and \([t^{0}](BC)=cC_{0}=1\).  Hence \(BC=1\).

\emph{Uniqueness.}  If \(BC=B\widetilde C=1\) with
\(C,\widetilde C\in\PLlau\), then
\(C=C(B\widetilde C)=(CB)\widetilde C=\widetilde C\).  This argument uses only
associativity and commutativity, so uniqueness holds in \(\PLlau\) -- and in
any commutative ring -- not merely in \(\PLpow\).

\emph{Laurent case.}  \(t^{m}\) is a unit with inverse \(t^{-m}\), so
\((t^{m}B)(t^{-m}C)=BC=1\).
\end{proof}

\begin{corollary}[Quotient recurrence]\label{plt:cor:div-quotient}
Let \(E\) be a field, \(A=\sum_{n\ge0}A_{n}t^{n}\) and
\(B=\sum_{n\ge0}B_{n}t^{n}\) in \(E[[t]]\) with \(B_{0}\ne0\).  Then
\(Q=A/B=\sum_{n\ge0}Q_{n}t^{n}\in E[[t]]\) is the unique solution of \(BQ=A\),
and
\begin{equation}\label{plt:eq:div-quotient-recurrence}
  Q_{n}=B_{0}^{-1}\Bigl(A_{n}-\sum_{j=1}^{n}B_{j}Q_{n-j}\Bigr),
  \qquad n\ge0,
\end{equation}
the sum being empty when \(n=0\).  When \(E=\K[L]\) is used instead of a field,
\cref{plt:eq:div-quotient-recurrence} still produces polynomial blocks provided
\(B_{0}\in\K^{\times}\), and then \(Q\in\PLpow\).
\end{corollary}

\begin{proof}
Existence and uniqueness follow from \cref{plt:thm:div-reciprocal} applied over
the coefficient ring \(E\) (or \(\K[L]\), where \(B_{0}\in\K^{\times}\) is
invertible): \(Q\DefinitionEquals A\,\recip{B}\) solves \(BQ=A\), and if
\(BQ=BQ'\) then \(Q=Q'\) after multiplying by \(\recip{B}\).  Comparing
coefficients of \(t^{n}\) in \(BQ=A\) gives
\(B_{0}Q_{n}+\sum_{j\ge1}B_{j}Q_{n-j}=A_{n}\), which is
\cref{plt:eq:div-quotient-recurrence}.  Since
\cref{plt:eq:div-quotient-recurrence} determines \(Q_{n}\) from
\(Q_{0},\dots,Q_{n-1}\), it also proves uniqueness directly.
\end{proof}

The next statement is the one an implementation actually needs: it says that a
finite computation with truncated inputs returns exactly the truncation of the
true quotient, and that the truncated quotient is characterised without
reference to the recurrence at all.

\begin{theorem}[Correctness of truncated division]
\label{plt:thm:div-truncated}
Let \(B\in\PLpow\) with \(B_{0}\in\K^{\times}\), let \(A\in\PLpow\), and let
\(N\ge0\).  Let \(Q_{0},\dots,Q_{N-1}\) be produced by
\cref{plt:eq:div-quotient-recurrence} and set
\(P\DefinitionEquals\sum_{n<N}Q_{n}t^{n}\).  Then
\begin{equation}\label{plt:eq:div-truncated-residual}
  BP=A+\FormalBigOAt{t^{N}}{t},
\end{equation}
and \(P\) is the \emph{only} element of \(\K[L][t]\) of \(t\)-degree \(<N\)
satisfying \cref{plt:eq:div-truncated-residual}.  Consequently
\(P=\Trunc{(A/B)}{N}\).
\end{theorem}

\begin{proof}
For \(n<N\) the coefficient \([t^{n}](BP)=\sum_{j=0}^{n}B_{j}Q_{n-j}\) involves
only \(Q_{0},\dots,Q_{n}\), all of which are among the computed blocks, and it
equals \(A_{n}\) by \cref{plt:eq:div-quotient-recurrence}.  This is
\cref{plt:eq:div-truncated-residual}.

Uniqueness: if \(P'\) is another such polynomial, put \(D=P-P'\), a polynomial
in \(t\) of degree \(<N\) with \(\ordt(BD)\ge N\).  If \(D\ne0\) then
\cref{plt:lem:div-valuation} gives \(\ordt(BD)=\ordt B+\ordt D=\ordt D<N\), a
contradiction; so \(D=0\).

Finally, \(B\bigl(A/B-P\bigr)=A-BP\) has \(\ordt\ge N\), and \(\ordt B=0\), so
\(\ordt(A/B-P)\ge N\) by \cref{plt:lem:div-valuation}.  As \(P\) has
\(t\)-degree \(<N\), this says exactly \(P=\Trunc{(A/B)}{N}\).
\end{proof}

% ed.: S6 states this correctness theorem with its local inclusive truncation
% [F]_{<=N}; the statement above is the exclusive normalisation demanded by the
% catalogue, and the two differ by exactly one block.
\begin{remark}[Truncation convention]\label{plt:rmk:div-truncation-convention}
Report S6 states \cref{plt:thm:div-truncated} in the inclusive form
\(B\sum_{n\le N}Q_{n}t^{n}=A+\FormalBigOAt{t^{N+1}}{t}\).  Because
\([F]_{\le N}=\Trunc{F}{N+1}\), that statement is the case \(N\mapsto N+1\) of
the one above.  Anyone transferring a precision budget between the two
conventions must account for the one-block shift; it is the commonest silent
off-by-one in this subject.
\end{remark}

\begin{proposition}[Exact input requirement: relative precision is preserved]
\label{plt:prop:div-relative-precision}
Let \(E\) be a field, and let \(A,B\in E((t))\) with \(A\ne0\ne B\) (the case
\(A=0\) being trivial).  Put \(a=\ordt A\) and \(b=\ordt B\), so that
\(\ordt(A/B)=a-b\).  For every \(M\ge0\), the first \(M\) blocks of \(A/B\),
namely the coefficients of \(t^{a-b},\dots,t^{a-b+M-1}\), are determined by the
first \(M\) blocks of \(A\) and the first \(M\) blocks of \(B\), and by nothing
else.  Equivalently, in absolute cutoffs,
\begin{equation}\label{plt:eq:div-guard}
  \Trunc{(A/B)}{N}
  \quad\text{is determined by}\quad
  \Trunc{A}{\,N+b\,}
  \ \text{and}\
  \Trunc{B}{\,N+2b-a\,}.
\end{equation}
\end{proposition}

\begin{proof}
Write \(\widetilde B=t^{-b}B\), so \(\ordt\widetilde B=0\) and
\(\widetilde B_{j}=B_{j+b}\).  Then \(A/B=t^{-b}\,(A/\widetilde B)\), and
\(G\DefinitionEquals A/\widetilde B\) has \(\ordt G=a\).  Applying
\cref{plt:eq:div-quotient-recurrence} to \(G\) after shifting indices by \(a\),
the block \(G_{a+i}\) is a function of \(A_{a},\dots,A_{a+i}\) and
\(\widetilde B_{0},\dots,\widetilde B_{i}\) alone.  Taking \(i<M\) gives the
first assertion, since \(\widetilde B_{0},\dots,\widetilde B_{M-1}\) are the
first \(M\) blocks of \(B\).  For \cref{plt:eq:div-guard}, the requested blocks
of \(A/B\) are those of index \(<N\), i.e.\ the first \(M=N-(a-b)\) blocks;
these need \(A\) through index \(a+M-1=N+b-1\) and \(B\) through index
\(b+M-1=N+2b-a-1\).
\end{proof}

% ed.: S6 states a "guard precision" principle in words ("compute the quotient
% beyond the final requested block") but never says by how much; the exact
% guard is the content of the proposition above.
\begin{remark}[Guard precision]\label{plt:rmk:div-guard}
\Cref{plt:prop:div-relative-precision} is the precise form of the informal rule
that a quotient must be computed ``beyond the requested block''.  The guard is
\(b=\ordt B\) blocks on the numerator and \(2b-a\) on the denominator; it is
zero exactly when \(a=b=0\), and it is what a symbolic system must propagate
through an expression tree instead of a single global cutoff.  A negative guard
means that fewer input blocks suffice, which happens when the denominator has a
Laurent principal part.
\end{remark}

\begin{example}[A reciprocal computed to five blocks]
\label{plt:ex:div-reciprocal-worked}
Let
\begin{equation}\label{plt:eq:div-recip-example-B}
  B=1+(L+1)t+L^{2}t^{2}+(L-2)t^{3}\in\PLpow .
\end{equation}
Here \(B_{0}=1\in\K^{\times}\) and \(B_{n}=0\) for \(n\ge4\).
\Cref{plt:eq:div-recip-c0,plt:eq:div-recip-cn} give
\begin{align*}
  C_{0}&=1,\\
  C_{1}&=-(L+1),\\
  C_{2}&=-\bigl((L+1)C_{1}+L^{2}C_{0}\bigr)=2L+1,\\
  C_{3}&=-\bigl((L+1)C_{2}+L^{2}C_{1}+(L-2)C_{0}\bigr)
        =L^{3}-L^{2}-4L+1,\\
  C_{4}&=-\bigl((L+1)C_{3}+L^{2}C_{2}+(L-2)C_{1}\bigr)
        =-L^{4}-2L^{3}+5L^{2}+2L-3,
\end{align*}
so that
\begin{equation}\label{plt:eq:div-recip-example-C}
  \recip{B}=1-(L+1)t+(2L+1)t^{2}+(L^{3}-L^{2}-4L+1)t^{3}
  +(-L^{4}-2L^{3}+5L^{2}+2L-3)t^{4}+\FormalBigOAt{t^{5}}{t}.
\end{equation}
Multiplying \cref{plt:eq:div-recip-example-B} by
\cref{plt:eq:div-recip-example-C} and truncating at \(t^{5}\) returns \(1\);
by \cref{plt:thm:div-truncated} this residual check certifies the five blocks
independently of the method used to produce them.
\end{example}

\begin{example}[A quotient with polynomial blocks]
\label{plt:ex:div-quotient-worked}
With
\[
  A=1+(L+1)t+(2L^{2}-L)t^{2}+\FormalBigOAt{t^{3}}{t},
  \qquad
  B=1+(2-L)t+(L^{2}+3)t^{2}+\FormalBigOAt{t^{3}}{t},
\]
\cref{plt:eq:div-quotient-recurrence} gives
\(Q_{0}=1\), \(Q_{1}=(L+1)-(2-L)=2L-1\), and
\[
  Q_{2}=(2L^{2}-L)-(2-L)(2L-1)-(L^{2}+3)=3L^{2}-6L-1,
\]
so \(A/B=1+(2L-1)t+(3L^{2}-6L-1)t^{2}+\FormalBigOAt{t^{3}}{t}\).  The block
\(Q_{3}\) is \emph{not} determined by the displayed data:
\cref{plt:prop:div-relative-precision} with \(a=b=0\) says three input blocks
determine exactly three output blocks.  If one additionally posits
\(A_{3}=B_{3}=0\), then \(Q_{3}=L^{3}-11L^{2}+5L+5\).
\end{example}

% ed.: S1 prints the block Q_3 of this example immediately after inputs given
% only modulo t^3; the hypothesis that the omitted blocks vanish is a genuine
% extra assumption and is stated here.
\begin{warningbox}[title={Do not read a block that the data does not contain}]
An input written with a formal tail \(\FormalBigOAt{t^{3}}{t}\) determines the
quotient only through \(t^{2}\).  Continuing the recurrence past that point
silently substitutes the assumption ``the omitted blocks are zero''.  Both S1
and S6 print one extra block of a worked example on exactly this unstated
convention; when it is intended it must be said.
\end{warningbox}

\begin{proposition}[Degree profile of a reciprocal and of a quotient]
\label{plt:prop:div-degree-profile}
Let \(\alpha\ge0\) and \(\beta\ge0\) be real.  Suppose \(B\in\PLpow\) has
\(B_{0}\in\K^{\times}\) and \(\degL B_{n}\le\alpha n\) for all \(n\ge1\).  Then
the blocks of \(\recip{B}\) satisfy
\begin{equation}\label{plt:eq:div-degree-recip}
  \degL C_{n}\le\alpha n\qquad(n\ge0).
\end{equation}
If moreover \(A\in\PLpow\) has \(\degL A_{n}\le\alpha n+\beta\), then the
quotient blocks satisfy \(\degL Q_{n}\le\alpha n+\beta\).
\end{proposition}

\begin{proof}
Strong induction on \(n\).  For \(n=0\), \(C_{0}=B_{0}^{-1}\in\K\) has degree
\(0\le0\).  For \(n\ge1\), each summand of \cref{plt:eq:div-recip-cn} obeys
\(\degL(B_{j}C_{n-j})\le\alpha j+\alpha(n-j)=\alpha n\) by the inductive
hypothesis and the additivity of \(\degL\) on products in \(\K[L]\); a finite
sum cannot raise the bound, and multiplying by the scalar \(-B_{0}^{-1}\)
changes nothing.  The quotient case is identical, with the extra term \(A_{n}\)
of degree \(\le\alpha n+\beta\) and \(\degL(B_{j}Q_{n-j})\le\alpha
j+\alpha(n-j)+\beta\).
\end{proof}

\begin{remark}
The bound in \cref{plt:eq:div-degree-recip} may be strict: cancellation can
lower a block's degree, and it does so in
\cref{plt:ex:div-reciprocal-worked}, where \(\degL B_{2}=2\) and \(\degL
C_{2}=1\).  An affine profile \(\dprofile{F}(n)\le\alpha n+\beta\) is a
hypothesis or a proved invariant, never part of the definition of the ring; the
Lambert expansion is the archetype with \(\alpha=1\), \(\beta=0\).
\end{remark}

\subsection{Geometric inversion and two-level summability}
\label{plt:sec:div-geometric}

The geometric formula \cref{plt:eq:div-geometric} is the conceptual
construction, and in \(\PLlau\) its justification is purely \(t\)-adic:
\(\ordt(U^{k})\ge k\to\infty\).  Over a completed logarithmic coefficient field
this reason is no longer available, and a second level of the argument is
needed.  Several sources normalise a divisor in \(\PLit\) as
\(B=b\,t^{\nu}L^{\mu}(1+S)\) with \(S\precdom1\) and then assert summability of
\(\sum_{k}(-S)^{k}\) ``by the support conditions''.  That is true but not for
the \(t\)-adic reason: \(S\) may have \(\ordt S=0\), so \(\ordt(S^{k})=0\) for
every \(k\).  The correct statement is the following.

\begin{lemma}[Two-level summability]\label{plt:lem:div-two-level}
For \(s\in\K((L^{-1}))\) let \(\omega(s)\) denote the \(L^{-1}\)-adic order, so
that \(s\in L^{-\omega(s)}\K[[L^{-1}]]\) and \(\omega(0)=+\infty\); \(\omega\)
is additive on products and satisfies
\(\omega(s+s')\ge\min\SetOf{\omega(s),\omega(s')}\).  Let
\(S=\sum_{n\ge0}s_{n}t^{n}\in\PLit\) satisfy
\begin{equation}\label{plt:eq:div-small-S}
  \ordt S\ge0
  \qquad\text{and}\qquad
  \omega(s_{0})\ge1 ,
\end{equation}
which is exactly the condition \(S\precdom1\) in the dominance order.  Fix
\(n\ge0\) and put
\(c_{n}\DefinitionEquals\max\SetOf{0,-\omega(s_{1}),\dots,-\omega(s_{n})}\), a
finite nonnegative integer.  Then for every \(k\ge0\),
\begin{equation}\label{plt:eq:div-two-level-order}
  \omega\bigl([t^{n}]S^{k}\bigr)\ge k-n\,(1+c_{n}).
\end{equation}
Consequently, for each fixed monomial \(t^{n}L^{-j}\) only the finitely many
indices \(k\le j+n(1+c_{n})\) contribute to \(\sum_{k\ge0}(-S)^{k}\); the
family is summable in \(\PLit\), and \((1+S)\sum_{k\ge0}(-S)^{k}=1\).
\end{lemma}

\begin{proof}
Because \(\ordt S\ge0\),
\[
  [t^{n}]S^{k}=\sum_{\substack{(n_{1},\dots,n_{k})\in\NaturalNumbers^{k}\\
                               n_{1}+\cdots+n_{k}=n}}
    s_{n_{1}}\cdots s_{n_{k}},
\]
a finite sum in which every index satisfies \(n_{i}\le n\).  Fix one tuple and
let \(r\) be the number of indices with \(n_{i}\ge1\); then \(r\le n\), since
each such index contributes at least \(1\) to the total \(n\).  The \(k-r\)
factors with \(n_{i}=0\) have \(\omega(s_{0})\ge1\), and the \(r\) remaining
factors have \(\omega(s_{n_{i}})\ge-c_{n}\) by definition of \(c_{n}\).
Additivity of \(\omega\) on products gives
\[
  \omega\bigl(s_{n_{1}}\cdots s_{n_{k}}\bigr)
  \ge(k-r)\cdot1+r\cdot(-c_{n})
  \ge k-n-n c_{n},
\]
using \(r\le n\) and \(c_{n}\ge0\).  Taking the minimum over the finitely many
tuples yields \cref{plt:eq:div-two-level-order}.

A power \(L^{-j}\) can occur in an element of order \(\omega\) only when
\(j\ge\omega\); hence \([t^{n}L^{-j}]S^{k}=0\) as soon as
\(k-n(1+c_{n})>j\), which proves the finiteness assertion and therefore
summability.  For the identity, fix \((n,j)\) and choose \(K\) larger than
\(j+n(1+c_{n})\).  Telescoping gives
\((1+S)\sum_{k\le K}(-S)^{k}=1-(-S)^{K+1}\), and by the finiteness just proved
the coefficient of \(t^{n}L^{-j}\) in \((-S)^{K+1}\) vanishes and the
coefficient of \(t^{n}L^{-j}\) in \((1+S)\sum_{k\le K}(-S)^{k}\) already equals
that in \((1+S)\sum_{k\ge0}(-S)^{k}\).  Hence the two sides agree at every
monomial.
\end{proof}

\begin{remark}
\Cref{plt:lem:div-two-level} is why the catalogue insists that a finite
representation in \(\PLit\) needs an outer \(t\)-cutoff \emph{and} an inner
\(L^{-1}\)-cutoff, or one explicitly declared staircase.  A single outer cutoff
does not bound the work: the coefficient of \(t^{0}L^{-j}\) still receives
contributions from \(j+1\) powers of \(S\).
\end{remark}

\subsection{Why unrestricted division needs the rational-log field}
\label{plt:sec:div-rational-field}

\begin{theorem}[Laurent series over a field]\label{plt:thm:div-Ett-field}
Let \(E\) be a field.  Then \(E((t))\) is a field.  In particular
\(\PLrat=\K(L)((t))\) and \(\PLit=\K((L^{-1}))((t))\) are fields.
\end{theorem}

\begin{proof}
\(E((t))\) is a commutative ring with \(1\), and it is an integral domain by
\cref{plt:lem:div-valuation}.  Let \(F\ne0\) with \(m=\ordt F\) and leading
coefficient \(f_{m}\in E^{\times}\).  Put \(U\DefinitionEquals
f_{m}^{-1}t^{-m}F-1\in tE[[t]]\).  Exactly as in the proof of
\cref{plt:thm:div-unit-criterion}, \(\ordt(U^{k})\ge k\) makes
\(\sum_{k\ge0}(-U)^{k}\) formally summable in \(E[[t]]\), and telescoping gives
\((1+U)\sum_{k\ge0}(-U)^{k}=1\).  Hence
\(F\cdot f_{m}^{-1}t^{-m}\sum_{k\ge0}(-U)^{k}=1\).  Finally \(\K(L)\) and
\(\K((L^{-1}))\) are fields: the first is the fraction field of the integral
domain \(\K[L]\); the second is the Laurent series field over \(\K\) in the
variable \(L^{-1}\), a field by the case already proved.
\end{proof}

The next theorem says that \(\PLrat\) is not merely \emph{a} field containing
\(\PLlau\), but the smallest one of the shape ``Laurent series in \(t\) over a
coefficient ring between \(\K[L]\) and \(\K(L)\)''.  This is the precise sense
in which unrestricted division of \emph{coefficient blocks} forces \(\K(L)\).

\begin{theorem}[Minimality of the rational-log coefficient field]
\label{plt:thm:div-minimal-field}
Let \(E\) be a subring of \(\K(L)\) with \(\K[L]\subseteq E\).  Then \(E((t))\)
is a field if and only if \(E=\K(L)\).  Consequently \(\PLrat\) is the unique
smallest field of the form \(E((t))\) containing \(\PLlau\) with
\(E\subseteq\K(L)\).
\end{theorem}

\begin{proof}
If \(E=\K(L)\), apply \cref{plt:thm:div-Ett-field}.  Conversely suppose
\(E((t))\) is a field and let \(e\in E\), \(e\ne0\).  Its inverse
\(G=\sum_{n}G_{n}t^{n}\in E((t))\) satisfies \(eG=1\), i.e.\ \(eG_{0}=1\) and
\(eG_{n}=0\) for \(n\ne0\).  Since \(E\subseteq\K(L)\) is an integral domain
and \(e\ne0\), we get \(G_{n}=0\) for \(n\ne0\) and \(G_{0}=e^{-1}\in E\).
Hence \(E\) is a field.  A subfield of \(\K(L)\) containing \(\K[L]\) contains
\(p^{-1}\) for every nonzero \(p\in\K[L]\) and hence every quotient \(p/q\),
so \(E=\K(L)\).
\end{proof}

\subsection{An integral domain that is not a field, made precise}
\label{plt:sec:div-domain}

``\(\PLlau\) is an integral domain but not a field'' is the standard slogan.
It is worth saying exactly how far from a field it is, and -- a point none of
the six sources records -- that its field of fractions is \emph{strictly
smaller} than \(\PLrat\).  Passing to \(\PLrat\) is therefore not the same
operation as ``forming the fraction field''.

\begin{lemma}[Denominator lemma]\label{plt:lem:div-denominator}
Let \(A,B\in\PLlau\) with \(B\ne0\), and let \(b_{0}\in\K[L]\) be the leading
block of \(B\).  Then, computed in \(\PLrat\),
\begin{equation}\label{plt:eq:div-denominator}
  \frac AB\in\K[L][b_{0}^{-1}]\,((t)),
\end{equation}
where \(\K[L][b_{0}^{-1}]=\SetOf{p/b_{0}^{k}:p\in\K[L],\,k\ge0}\subseteq\K(L)\).
In particular every element of the fraction field of \(\PLlau\) has all of its
coefficient blocks in a \emph{single} localisation \(\K[L][q^{-1}]\) of
\(\K[L]\) at one nonzero polynomial \(q\).
\end{lemma}

\begin{proof}
Write \(m=\ordt B\) and \(S\DefinitionEquals\K[L][b_{0}^{-1}]\), a subring of
\(\K(L)\).  Then \(t^{-m}B=\sum_{k\ge0}b_{k}t^{k}\) with \(b_{k}\in\K[L]\), so
\(b_{0}^{-1}t^{-m}B=1+U\) with \(U\DefinitionEquals\sum_{k\ge1}(b_{k}/b_{0})
t^{k}\in tS[[t]]\).  By the proof of \cref{plt:thm:div-Ett-field} applied to
the ring \(S\) -- only summability and telescoping are used, not that \(S\) is
a field -- \((1+U)^{-1}=\sum_{k\ge0}(-U)^{k}\in S[[t]]\).  Hence
\(B^{-1}=t^{-m}b_{0}^{-1}(1+U)^{-1}\in S((t))\) and, since
\(A\in\PLlau\subseteq S((t))\), also \(A/B\in S((t))\).  The last sentence is
this statement with \(q=b_{0}\).
\end{proof}

\begin{theorem}[The fraction field of \(\PLlau\) is strictly smaller than
\(\PLrat\)]\label{plt:thm:div-frac-strict}
Let \(\K\) have characteristic zero.  Then
\begin{equation}\label{plt:eq:div-frac-strict}
  \PLlau\subsetneq\operatorname{Frac}(\PLlau)\subsetneq\PLrat .
\end{equation}
An explicit witness for the second strict inclusion is
\begin{equation}\label{plt:eq:div-frac-witness}
  F\DefinitionEquals\sum_{n\ge0}\frac{t^{n}}{L-n}\in\PLrat .
\end{equation}
\end{theorem}

\begin{proof}
The first inclusion is strict because \(L\) is a nonzero nonunit
(\cref{plt:ex:div-nonunit}).  For the second, \(\operatorname{Frac}(\PLlau)\)
embeds in the field \(\PLrat\) because \(\PLrat\) contains \(\PLlau\), so we
must only show \(F\notin\operatorname{Frac}(\PLlau)\).  Certainly
\(F\in\PLrat\): its blocks \((L-n)^{-1}\) lie in \(\K(L)\) and its support is
bounded below by \(0\).

Suppose \(F=A/B\) with \(A,B\in\PLlau\), \(B\ne0\).  By
\cref{plt:lem:div-denominator} there is a nonzero \(q\in\K[L]\) with every
block of \(F\) in \(\K[L][q^{-1}]\).  Fix \(n\ge0\).  Then
\((L-n)^{-1}=p/q^{k}\) for some \(p\in\K[L]\) and \(k\ge0\), i.e.\
\(q^{k}=(L-n)p\).  Since \(\K\) has characteristic zero, the elements
\(0,1,2,\dots\) of \(\K\) are pairwise distinct, so \(L-n\) is a nonassociate
irreducible of the unique factorisation domain \(\K[L]\) for each \(n\), and
\((L-n)\mid q^{k}\) forces \((L-n)\mid q\).  As this holds for every
\(n\ge0\), the polynomial \(q\) has infinitely many roots, hence \(q=0\), a
contradiction.
\end{proof}

\begin{corollary}[The precise structure]\label{plt:cor:div-domain-not-field}
\(\PLlau\) is an integral domain whose unit group is
\cref{plt:eq:div-unit-group}.  It is not a field, and its nonunits do not form
an ideal, so it is not a local ring: \(L\) and \(1-L\) are nonunits whose sum
is the unit \(1\).  The obstruction to division is therefore not concentrated
at one maximal ideal but is spread over the primes of \(\K[L]\): for each monic
irreducible \(\pi\in\K[L]\), the set \(\pi\PLlau\) is a proper nonzero ideal,
and \(\bigcap_{\pi}\pi\PLlau=0\).
\end{corollary}

\begin{proof}
The first two sentences have been proved.  For the last, \(\pi\PLlau\) is an
ideal; it is nonzero and proper because \(\pi\) is a nonzero nonunit
(\cref{plt:thm:div-unit-criterion}).  If \(F\ne0\) lay in \(\pi\PLlau\) for
every \(\pi\), its leading block \(p_{\ordt F}\) would be divisible in
\(\K[L]\) by every monic irreducible -- by \cref{plt:lem:div-valuation} the
leading block of \(\pi G\) is \(\pi\) times the leading block of \(G\) -- and
\(\K[L]\) has infinitely many pairwise nonassociate irreducibles (Euclid's
argument applies verbatim), forcing \(p_{\ordt F}=0\).
\end{proof}

\subsection{Conditional divisibility}
\label{plt:sec:div-conditional}

A nonunit divisor still divides some elements exactly.  The following
proposition makes the test explicit; it also shows that ``being divisible'' is
a genuinely global condition, so that no unconditional division theorem in
\(\PLlau\) is available.

\begin{proposition}[Coefficientwise divisibility test]
\label{plt:prop:div-conditional}
Let \(B\in\PLlau\) be nonzero with \(\ordt B=0\) and leading block \(B_{0}\in
\K[L]\), and let \(A\in\PLpow\).  Let \(Q=A/B\in\K(L)[[t]]\) have blocks
\(Q_{n}\) given by \cref{plt:eq:div-quotient-recurrence} over \(E=\K(L)\).
Then:
\begin{enumerate}
\item \(A/B\in\PLpow\) if and only if \(Q_{n}\in\K[L]\) for every \(n\ge0\);
\item given \(Q_{0},\dots,Q_{n-1}\in\K[L]\), the block \(Q_{n}\) lies in
      \(\K[L]\) if and only if \(B_{0}\) divides
      \(A_{n}-\sum_{j=1}^{n}B_{j}Q_{n-j}\) in \(\K[L]\);
\item the set \(\SetOf{A\in\PLlau:A/B\in\PLlau}\) is exactly the principal
      ideal \(B\PLlau\), which is nonzero, and proper as soon as \(B\) is not a
      unit.
\end{enumerate}
\end{proposition}

\begin{proof}
(1) is the definition of \(\PLpow\) as the set of series with polynomial
blocks, together with the uniqueness in \cref{plt:cor:div-quotient}: the blocks
of \(A/B\) computed over \(\K(L)\) are the only candidates.  (2) is immediate
from \cref{plt:eq:div-quotient-recurrence}, since \(Q_{n}=B_{0}^{-1}R_{n}\)
with \(R_{n}=A_{n}-\sum_{j\ge1}B_{j}Q_{n-j}\in\K[L]\) under the stated
hypothesis.  For (3): \(A/B\in\PLlau\) means \(A=B\cdot(A/B)\) with
\(A/B\in\PLlau\), i.e.\ \(A\in B\PLlau\); conversely \(A=BC\) with
\(C\in\PLlau\) gives \(A/B=C\).  The ideal \(B\PLlau\) contains \(B\ne0\), and
it equals \(\PLlau\) exactly when \(1\in B\PLlau\), i.e.\ when \(B\) is a unit.
\end{proof}

\begin{example}[Exact division by a nonunit, and a near miss]
\label{plt:ex:div-conditional}
Take \(B=L+t\), a nonzero nonunit.  Then
\begin{equation}\label{plt:eq:div-exact-nonunit}
  \frac{L^{2}-t^{2}}{L+t}=L-t\in\PLpow ,
\end{equation}
verified by the single multiplication \((L+t)(L-t)=L^{2}-t^{2}\).  So a nonunit
can divide exactly.  Change one sign and the divisibility fails at the third
step: with \(A=L^{2}+t^{2}\), the test of
\cref{plt:prop:div-conditional}\,(2) gives
\[
  Q_{0}=\frac{L^{2}}{L}=L,\qquad
  Q_{1}=\frac{0-1\cdot L}{L}=-1,\qquad
  Q_{2}=\frac{1-1\cdot(-1)}{L}=\frac2L\notin\K[L],
\]
consistently with the exact identity
\(\dfrac{L^{2}+t^{2}}{L+t}=L-t+\dfrac{2t^{2}}{L+t}\).  Two successful steps
prove nothing about the third.
\end{example}

\begin{remark}[The test is a semi-decision procedure]
\label{plt:rmk:div-semidecision}
\Cref{plt:prop:div-conditional}\,(2) refutes divisibility in finite time
whenever it fails at some finite order.  It never confirms divisibility in
finite time on its own, because the condition quantifies over all \(n\).
A positive answer needs a finite certificate: either a candidate quotient
\(C\in\PLlau\) with only finitely many blocks -- so that \(BC=A\) is a finite
identity, as in \cref{plt:eq:div-exact-nonunit} -- or an a priori bound
reducing the infinitely many conditions to finitely many, or a structural
argument.  Absent such a certificate, a computation that ``keeps coming out
polynomial'' is evidence, not proof.
\end{remark}

\begin{checkpointbox}
Three closure questions must be kept apart.
\begin{enumerate}
\item Does the quotient exist in \emph{some} field?  Always, in \(\PLrat\).
\item Does it have rational blocks and nothing worse?  Always, again in
      \(\PLrat\); this is \cref{plt:thm:div-Ett-field}.
\item Does it have polynomial blocks?  Only under the unit criterion, or by a
      certified accident as in \cref{plt:ex:div-conditional}.
\end{enumerate}
Answering (1) and reporting it as an answer to (3) is the standard failure
mode.
\end{checkpointbox}

\subsection{Newton iteration and the exact precision-doubling law}
\label{plt:sec:div-newton}

The reciprocal recurrence produces one block per step.  Newton's method for
the equation \(C^{-1}-B=0\) produces the map
\begin{equation}\label{plt:eq:div-newton-step}
  C\longmapsto C(2-BC),
\end{equation}
which doubles the number of correct blocks.  Every source states
\cref{plt:eq:div-newton-step} and the identity \(1-BC(2-BC)=(1-BC)^{2}\); none
states which truncations of \(B\) and of \(C\) suffice to compute the doubled
output, which is precisely what makes the scheme an algorithm rather than an
identity.  Part (4) below supplies that statement.

\begin{theorem}[Precision doubling]\label{plt:thm:div-newton}
Let \(B\in\PLpow\) with \(B_{0}\in\K^{\times}\), let \(C\in\PLpow\), let
\(N\ge1\), and suppose
\begin{equation}\label{plt:eq:div-newton-invariant}
  \ordt(1-BC)\ge N .
\end{equation}
Put \(C'\DefinitionEquals C(2-BC)\).  Then:
\begin{enumerate}
\item \(1-BC'=(1-BC)^{2}\), hence \(\ordt(1-BC')\ge2N\);
\item \(\ordt\bigl(\recip{B}-C'\bigr)=\ordt(1-BC')\ge 2N\), and more
      precisely \(\ordt(\recip{B}-C)=\ordt(1-BC)\);
\item \(\ordt\bigl(1-B\Trunc{C'}{2N}\bigr)\ge2N\), so truncating the output at
      order \(2N\) preserves the invariant;
\item \(\Trunc{C'}{2N}\) depends on \(B\) and \(C\) only through
      \(\Trunc{B}{2N}\) and \(\Trunc{C}{N}\).  Explicitly, if
      \(\overline B\in\PLpow\) and \(\overline C\in\PLpow\) satisfy
      \(\ordt(B-\overline B)\ge2N\) and \(\ordt(C-\overline C)\ge N\), then
      \(\ordt\bigl(C'-\overline C(2-\overline B\,\overline C)\bigr)\ge2N\).
\end{enumerate}
\end{theorem}

\begin{proof}
(1)  \(BC'=BC(2-BC)=2BC-(BC)^{2}\), so
\(1-BC'=1-2BC+(BC)^{2}=(1-BC)^{2}\).  Valuations add
(\cref{plt:lem:div-valuation}), so \(\ordt(1-BC')\ge2N\).

(2)  \(B\) is a unit by \cref{plt:thm:div-unit-criterion} and
\(\ordt(\recip{B})=-\ordt B=0\).  From
\(\recip{B}-C'=\recip{B}\,(1-BC')\) and additivity of \(\ordt\) we get
\(\ordt(\recip{B}-C')=\ordt(1-BC')\).  The same computation with \(C\) in
place of \(C'\) gives the second identity.

(3)  \(C'-\Trunc{C'}{2N}\in t^{2N}\PLpow\), so
\(B(C'-\Trunc{C'}{2N})\in t^{2N}\PLpow\) because \(\ordt B=0\); adding this to
(1) gives \(\ordt(1-B\Trunc{C'}{2N})\ge2N\).

% ed.: the input-requirement statement (4) is absent from S2, S3, S5 and S6,
% each of which asserts doubling without saying which truncations suffice.
(4)  Put \(\delta=C-\overline C\) and \(\varepsilon=B-\overline B\), so
\(\ordt\delta\ge N\) and \(\ordt\varepsilon\ge2N\).  Since
\(C'=2C-BC^{2}\) and \(\overline C(2-\overline B\,\overline C)
=2\overline C-\overline B\,\overline C^{2}\),
\[
  C'-\overline C(2-\overline B\,\overline C)
  =2\delta-\bigl(BC^{2}-\overline B\,\overline C^{2}\bigr)
  =2\delta-\varepsilon C^{2}
    -\overline B\,(C+\overline C)\,\delta
  =\delta\bigl(2-\overline B(C+\overline C)\bigr)-\varepsilon C^{2},
\]
using \(BC^{2}-\overline B\,\overline C^{2}
=\varepsilon C^{2}+\overline B(C^{2}-\overline C^{2})\) and
\(C^{2}-\overline C^{2}=(C+\overline C)\delta\).  The second term has
\(\ordt\ge2N\).  For the first,
\[
  2-\overline B(C+\overline C)
  =\bigl(1-\overline BC\bigr)+\bigl(1-\overline B\,\overline C\bigr),
\]
and \(1-\overline BC=(1-BC)+\varepsilon C\) has \(\ordt\ge N\), while
\(1-\overline B\,\overline C=(1-BC)+\varepsilon C+B\delta-\varepsilon\delta\)
also has \(\ordt\ge N\).  Hence
\(\ordt\bigl(\delta(2-\overline B(C+\overline C))\bigr)\ge N+N=2N\), and the
whole difference has \(\ordt\ge2N\).
\end{proof}

\begin{algorithm}[Newton reciprocal to exclusive order \(N\)]
\label{plt:alg:div-newton}
Given \(B\in\PLpow\) with \(B_{0}\in\K^{\times}\), known through order
\(t^{N-1}\):
\begin{lstlisting}[style=pseudo]
function newton_reciprocal(B, N):
    require B[0] is a nonzero scalar in K
    C := 1/B[0]                    # exact through order t^0
    m := 1
    while m < N:
        m2 := min(2*m, N)
        E  := Trunc(1 - B*C, m2)   # valuation >= m by the invariant
        C  := Trunc(C*(1 + E), m2)
        m  := m2
    return C                       # equals Trunc_{<N} of the reciprocal
\end{lstlisting}
\end{algorithm}

\begin{corollary}[The iterates are exact truncations]
\label{plt:cor:div-newton-exact}
Let \(C^{(0)}=B_{0}^{-1}\) and \(C^{(k+1)}=\Trunc{\bigl(C^{(k)}
(2-BC^{(k)})\bigr)}{2^{k+1}}\).  Then for every \(k\ge0\),
\begin{equation}\label{plt:eq:div-newton-exact}
  C^{(k)}=\Trunc{\recip{B}}{2^{k}} .
\end{equation}
In particular \cref{plt:alg:div-newton} returns \(\Trunc{\recip{B}}{N}\), and
the intermediate steps require no blocks of \(B\) beyond index \(2^{k+1}-1\)
at step \(k\).
\end{corollary}

\begin{proof}
Induction on \(k\).  For \(k=0\), \(1-BC^{(0)}=1-B_{0}^{-1}B\) has zero
constant block, so \(\ordt(1-BC^{(0)})\ge1\), and by
\cref{plt:thm:div-newton}\,(2) \(\ordt(\recip{B}-C^{(0)})\ge1\); as
\(C^{(0)}\) has \(t\)-degree \(0<1\), it equals \(\Trunc{\recip{B}}{1}\).
Assume \(C^{(k)}=\Trunc{\recip{B}}{2^{k}}\); then
\(\ordt(\recip{B}-C^{(k)})\ge2^{k}\), so \(\ordt(1-BC^{(k)})\ge2^{k}\) by
\cref{plt:thm:div-newton}\,(2).  Parts (1) and (3) give
\(\ordt(1-BC^{(k+1)})\ge2^{k+1}\), hence
\(\ordt(\recip{B}-C^{(k+1)})\ge2^{k+1}\) by (2) again; since \(C^{(k+1)}\) has
\(t\)-degree \(<2^{k+1}\), it is \(\Trunc{\recip{B}}{2^{k+1}}\).  The
statement about required blocks of \(B\) is
\cref{plt:thm:div-newton}\,(4).
\end{proof}

% ed.: S2 justifies the doubling by "the valuation is multiplicative"; the
% valuation is additive on products, which is what the argument uses.
\begin{remark}[A wording to avoid]\label{plt:rmk:div-valuation-additive}
The step \(\ordt(E^{2})\ge2N\) holds because \(\ordt\) is \emph{additive} on
products (\cref{plt:lem:div-valuation}), not because it is multiplicative;
\(\ordt\) is a homomorphism from the multiplicative monoid of nonzero elements
to \((\IntegerNumbers,+)\).
\end{remark}

\begin{remark}[Where the iteration lives]
\Cref{plt:thm:div-newton} used only that \(\ordt\) is additive and that
\(B\) has an invertible leading block; it therefore holds verbatim over the
coefficient rings \(\K(L)\) and \(\K((L^{-1}))\), and in \(\PLit\) after
normalising the leading Laurent-log coefficient.  It is an exact algebraic
identity, not an approximation: unlike numerical Newton iteration, there is no
convergence hypothesis and no loss of significance -- by
\cref{plt:cor:div-newton-exact} each iterate is literally a truncation of the
answer, so the logarithmic degrees appearing in the iterates are exactly those
of \(\recip{B}\) and cannot swell beyond them, provided one truncates at every
step.
\end{remark}

\subsection{Cost: the recurrence against the doubling scheme}
\label{plt:sec:div-cost}

We count multiplications in the coefficient ring \(\K[L]\); additions and
multiplications by the scalar \(B_{0}^{-1}\in\K^{\times}\) are not counted.
Let \(M(N)\) denote the number of such block multiplications used by whatever
routine computes \(\Trunc{(PQ)}{N}\) from \(\Trunc{P}{N}\) and
\(\Trunc{Q}{N}\).

\begin{proposition}[Cost of the recurrence]\label{plt:prop:div-cost-naive}
Computing \(\Trunc{\recip{B}}{N}\) by
\cref{plt:eq:div-recip-c0,plt:eq:div-recip-cn} uses exactly
\begin{equation}\label{plt:eq:div-cost-naive}
  \binom{N}{2}=\frac{N(N-1)}{2}
\end{equation}
block multiplications, and computing \(\Trunc{(A/B)}{N}\) by
\cref{plt:eq:div-quotient-recurrence} uses the same number.
\end{proposition}

\begin{proof}
Step \(n\) forms the products \(B_{j}C_{n-j}\) for \(j=1,\dots,n\), that is
\(n\) block multiplications, for \(n=0,1,\dots,N-1\).  The total is
\(\sum_{n=0}^{N-1}n=N(N-1)/2\).
\end{proof}

\begin{proposition}[Cost of the doubling scheme]
\label{plt:prop:div-cost-newton}
Let \(N=2^{K}\).
\begin{enumerate}
\item If \(M\) is nondecreasing and satisfies \(2M(n)\le M(2n)\), then
      \cref{plt:alg:div-newton} uses at most \(4M(N)\) block multiplications.
\item With schoolbook truncated multiplication, and exploiting the invariant
      \(\ordt(1-BC)\ge m\) so that the first \(m\) blocks of \(BC\) are known
      in advance, \cref{plt:alg:div-newton} uses exactly
      \begin{equation}\label{plt:eq:div-cost-newton-exact}
        \frac{N^{2}+N-2}{2}
      \end{equation}
      block multiplications, which exceeds
      \cref{plt:eq:div-cost-naive} by exactly \(N-1\).
\end{enumerate}
\end{proposition}

\begin{proof}
(1)  The step from precision \(m\) to \(2m\) forms two truncated products at
order \(2m\), so costs at most \(2M(2m)\).  Summing over
\(m=1,2,4,\dots,2^{K-1}\) gives at most \(\sum_{j=1}^{K}2M(2^{j})\).  The
hypothesis \(2M(n)\le M(2n)\) gives \(M(2^{j})\le M(2^{K})/2^{K-j}\), so the
sum is at most \(2M(N)\sum_{j=1}^{K}2^{\,j-K}\le4M(N)\).

(2)  Consider the step from \(m\) to \(2m\), with \(C=\Trunc{\recip{B}}{m}\)
supported on indices \(0,\dots,m-1\).
\emph{Forming \(E=1-BC\) modulo \(t^{2m}\).}  By the invariant, the blocks of
\(BC\) of index \(<m\) are known to be \(1,0,\dots,0\); only indices
\(m,\dots,2m-1\) must be computed.  The needed products are \(B_{j}C_{i}\) with
\(0\le i<m\) and \(m\le i+j<2m\), and for each of the \(m\) values of \(i\)
there are exactly \(m\) admissible \(j\).  Cost: \(m^{2}\).
\emph{Forming \(C(1+E)\) modulo \(t^{2m}\).}  \(E\) is supported on indices
\(m,\dots,2m-1\), and \(C\) on \(0,\dots,m-1\); the products \(C_{i}E_{j}\)
with \(i+j<2m\) number \(\sum_{j=m}^{2m-1}(2m-j)=\sum_{s=1}^{m}s=m(m+1)/2\).
Total for the step: \(m^{2}+m(m+1)/2=(3m^{2}+m)/2\).  Summing over
\(m=2^{k}\), \(k=0,\dots,K-1\),
\[
  \sum_{k=0}^{K-1}\frac{3\cdot4^{k}+2^{k}}{2}
  =\frac{3}{2}\cdot\frac{4^{K}-1}{3}+\frac{2^{K}-1}{2}
  =\frac{N^{2}-1}{2}+\frac{N-1}{2}
  =\frac{N^{2}+N-2}{2}.
\]
Subtracting \cref{plt:eq:div-cost-naive} gives \((2N-2)/2=N-1\).
\end{proof}

% ed.: S3 asserts that Newton inversion "is often faster than coefficient
% recursion for large N even with naive multiplication"; the exact counts above
% show the opposite at the level of coefficient multiplications, and the honest
% statement is the corollary.
\begin{corollary}[When doubling pays]\label{plt:cor:div-doubling-pays}
The doubling scheme improves on the recurrence exactly when block
multiplication is subquadratic.  Precisely: if \(M(N)/N^{2}\to0\) as
\(N\to\infty\) (with \(M\) as in \cref{plt:prop:div-cost-newton}\,(1)), then
the doubling scheme uses \(o(N^{2})\) block multiplications against
\(\Theta(N^{2})\) for the recurrence; whereas with schoolbook block
multiplication it uses exactly \(N-1\) \emph{more} block multiplications than
the recurrence.
\end{corollary}

\begin{proof}
Combine \cref{plt:prop:div-cost-naive,plt:prop:div-cost-newton}; the first
claim is part (1) with the stated growth hypothesis, the second is part (2).
\end{proof}

% ed.: S6 gives the cost as O(N^2 d^2) "if the typical coefficient degree is
% d"; under the affine log-degree profile that actually occurs there is no such
% constant d, and the honest count is quartic in N.
\begin{remark}[Scalar cost under an affine degree profile]
\label{plt:rmk:div-scalar-cost}
Block multiplications are not unit-cost.  Assume the hypothesis of
\cref{plt:prop:div-degree-profile} with \(\alpha>0\) an integer, so
\(\degL B_{n}\le\alpha n\) and \(\degL C_{n}\le\alpha n\), and count
\(\K\)-multiplications with schoolbook polynomial arithmetic in \(L\).  The
recurrence then costs at most
\[
  \sum_{n=1}^{N-1}\sum_{j=1}^{n}\bigl(\alpha j+1\bigr)
       \bigl(\alpha(n-j)+1\bigr)
  =\Theta\bigl(\alpha^{2}N^{4}\bigr),
  \qquad\text{with leading term }\frac{\alpha^{2}N^{4}}{24},
\]
because \(\sum_{j=1}^{n}j(n-j)=\binom{n+1}{3}\sim n^{3}/6\) and
\(\sum_{n<N}\binom{n+1}{3}=\binom{N+1}{4}\sim N^{4}/24\).  Quoting a bound
\(\BigO(N^{2}d^{2})\)
with a fixed \(d\) therefore understates the cost whenever the log-degree
profile grows; substituting the honest \(d\asymp\alpha N\) recovers
\(\alpha^{2}N^{4}\).  This quartic behaviour, not the quadratic block count, is
what makes truncating after \emph{every} multiplication the decisive
implementation discipline; see \cite{vanderhoeven-book} for the relaxed and
fast-multiplication variants.
\end{remark}

\subsection{The embedding \texorpdfstring{\(\PLrat\hookrightarrow\PLit\)}%
{K(L)((t)) into K((1/L))((t))}}
\label{plt:sec:div-embedding}

The third response to a nonconstant leading block is to expand each rational
coefficient at \(L=\infty\).  This resolves a finer asymptotic scale, and it is
the natural home of a multiseries implementation \cite{salvy-shackell-1998};
it is also the step most likely to be taken without noticing that the object
has changed.

\begin{definition}[Iterated Laurent-log field]\label{plt:def:div-iterated}
\(\K((L^{-1}))\) is the field of formal Laurent series in the variable
\(L^{-1}\): elements \(\sum_{j\ge j_{0}}a_{j}L^{-j}\) with
\(j_{0}\in\IntegerNumbers\) and \(a_{j}\in\K\), so that only finitely many
positive powers of \(L\) occur but arbitrarily many negative ones may.  Then
\(\PLit\DefinitionEquals\K((L^{-1}))((t))\), whose elements are
\(\sum_{n\ge n_{0}}t^{n}\sum_{j\ge j_{n}}a_{n,j}L^{-j}\).
\end{definition}

\begin{theorem}[The expansion map]\label{plt:thm:div-embedding}
There is exactly one \(\K\)-algebra homomorphism
\(\iota_{0}:\K(L)\to\K((L^{-1}))\) with \(\iota_{0}(L)=L\), namely Laurent
expansion at \(L=\infty\), and it is injective.  Applying it blockwise defines
an injective ring homomorphism \(\iota:\PLrat\to\PLit\) that restricts to the
identity on \(\PLlau\).  The map \(\iota\) is \emph{not} surjective; its image
is therefore a proper subfield of \(\PLit\).
\end{theorem}

\begin{proof}
\emph{Existence on coefficients.}  A polynomial \(\sum_{i\le d}c_{i}L^{i}\)
already lies in \(\K((L^{-1}))\) (as \(\sum_{j\ge-d}c_{-j}L^{-j}\)), giving an
injective ring map \(\K[L]\hookrightarrow\K((L^{-1}))\).  Since
\(\K((L^{-1}))\) is a field (\cref{plt:thm:div-Ett-field}), every nonzero
polynomial maps to a unit, so by the universal property of the localisation
\(\K(L)=\operatorname{Frac}\K[L]\) the map extends uniquely to a ring
homomorphism \(\iota_{0}:\K(L)\to\K((L^{-1}))\); a ring homomorphism out of a
field is injective.  Concretely, for \(q=q_{d}L^{d}(1+w)\) with
\(q_{d}\in\K^{\times}\) and \(w\in L^{-1}\K[L^{-1}]\),
\(\iota_{0}(1/q)=q_{d}^{-1}L^{-d}\sum_{k\ge0}(-w)^{k}\), the series being
summable because \(w\) has positive \(L^{-1}\)-order.

Uniqueness of \(\iota_{0}\) is the uniqueness clause in that universal
property: a \(\K\)-algebra map on \(\K(L)\) is determined by the image of
\(L\).

\emph{Extension to Laurent series.}  Define
\(\iota\bigl(\sum_{n\ge n_{0}}f_{n}t^{n}\bigr)
=\sum_{n\ge n_{0}}\iota_{0}(f_{n})t^{n}\).  The support stays bounded below, so
the image is in \(\PLit\).  Applying a ring homomorphism blockwise to Laurent
series is again a ring homomorphism, because each block of a sum or a Cauchy
product is a finite algebraic expression in the blocks of the arguments; and
\(\iota\) is injective because \(\iota_{0}\) is.  On \(\PLlau\) every block is
a polynomial and \(\iota_{0}\) is the inclusion, so \(\iota\) is the identity
there.

\emph{Not surjective.}  Since \(\operatorname{char}\K=0\), the binomial series
\(\sigma\DefinitionEquals\sum_{k\ge0}\binom{1/2}{k}L^{-k}\in\K[[L^{-1}]]\) is
defined and satisfies \(\sigma^{2}=1+L^{-1}\), the identity holding
coefficientwise because it holds in \(\RationalNumbers[[u]]\) for
\(u=L^{-1}\).  Suppose \(\sigma=\iota_{0}(f)\) for some \(f\in\K(L)\).  Then
\(f^{2}=(L+1)/L\), so \(\bigl(Lf\bigr)^{2}=L(L+1)\) with \(Lf\in\K(L)\).
Writing \(Lf=u/v\) in lowest terms gives \(u^{2}=L(L+1)v^{2}\); comparing
multiplicities of the irreducible \(L\) on the two sides of this equation in
the unique factorisation domain \(\K[L]\) yields an even number on the left and
an odd one on the right, a contradiction.  Hence \(\sigma\) is not in the image
of \(\iota_{0}\), and \(\sigma\) regarded as a \(t^{0}\)-block is not in the
image of \(\iota\).
\end{proof}

\begin{example}[The exact coefficient becomes an infinite series]
\label{plt:ex:div-embedding}
\begin{equation}\label{plt:eq:div-embedding-example}
  \iota\!\left(\frac1{L+1}\right)
  =L^{-1}-L^{-2}+L^{-3}-\cdots
  =\sum_{j\ge1}(-1)^{j-1}L^{-j},
  \qquad
  \iota\!\left(\frac1{L+t}\right)
  =\sum_{n\ge0}(-1)^{n}L^{-n-1}t^{n}.
\end{equation}
The left-hand sides are single finite data; the right-hand sides are infinite.
\end{example}

\begin{proposition}[The iteration order is part of the type]
\label{plt:prop:div-iteration-order}
Regard both \(\K((L^{-1}))((t))\) and \(\K((t))((L^{-1}))\) as sets of families
\((a_{n,j})_{(n,j)\in\IntegerNumbers^{2}}\) in \(\K\) representing
\(\sum a_{n,j}t^{n}L^{-j}\), the first with the support condition ``\(n\) is
bounded below, and for each \(n\) the set of \(j\) with \(a_{n,j}\ne0\) is
bounded below'', the second with the two roles of \(n\) and \(j\) exchanged.
Then neither set contains the other.  Witnesses:
\begin{equation}\label{plt:eq:div-iteration-order}
  \sum_{n\ge0}L^{n}t^{n}\in\K((L^{-1}))((t))\setminus\K((t))((L^{-1})),
  \qquad
  \sum_{j\ge0}L^{-j}t^{-j}\in\K((t))((L^{-1}))\setminus\K((L^{-1}))((t)).
\end{equation}
\end{proposition}

\begin{proof}
The first element has \(a_{n,-n}=1\) for \(n\ge0\) and all other entries zero.
Its outer index \(n\) is bounded below by \(0\), and for each \(n\) exactly one
\(j\) occurs, so the first support condition holds; indeed the element lies
already in \(\PLpow\).  But for the second condition one must read the family
by inner index: the set of \(j\) carrying a nonzero entry is
\(\SetOf{0,-1,-2,\dots}\), which is not bounded below, so the element is not in
\(\K((t))((L^{-1}))\).

The second element has \(a_{-j,j}=1\) for \(j\ge0\) and all other entries zero.
Its inner index \(j\) is bounded below by \(0\), and for each \(j\) exactly one
\(n\) occurs, so it lies in \(\K((t))((L^{-1}))\).  But the set of \(n\)
carrying a nonzero entry is \(\SetOf{0,-1,-2,\dots}\), so it has infinitely
many positive powers of \(X\) and is excluded from \(\K((L^{-1}))((t))\).
\end{proof}

% ed.: the catalogue asserts that the two iteration orders differ "in general";
% no source proves it, and the witnesses above supply the proof.
\begin{warningbox}[title={Expanding at \(L=\infty\) changes the object}]
\label{plt:rmk:div-embedding-warning}
Three distinct things go wrong if \(\iota\) is applied silently.
\begin{enumerate}
\item \emph{Finiteness.}  An exact block such as \(1/(L+1)\) becomes an
      infinite series.  A finite representation in \(\PLit\) therefore needs an
      outer \(t\)-cutoff \emph{and} an inner \(L^{-1}\)-cutoff, or one declared
      staircase in the lexicographic monomial order.  An outer cutoff alone
      does not make the data finite.
\item \emph{Algebra.}  \(\iota\) is injective but not surjective
      (\cref{plt:thm:div-embedding}), so \(\PLrat\) sits inside \(\PLit\) as a
      proper subfield: an identity established in \(\PLit\) need not descend to
      \(\PLrat\), and an element of \(\PLit\) need not be a rational function
      of \(L\) at all.  Moreover the leading data change meaning -- in
      \(\PLlau\) each block has a finite \(\degL\), while in \(\PLit\) the
      relevant invariant is the \(L^{-1}\)-adic order.
\item \emph{Asymptotics.}  Let \(\K\subseteq\ComplexNumbers\) and let
      \(f\in\K(L)\) have largest pole modulus \(\rho(f)\).  The numerical
      series \(\iota_{0}(f)\) converges for \(\AbsoluteValue{L}>\rho(f)\), that
      is for \(X>\EulerE^{\rho(f)}\) on the positive ray, and diverges for
      \(\AbsoluteValue{L}<\rho(f)\).  There is in general \emph{no} threshold
      valid for all blocks at once: for
      \(F=\sum_{n\ge0}t^{n}/(L-n)\in\PLrat\) the \(n\)-th block requires
      \(X>\EulerE^{n}\).  Hence ``\(F\) equals its \(L^{-1}\)-expansion for
      large \(X\)'' is not a theorem about \(F\); it is a statement about one
      block at a time, and any uniform claim needs its own proof with the
      regime printed as \(\BigOAt{\cdot}{X\to+\infty}\).
\end{enumerate}
\end{warningbox}

\begin{remark}[When to expand and when not to]
Keeping \(1/(L+1)\) unexpanded is exact, finite, and preserves cancellation;
expanding it resolves the finer scale
\(t\precdom\cdots\precdom L^{-2}\precdom L^{-1}\precdom1\) and is required as
soon as one must compare two blocks of the same \(t\)-order.  The decision is a
modelling decision and should be recorded, not inferred: a system that expands
rational coefficients automatically will destroy exact cancellations and swell
expressions, while one that never expands cannot answer questions about the
inner scale.  The same tension appears in every multiseries implementation
\cite{salvy-shackell-1998,salvy-shackell-1999}.
\end{remark}

\begin{summarybox}
An element of \(\PLlau\) is invertible there exactly when its leading
coefficient block is a nonzero scalar; the leading \emph{block} is a
polynomial, and factoring out the dominant monomial does not by itself invert.
Reciprocals of units exist, are unique, and are computed by a triangular
recurrence whose truncations are certified by the residual
\(B\Trunc{C}{N}=1+\FormalBigOAt{t^{N}}{t}\); relative precision is exactly
preserved, so \(M\) leading input blocks give \(M\) leading output blocks.
Newton's map \(C\mapsto C(2-BC)\) doubles the correct order and, with the input
requirement of \cref{plt:thm:div-newton}\,(4), yields exact truncations of the
reciprocal; it beats the recurrence only when block multiplication is
subquadratic.  Unrestricted division of coefficient blocks forces \(\PLrat\),
the smallest Laurent field over an intermediate coefficient ring, which
strictly contains the fraction field of \(\PLlau\).  Expanding each rational
coefficient at \(L=\infty\) embeds \(\PLrat\) into \(\PLit\) injectively but
not surjectively, and converts one exact datum into an infinite series with no
uniform analytic threshold.
\end{summarybox}

\clearpage
\appendix
\part*{Appendices}
\addcontentsline{toc}{part}{Appendices}

% ---- fragment 17-tables ----
\section{Tables of the Lambert polynomials}
\label{plt:sec:lambert-tables}

All entries below were produced from the operator definition of
\cref{plt:def:lambert-polynomials} in exact rational arithmetic and then
verified three independent ways: against the Stirling closed form of
\cref{plt:thm:lambert-stirling}, against the differential recurrence of
\cref{plt:thm:lambert-recurrence}, and --- the only check that does not
presuppose any of the derived identities --- by substituting the resulting
series into the defining equation
\(\WrightOmega+\log\WrightOmega=X\) and confirming that every coefficient block
through outer order \(t^{8}\) vanishes identically.

\subsection{The scaled integer family}

By \cref{plt:cor:lambert-integrality}, \(n!\LamP{n}(u)\in\IntegerNumbers[u]\)
for \(n\ge1\); the table lists that integer polynomial together with the
factorial that divides it, which is the storage normalization
\(\LamPsc{n}=n!\LamP{n}\).  Read the row as
\(\LamP{n}(u)=\bigl(\text{entry}\bigr)/n!\).  The sources tabulate these
polynomials only through \(n=8\); rows \(9\)--\(12\) are new here.

\begin{longtable}{@{}r r p{0.70\textwidth}@{}}
\caption{The Lambert polynomials $\LamP{n}(u)=\LamPsc{n}(u)/n!$.}\label{plt:tab:lambert-polynomials}\\
\toprule
$n$ & $n!$ & \multicolumn{1}{c}{$\LamPsc{n}(u)=n!\,\LamP{n}(u)$}\\
\midrule
\endfirsthead
\toprule
$n$ & $n!$ & \multicolumn{1}{c}{$\LamPsc{n}(u)=n!\,\LamP{n}(u)$}\\
\midrule
\endhead
$0$ & $1$ & $-u$\\[0.25em]
$1$ & $1$ & $u$\\[0.25em]
$2$ & $2$ & $u^{2}-2u$\\[0.25em]
$3$ & $6$ & $2u^{3}-9u^{2}+6u$\\[0.25em]
$4$ & $24$ & $6u^{4}-44u^{3}+72u^{2}-24u$\\[0.25em]
$5$ & $120$ & $24u^{5}-250u^{4}+700u^{3}-600u^{2}+120u$\\[0.25em]
$6$ & $720$ & $120u^{6}-1{,}644u^{5}+6{,}750u^{4}-10{,}200u^{3}+5{,}400u^{2}-720u$\\[0.25em]
$7$ & $5{,}040$ & $720u^{7}-12{,}348u^{6}+68{,}208u^{5}-154{,}350u^{4}+147{,}000u^{3}-52{,}920u^{2}+5{,}040u$\\[0.25em]
$8$ & $40{,}320$ & $5{,}040u^{8}-104{,}544u^{7}+735{,}392u^{6}-2{,}274{,}384u^{5}+3{,}292{,}800u^{4}-2{,}163{,}840u^{3}+564{,}480u^{2}-40{,}320u$\\[0.25em]
$9$ & $362{,}880$ & $40{,}320u^{9}-986{,}256u^{8}+8{,}504{,}928u^{7}-33{,}911{,}136u^{6}+67{,}885{,}776u^{5}-68{,}584{,}320u^{4}+33{,}022{,}080u^{3}-6{,}531{,}840u^{2}+362{,}880u$\\[0.25em]
$10$ & $3{,}628{,}800$ & $362{,}880u^{10}-10{,}265{,}760u^{9}+105{,}543{,}000u^{8}-521{,}049{,}600u^{7}+1{,}357{,}398{,}000u^{6}-1{,}913{,}375{,}520u^{5}+1{,}428{,}840{,}000u^{4}-526{,}176{,}000u^{3}+81{,}648{,}000u^{2}-3{,}628{,}800u$\\[0.25em]
$11$ & $39{,}916{,}800$ & $3{,}628{,}800u^{11}-116{,}915{,}040u^{10}+1{,}402{,}893{,}360u^{9}-8{,}325{,}405{,}000u^{8}+27{,}062{,}085{,}600u^{7}-50{,}009{,}929{,}200u^{6}+52{,}481{,}610{,}720u^{5}-30{,}187{,}080{,}000u^{4}+8{,}781{,}696{,}000u^{3}-1{,}097{,}712{,}000u^{2}+39{,}916{,}800u$\\[0.25em]
$12$ & $479{,}001{,}600$ & $39{,}916{,}800u^{12}-1{,}446{,}526{,}080u^{11}+19{,}921{,}172{,}832u^{10}-138{,}940{,}660{,}320u^{9}+546{,}429{,}272{,}400u^{8}-1{,}267{,}789{,}406{,}400u^{7}+1{,}754{,}714{,}586{,}240u^{6}-1{,}426{,}718{,}240{,}640u^{5}+652{,}040{,}928{,}000u^{4}-153{,}679{,}680{,}000u^{3}+15{,}807{,}052{,}800u^{2}-479{,}001{,}600u$\\[0.25em]
\bottomrule
\end{longtable}

\begin{remark}[The $n=0$ row is not an exception]
\label{plt:rmk:lambert-zero-row}
\(\LamP{0}(u)=-u\) has degree~\(1\), not~\(0\).  The degree law
\(\deg\LamP{n}=n\) of \cref{plt:cor:lambert-coefficients} is therefore stated
only for \(n\ge1\), and the integrality normalization
\(\LamPsc{n}=n!\LamP{n}\) is likewise a statement about \(n\ge1\); the table's
\(n=0\) row records the value, not an instance of either law.
\end{remark}

\subsection{The coefficient triangle}

Write \(\LamP{n}(u)=\sum_{m=1}^{n}\lambda_{n,m}u^{m}\) for \(n\ge1\).  By
\cref{plt:thm:lambert-stirling},
\[
 \lambda_{n,m}=\frac{(-1)^{n-m}}{m!}\stirlingone{n}{n-m+1},
\]
so the signs alternate along each row and the magnitudes are unsigned Stirling
cycle numbers divided by a factorial.  The next table lists the integers
\(m!\,(-1)^{n-m}\lambda_{n,m}=\stirlingone{n}{n-m+1}\), which is the row of the
Stirling triangle read from the right.

\begin{table}[htbp]\centering\small
\caption{$\stirlingone{n}{n-m+1}=m!\,\AbsoluteValue{\lambda_{n,m}}$, the unsigned coefficient data of $\LamP{n}$.}\label{plt:tab:lambert-triangle}
\begin{tabular}{@{}rrrrrrrrrr@{}}
\toprule
$n\backslash m$ & $1$ & $2$ & $3$ & $4$ & $5$ & $6$ & $7$ & $8$ & $9$\\
\midrule
$1$ & $1$ &  &  &  &  &  &  &  & \\
$2$ & $1$ & $1$ &  &  &  &  &  &  & \\
$3$ & $1$ & $3$ & $2$ &  &  &  &  &  & \\
$4$ & $1$ & $6$ & $11$ & $6$ &  &  &  &  & \\
$5$ & $1$ & $10$ & $35$ & $50$ & $24$ &  &  &  & \\
$6$ & $1$ & $15$ & $85$ & $225$ & $274$ & $120$ &  &  & \\
$7$ & $1$ & $21$ & $175$ & $735$ & $1{,}624$ & $1{,}764$ & $720$ &  & \\
$8$ & $1$ & $28$ & $322$ & $1{,}960$ & $6{,}769$ & $13{,}132$ & $13{,}068$ & $5{,}040$ & \\
$9$ & $1$ & $36$ & $546$ & $4{,}536$ & $22{,}449$ & $67{,}284$ & $118{,}124$ & $109{,}584$ & $40{,}320$\\
\bottomrule
\end{tabular}
\end{table}

The first column is \(\stirlingone{n}{n}=1\), which is
\cref{plt:cor:lambert-coefficients}(iv) in the form
\(\lambda_{n,1}=(-1)^{n-1}\).  The last entry of row~\(n\) is
\(\stirlingone{n}{1}=(n-1)!\), giving the leading coefficient
\(\lambda_{n,n}=1/n\).  The penultimate entry is
\(\stirlingone{n}{2}=(n-1)!H_{n-1}\), giving
\(\lambda_{n,n-1}=-H_{n-1}\).

\subsection{Numerical check values}

The following exact values are convenient unit tests for an
implementation of the recurrence; each is \(\LamP{n}(1)\), which by
\cref{plt:thm:lambert-stirling} equals
\(\sum_{m=1}^{n}(-1)^{n-m}\stirlingone{n}{n-m+1}/m!\).

\begin{table}[htbp]\centering\small
\caption{Exact values $\LamP{n}(1)$.}\label{plt:tab:lambert-at-one}
\begin{tabular}{@{}r l r l@{}}
\toprule
$n$ & $\LamP{n}(1)$ & $n$ & $\LamP{n}(1)$\\
\midrule
$0$ & $-1$ & $7$ & $\frac{15}{56}$\\
$1$ & $1$ & $8$ & $\frac{457}{1260}$\\
$2$ & $-\frac{1}{2}$ & $9$ & $-\frac{49}{90}$\\
$3$ & $-\frac{1}{6}$ & $10$ & $-\frac{391}{2016}$\\
$4$ & $\frac{5}{12}$ & $11$ & $\frac{32213}{36960}$\\
$5$ & $-\frac{1}{20}$ & $12$ & $-\frac{299363}{1425600}$\\
$6$ & $-\frac{49}{120}$ &  & \\
\bottomrule
\end{tabular}
\end{table}

% ---- fragment 19-notation ----
\section{Notation, and a concordance with the six sources}
\label{plt:sec:notation}

\subsection{The normative source}

Every mathematical symbol in this volume is defined normatively in the corpus
notation catalogue, \emph{FabiusFunction Mathematical Notation Catalogue}, whose
entries \texttt{polynomial-logarithmic-ambient},
\texttt{polynomial-logarithmic-leading-data},
\texttt{polynomial-logarithmic-truncation},
\texttt{polynomial-logarithmic-differential},
\texttt{polynomial-logarithmic-composition},
\texttt{polynomial-logarithmic-reversion},
\texttt{wright-omega-lambert-polynomials}, and
\texttt{lambert-branch-expansions} cover exactly the material of this volume.
This document introduces no new mathematical notation.  Its preamble defines
short local aliases --- \(\PLpow\), \(\PLlau\), \(\PLrat\), \(\PLit\),
\(\Gnear\), \(\ordt\), \(\lm\), \(\lc\) and the rest --- and every one of them
expands to a catalogue command; the aliases exist to keep displayed formulas
readable, not to name anything the catalogue does not already name.

\subsection{Symbols used throughout}

\begin{longtable}{@{}l l p{0.56\textwidth}@{}}
\caption{Principal notation.}\label{plt:tab:notation}\\
\toprule
Symbol & Alias & Meaning\\
\midrule
\endfirsthead
\toprule
Symbol & Alias & Meaning\\
\midrule
\endhead
\(X\) & --- & the large variable; the default regime is \(X\to+\infty\) along
  the positive real ray\\
\(t\) & --- & \(X^{-1}\), the outer generator\\
\(L\) & --- & \(\log X\), the logarithmic generator, formally independent of
  \(t\)\\
\(\K\) & \verb|\K| & the coefficient field, of characteristic zero\\
\(\PLpow\) & \verb|\PLpow| & \(\K[L][[t]]\), the polynomial-log power-series
  ring\\
\(\PLlau\) & \verb|\PLlau| & \(\K[L]((t))\), its Laurent extension: finitely
  many positive powers of \(X\)\\
\(\PLrat\) & \verb|\PLrat| & \(\K(L)((t))\), the rational-log Laurent field\\
\(\PLit\) & \verb|\PLit| & \(\K((L^{-1}))((t))\), the iterated Laurent-log
  field\\
\(\Gnear\) & \verb|\Gnear| & \(\SetOf{X+A:A\in\PLpow}\), the near-identity
  group under composition\\
\(\ordt F\) & \verb|\ordt| & the outer valuation, with
  \(\ordt 0=+\infty\)\\
\(\lm(F),\ \lc(F)\) & \verb|\lm|, \verb|\lc| & leading monomial and scalar
  leading coefficient; the leading \emph{block} is the whole polynomial
  \(p_{\ordt F}\), a third and different object\\
\(\degL p\) & \verb|\degL| & degree in \(L\)\\
\(\dprofile{F}(n)\) & \verb|\dprofile| & the log-degree profile, \(-\infty\)
  where the block vanishes\\
\(\Trunc{F}{N}\) & \verb|\Trunc| & the \emph{exclusive} truncation
  \(\sum_{n<N}p_n(L)t^n\)\\
\(\FormalBigOAt{t^N}{t}\) & --- & a formal valuation tail: coefficient
  agreement below outer order \(N\), and nothing analytic\\
\(\BigOAt{A(X)}{X\to\infty}\) & --- & an analytic estimate, with the regime
  printed\\
\(\DX\) & \verb|\DX| & the distinguished derivation
  \(t\partial_L-t^{2}\partial_t\)\\
\(\shiftop{n}{k}\) & \verb|\shiftop| & \(\prod_{j=0}^{k-1}(\partial_L-n-j)\),
  empty and equal to the identity at \(k=0\)\\
\(\rev{F}\) & \verb|\rev| & compositional inverse,
  \(F\compose\rev{F}=\rev{F}\compose F=X\)\\
\(\recip{F}\) & \verb|\recip| & multiplicative inverse, \(F\recip{F}=1\)\\
\(\RevErr{F}{H}\) & \verb|\RevErr| & the reversion residual
  \(F\compose H-X\); \emph{both} arguments are part of the object\\
\(\WrightOmega\) & --- & Wright omega, the inverse of \(Y\mapsto Y+\log Y\)\\
\(\LamP{n}\) & \verb|\LamP| & the canonical zero-based Lambert polynomials\\
\(\LamPsc{n}\) & \verb|\LamPsc| & the scaled family \(n!\LamP{n}\), for
  \(n\ge1\)\\
\(\stirlingone{n}{k}\) & \verb|\stirlingone| & unsigned Stirling numbers of the
  first kind (cycle numbers)\\
\(\stirlingsigned{n}{k}\) & \verb|\stirlingsigned| & signed Stirling numbers of
  the first kind\\
\(\precdom,\succdom\) & --- & the dominance order on monomials\\
\bottomrule
\end{longtable}

\subsection{Concordance with the absorbed sources}

The six sources used mutually incompatible local abbreviations, and in two
cases the \emph{same} abbreviation for different objects.  A reader returning
to a source through the Git history needs the following table.

\begin{longtable}{@{}l l l@{}}
\caption{Local abbreviations of the absorbed sources and their meaning here.
S1--S6 are as identified in \cref{plt:tab:provenance}.}
\label{plt:tab:concordance}\\
\toprule
Source & Local & This volume\\
\midrule
\endfirsthead
\toprule
Source & Local & This volume\\
\midrule
\endhead
S1 & \verb|\Plog|, \verb|\Rlog|, \verb|\Hlog| &
  \(\PLlau\), \(\PLrat\), \(\PLit\)\\
S1 & \verb|\val|, \verb|\trunc| & \(\ordt\), exclusive \(\Trunc{\cdot}{N}\)\\
S1 & \verb|\invcomp| & \(\rev{\cdot}\)\\
S2 & \verb|\val| & \(\ordt\)\\
S2 & \verb|\rev|, \verb|\recip| & \(\rev{\cdot}\), \(\recip{\cdot}\)
  (already the catalogue commands)\\
S2 & \verb|\Dop| \(=\mathrm{D}\) & \(\DX\)\\
S2 & \verb|\normalexp| & normal-ordered shifted exponential; expanded inline
  here\\
S3 & \verb|\Rring|, \verb|\Pring|, \verb|\Fring|, \verb|\Ggroup| &
  \(\PLpow\), \(\PLlau\), \(\PLit\), \(\Gnear\)\\
S3 & \verb|\ordt|, \verb|\D| \(=\mathrm{D}\) & \(\ordt\), \(\DX\)\\
S4 & \verb|\PL| & \(\PLlau\) --- \emph{not} \(\PLpow\); see below\\
S4 & \verb|\RL|, \verb|\HL|, \verb|\Gnear| & \(\PLrat\), \(\PLit\),
  \(\Gnear\)\\
S4 & \verb|\ordx|, \verb|\degell| & \(\ordt\), \(\degL\)\\
S5 & \verb|\PL| & \(\PLpow\) --- \emph{not} \(\PLlau\); see below\\
S5 & \verb|\PLL|, \verb|\QPL| & \(\PLlau\), \(\PLrat\)\\
S5 & \verb|\D| \(=\mathcal{D}\) & \(\DX\)\\
S6 & \verb|\PLpoly|, \verb|\PLfield| & \(\PLlau\), \(\PLit\)\\
S6 & \verb|\EulerDerivation| & \(\EulerOp=X\DX\)\\
S6 & \verb|\trunc{F}{N}| \(=[F]_{\le N}\) & \emph{inclusive}; equals
  \(\Trunc{F}{N+1}\)\\
S6 & \verb|\compinv| & \(\rev{\cdot}\)\\
\bottomrule
\end{longtable}

\begin{warningbox}[title={Two abbreviations that collide across sources}]
\label{plt:rmk:concordance-collisions}
\begin{enumerate}
\item \verb|\PL| denotes the \emph{Laurent} ring \(\K[L]((t))\) in S4 and the
  \emph{power-series} ring \(\K[L][[t]]\) in S5.  A statement transported from
  one to the other without checking which ring is meant can silently acquire or
  lose the finitely many positive powers of \(X\); the unit criterion and the
  antiderivative obstruction both change.  In S4 the power-series subring was
  written \([\PLlau]_{\ge0}\), which is \(\PLpow\).
\item \verb|\trunc| is the exclusive \(\operatorname{Trunc}_{<N}\) in S1 and the
  inclusive \([\,\cdot\,]_{\le N}\) in S6.  Every cutoff index transported from
  S6 into this volume has been shifted by one block.  A residual certificate
  stated with the wrong convention is off by exactly one order, which is
  precisely the amount that makes a correct algorithm look wrong and an
  incorrect one look right.
\end{enumerate}
\end{warningbox}

\subsection{Conventions that are not merely notational}

Three conventions are recorded here because getting them wrong changes
theorems, not only glyphs.

\begin{enumerate}
\item \textbf{Empty products and sums.}  \(\shiftop{n}{0}=1\); the
  Lagrange--B\"urmann coefficient product is empty at \(r=1\); \(C_{0,0}=1\)
  and \(C_{m,0}=0\) for \(m>0\); \(\stirlingsigned{0}{0}=1\).  Every statement
  below is written so that these conventions make the boundary case correct
  rather than exceptional.
\item \textbf{The three \(\BigO\) roles.}  A formal valuation tail, an analytic
  estimate, and an algorithmic cost bound are three different assertions and
  are printed with three different tokens.  None may be inferred from another.
  In particular, \(\FormalBigOAt{t^{N}}{t}\) carries no uniformity in \(L\)
  whatsoever, and a bound uniform in \(L\) is a separate theorem with separate
  hypotheses.
\item \textbf{\(t^{N}\PLpow\) is not an ideal of \(\PLlau\).}  In the Laurent
  ring \(t\) is a unit, so the principal ideal it generates is the whole ring.
  Quotient-algebra language is therefore available on \(\PLpow\) and not on
  \(\PLlau\); on the latter, use \(\ordt(F-G)\ge N\), the formal tail token, or
  equality of exclusive truncations.
\end{enumerate}

% ---- fragment 98-provenance ----
\section{Provenance of the consolidation}
\label{plt:sec:provenance}

\subsection{What was merged}

This volume is an editorial consolidation, carried out on 2~September 2026, of
six independently written article packages that arrived together as bare
directories in the direct-arrival commit
\texttt{730e1763291099cd50ca1e20ed2c62c38d95ab4f} and were filed on
1~September 2026 as standalone quick-gate receipts.  No archive or checksum
ledger accompanied them.  All six treat the same subject -- the algebra,
division, composition, and asymptotic reversion of polynomial--logarithmic
transseries, with Lambert \(\LambertW\) as the guiding example -- from
independent starting points, with substantial overlap and no containment
relation among them.

\emph{Editorial} merge means: every result is stated once, in the strongest
form that the six sources jointly support and that we can prove; the best
available proof is retained and gaps are filled; conflicting conventions are
resolved against the corpus notation catalogue; errors are corrected and the
corrections are marked in the source with \verb|% ed.:| comments; and every
source-specific layer that is not redundant is preserved.  Nothing was
discarded merely for being repeated: repetition was collapsed, content was not.

\subsection{Absorbed sources and their receipts}

Each source is recorded twice: as filed at intake, and as actually absorbed.
The two differ by exactly three bytes per file because a regrouping on
2~September 2026 moved the packages one directory level deeper and the shared
notation input path was repaired from four to five levels
(\cref{plt:tab:provenance}).  No mathematical content changed between the two
states.

\begin{table}[htbp]
\centering
\small
\caption{The six absorbed sources.  ``Intake'' is the receipt recorded in the
group manifest on 1~September 2026; ``absorbed'' is the state merged into this
volume.  The retained arrival PDFs were historical publication checkpoints
that already predated their own sources and were not rebuilt for this merge.}
\label{plt:tab:provenance}
\begin{tabularx}{\textwidth}{@{}lrrX@{}}
\toprule
Source & Lines & Bytes & SHA-256 \\
\midrule
\multicolumn{4}{@{}l}{\textbf{S1} \(=\)
  \texttt{Polynomial-Logarithmic-Transseries-1.tex}} \\
\quad intake    & 4{,}020 & 187{,}071 &
  \texttt{01a03e0934d05b8b94369da9656fc0da4dde9b7c559b4e0738075484e14f4e61} \\
\quad absorbed  & 4{,}020 & 187{,}074 &
  \texttt{4fd42b7e344c485f8d878279dc3b9a6d5dc326a9c739114e443fd9c07a5eb2b8} \\
\addlinespace
\multicolumn{4}{@{}l}{\textbf{S2} \(=\)
  \texttt{polynomial\_logarithmic\_transseries-2.tex}} \\
\quad intake    & 5{,}006 & 173{,}396 &
  \texttt{aa25baa03099472df765d35c03a2bb265dbe9af5aab500ec0a4a04a378ef83e5} \\
\quad absorbed  & 5{,}006 & 173{,}399 &
  \texttt{d172b81453a805366f056f6114bb6df6fdf781146641320185530940754bdefe} \\
\addlinespace
\multicolumn{4}{@{}l}{\textbf{S3} \(=\)
  \texttt{Polynomial\_Logarithmic\_Transseries-3.tex}} \\
\quad intake    & 4{,}249 & 150{,}182 &
  \texttt{0962c15683cb4610d45961c227f6e9b102d66ee37073432d1943d31a1e6ad348} \\
\quad absorbed  & 4{,}249 & 150{,}185 &
  \texttt{d2a8bd25d71cc41fc81c1c75fb66ca75496b8425ac496ce105d3220e1ea07e3f} \\
\addlinespace
\multicolumn{4}{@{}l}{\textbf{S4} \(=\)
  \texttt{Polynomial-Logarithmic-Transseries-4.tex}} \\
\quad intake    & 3{,}132 & 120{,}607 &
  \texttt{387ca51f94292a996fb339ec6dfb3477f68575255b2aaaf539fcffb1c32ff564} \\
\quad absorbed  & 3{,}132 & 120{,}610 &
  \texttt{37026b625cf7b2fd106005b8c89f720b8b724f99a5446ed84ca52531c6e2b18f} \\
\addlinespace
\multicolumn{4}{@{}l}{\textbf{S5} \(=\)
  \texttt{Polynomial\_Logarithmic\_Transseries-5.tex}} \\
\quad intake    & 2{,}443 & 106{,}141 &
  \texttt{1149ae6884e63dc77b747c786d0455419b6c79cdd3b33dc915bf3875cc8d26cc} \\
\quad absorbed  & 2{,}443 & 106{,}144 &
  \texttt{187bd510f33ea4296f0247b286d0fcae6200ce0b5bd175c6c8706a37652dc17a} \\
\addlinespace
\multicolumn{4}{@{}l}{\textbf{S6} \(=\)
  \texttt{polynomial\_logarithmic\_transseries-6.tex}} \\
\quad intake    & 4{,}388 & 155{,}846 &
  \texttt{df4e4bc5932a59e4ed265b32d8463edc65f67cbd2735b34902acf479647f8491} \\
\quad absorbed  & 4{,}388 & 155{,}849 &
  \texttt{16978aae6fb1e458e236cac9807ab81bb5314fdb087a2504a416867a4d43b4c4} \\
\bottomrule
\end{tabularx}
\end{table}

The six retained arrival PDFs, which are historical artefacts and are not
renders of the absorbed sources, had SHA-256 values
\texttt{a4fc4af0\ldots69e886} (S1, 119 custom 522-by-738-point pages),
\texttt{5e9ff596\ldots bfc68e} (S2, 102 custom pages),
\texttt{3f7c4bc1\ldots58a4af} (S3, 87 Letter pages),
\texttt{c2d75b35\ldots45a3eb} (S4, 47 A4 pages),
\texttt{189e95ab\ldots58a2db} (S5, 44 Letter pages), and
\texttt{b5142bad\ldots467aa} (S6, 100 A4 pages).  Git history is the archive
for all twelve files.

\subsection{What each source contributed}

No source was a superset of the others, and each contributed at least one layer
that no other supplied.

\begin{itemize}
\item \textbf{S1} contributed the block-algebra presentation of the scale, the
  systematic comparison of four reversal algorithms, the residual-certificate
  and reproducible-verification discipline, the Dulac and power--log normal-form
  literature, and the largest exercise set.
\item \textbf{S2} contributed the \(t\)-adic completeness treatment, the
  normal-ordered composition formalism, the power--log series near a finite
  point with the Dulac conventions, the derivation of the universal composition
  formulas, and the pitfalls-and-diagnostics decision procedure.
\item \textbf{S3} contributed the formal summability criterion, the degree law
  for polynomial perturbations, the Lagrange--B\"urmann transport formula, the
  fully formal proof of Lagrange--B\"urmann inversion at infinity, and the
  reusable inverse templates beyond Lambert.
\item \textbf{S4} contributed the sharpest statement of the unit criterion and
  the field property, the finite reversal certificate, and closed-form
  coefficient extraction.
\item \textbf{S5} contributed the tightest and most completely proved core --
  it is the spine of Parts~I--IV -- together with the Laplace transform of the
  Lambert polynomials, the resulting exponential-moment cancellations, and the
  generalized Lambert expansion.
\item \textbf{S6} contributed the filtered Lipschitz gain underlying the
  fixed-point construction, the composed-observable form of
  Lagrange--B\"urmann, the systematic change of variable between infinity and a
  finite point, and the failure-mode catalogue.
\end{itemize}

\subsection{Conventions reconciled}

The six sources used mutually incompatible local conventions.  All are resolved
here in favour of the corpus notation catalogue; the concordance in
\cref{plt:sec:notation} records the mapping in full.  The reconciliations that
change a formula rather than only a glyph are:

\begin{itemize}
\item \textbf{Truncation.}  S6 used the inclusive cutoff \([F]_{\le N}\); every
  other source and the catalogue use the exclusive
  \(\Trunc{F}{N}\).  The relation is
  \([F]_{\le N}=\Trunc{F}{N+1}\), so every S6 statement transported here has
  had its cutoff index shifted by one.
\item \textbf{Lambert polynomial indexing.}  The sources split between a
  zero-based family with \(\LamP{0}(u)=-u\) and a one-based family that
  separates the logarithmic block.  The catalogue's zero-based normalization is
  used throughout; the degree law \(\deg\LamP{n}=n\) is therefore stated only
  for \(n\ge1\), since \(\LamP{0}=-u\) has degree~\(1\).
\item \textbf{Coefficient extraction.}  Two sources printed \([M]F\) and two
  printed \([F]_{M}\) for the same operation.  This volume writes the series to
  the right of the bracket throughout.
\item \textbf{The three \(\BigO\) roles.}  Formal valuation tails, analytic
  estimates, and algorithmic cost bounds are printed with three different
  tokens and are never interchanged.  Several source statements that used one
  glyph for all three have been split accordingly.
\end{itemize}

\subsection{Status of the claims}

Every theorem, proposition, lemma, and corollary in this volume carries a
proof.  Where a source asserted a result without proof, the proof is supplied
here or the assertion is demoted to an explicitly labelled conjecture or open
problem.  Where a source claimed analytic validity on the strength of formal
algebra alone, the claim has been weakened to what the algebra establishes and
the missing analytic hypothesis is named.  Corrections to the sources are
marked in this document's own source with \verb|% ed.:| comments at the point
of repair, and are collected in \cref{plt:sec:repairs}.

None of the proofs in this volume is a machine-checked proof.  The corpus
formalizes the real Lambert pair and its branch-point geometry in Lean, and the
concordance in \cref{plt:sec:notation} records which statements here correspond
to declarations that exist there; the remaining results are human-proved
frontier mathematics.

\clearpage
\begin{thebibliography}{99}
\addcontentsline{toc}{section}{References}

\bibitem{adh-book}
M.~Aschenbrenner, L.~van den Dries, and J.~van der Hoeven,
\emph{Asymptotic Differential Algebra and Model Theory of Transseries},
Annals of Mathematics Studies, vol.~195, Princeton University Press,
Princeton, 2017.
\href{https://doi.org/10.1515/9781400885411}{doi:10.1515/9781400885411}.

\bibitem{cohen-branch-identities}
H.~Cohen,
``Lambert \(\LambertW\)-function branch identities,''
arXiv:2012.11698, revised 2021.

\bibitem{comtet1970}
L.~Comtet,
``Inversion de \(y^{\alpha}\EulerE^{y}\) et \(y\log y\),''
\emph{Comptes Rendus de l'Acad\'emie des Sciences Paris, S\'erie A}
\textbf{270} (1970), 1085--1088.

\bibitem{comtet-book}
L.~Comtet,
\emph{Advanced Combinatorics: The Art of Finite and Infinite Expansions},
D.~Reidel, Dordrecht, 1974.

\bibitem{corless-et-al-1996}
R.~M. Corless, G.~H. Gonnet, D.~E.~G. Hare, D.~J. Jeffrey, and D.~E. Knuth,
``On the Lambert \(\LambertW\) function,''
\emph{Advances in Computational Mathematics} \textbf{5} (1996), 329--359.
\href{https://doi.org/10.1007/BF02124750}{doi:10.1007/BF02124750}.

\bibitem{corless-jeffrey-knuth-1997}
R.~M. Corless, D.~J. Jeffrey, and D.~E. Knuth,
``A sequence of series for the Lambert \(\LambertW\) function,''
in \emph{Proceedings of the 1997 International Symposium on Symbolic and
Algebraic Computation (ISSAC~'97)}, ACM, New York, 1997, 197--204.
\href{https://doi.org/10.1145/258726.258783}{doi:10.1145/258726.258783}.

\bibitem{wright-omega}
R.~M. Corless and D.~J. Jeffrey,
``The Wright \(\omega\) function,''
in \emph{Artificial Intelligence, Automated Reasoning, and Symbolic
Computation}, Lecture Notes in Computer Science, vol.~2385, Springer, Berlin,
2002, 76--89.
\href{https://doi.org/10.1007/3-540-45470-5_10}{doi:10.1007/3-540-45470-5\_10}.

\bibitem{costin-book}
O.~Costin,
\emph{Asymptotics and Borel Summability},
Chapman \& Hall/CRC Monographs and Surveys in Pure and Applied Mathematics,
vol.~141, Boca Raton, 2009.

\bibitem{dahn-goering-1987}
B.~I. Dahn and P.~G\"oring,
``Notes on exponential-logarithmic terms,''
\emph{Fundamenta Mathematicae} \textbf{127} (1987), no.~1, 45--50.
\href{https://doi.org/10.4064/fm-127-1-45-50}{doi:10.4064/fm-127-1-45-50}.

\bibitem{debruijn}
N.~G. de Bruijn,
\emph{Asymptotic Methods in Analysis}, 3rd ed.,
Dover Publications, New York, 1981.

\bibitem{dlmf-W}
NIST Digital Library of Mathematical Functions,
\S4.13, ``Lambert \(\LambertW\)-function,''
F.~W.~J. Olver, A.~B. Olde Daalhuis, D.~W. Lozier, B.~I. Schneider,
R.~F. Boisvert, C.~W. Clark, B.~R. Miller, B.~V. Saunders, H.~S. Cohl, and
M.~A. McClain, eds., release 1.2.7 of 15 June 2026,
\url{https://dlmf.nist.gov/4.13}, accessed 31 August 2026.

\bibitem{ecalle}
J.~\'Ecalle,
\emph{Les Fonctions R\'esurgentes}, 3 vols.,
Publications Math\'ematiques d'Orsay, Orsay, 1981--1985.

\bibitem{edgar-beginners}
G.~A. Edgar,
``Transseries for beginners,''
\emph{Real Analysis Exchange} \textbf{35} (2009/2010), no.~2, 253--310.
\href{https://doi.org/10.14321/realanalexch.35.2.0253}{doi:10.14321/realanalexch.35.2.0253};
arXiv:0801.4877.

\bibitem{edgar-composition}
G.~A. Edgar,
``Transseries: composition, recursion, and convergence,''
arXiv:0909.1259, 2009.

% ed.: the six sources cited this work inconsistently as a 2010 preprint and as
% a 2013 journal article; the published version is given, with the preprint
% retained as a secondary pointer.
\bibitem{edgar-fractional}
G.~A. Edgar,
``Fractional iteration of series and transseries,''
\emph{Transactions of the American Mathematical Society} \textbf{365} (2013),
no.~11, 5805--5832;
arXiv:1002.2378.

\bibitem{flajolet-sedgewick}
P.~Flajolet and R.~Sedgewick,
\emph{Analytic Combinatorics},
Cambridge University Press, Cambridge, 2009.

\bibitem{gessel-lagrange}
I.~M. Gessel,
``Lagrange inversion,''
\emph{Journal of Combinatorial Theory, Series A} \textbf{144} (2016),
212--249;
arXiv:1609.05988.

\bibitem{graham-knuth-patashnik}
R.~L. Graham, D.~E. Knuth, and O.~Patashnik,
\emph{Concrete Mathematics: A Foundation for Computer Science}, 2nd ed.,
Addison--Wesley, Reading, MA, 1994.

\bibitem{hahn}
H.~Hahn,
``\"Uber die nichtarchimedischen Gr\"o\ss ensysteme,''
\emph{Sitzungsberichte der Kaiserlichen Akademie der Wissenschaften Wien,
Mathematisch-Naturwissenschaftliche Klasse, Abteilung IIa}
\textbf{116} (1907), 601--655.

\bibitem{hardy-orders}
G.~H. Hardy,
\emph{Orders of Infinity: The `Infinit\"arcalc\"ul' of Paul du Bois-Reymond},
Cambridge Tracts in Mathematics and Mathematical Physics, no.~12,
Cambridge University Press, Cambridge, 1910.

\bibitem{jeffrey-et-al-1995}
D.~J. Jeffrey, R.~M. Corless, D.~E.~G. Hare, and D.~E. Knuth,
``Sur l'inversion de \(y^{\alpha}\EulerE^{y}\) au moyen des nombres de
Stirling associ\'es,''
\emph{Comptes Rendus de l'Acad\'emie des Sciences Paris, S\'erie I,
Math\'ematique} \textbf{320} (1995), 1449--1452;
English version, arXiv:math/9512230.

\bibitem{mantova-monotonicity}
V.~Mantova,
``Monotonicity and a Taylor approximation theorem for transseries,''
arXiv:2601.07747, version~2, 2026.

\bibitem{mardesic-normalforms}
P.~Marde\v{s}i\'c, M.~Resman, J.-P. Rolin, and V.~\v{Z}upanovi\'c,
``Normal forms and embeddings for power-log transseries,''
\emph{Advances in Mathematics} \textbf{303} (2016), 888--953.
\href{https://doi.org/10.1016/j.aim.2016.08.030}{doi:10.1016/j.aim.2016.08.030}.

\bibitem{mardesic-length}
P.~Marde\v{s}i\'c, M.~Resman, J.-P. Rolin, and V.~\v{Z}upanovi\'c,
``The length of \(\varepsilon\)-neighborhoods of orbits of Dulac maps,''
arXiv:1606.02581, version~3, 2018.

\bibitem{olver-asymptotics}
F.~W.~J. Olver,
\emph{Asymptotics and Special Functions},
AKP Classics, A~K Peters, Wellesley, MA, 1997.

\bibitem{peran-logtransseries}
D.~Peran,
``Normalization of strongly hyperbolic logarithmic transseries and complex
Dulac germs,''
arXiv:2302.14527, 2023.

\bibitem{resman-dulac}
P.~Marde\v{s}i\'c and M.~Resman,
``Analytic moduli for parabolic Dulac germs,''
\emph{Russian Mathematical Surveys} \textbf{76} (2021), no.~3, 389--460.

\bibitem{salvy-shackell-1992}
B.~Salvy and J.~Shackell,
``Asymptotic expansions of functional inverses,''
in P.~S. Wang, ed., \emph{Symbolic and Algebraic Computation: Proceedings of
ISSAC~'92}, ACM Press, New York, 1992, 130--137.

\bibitem{salvy-shackell-1998}
B.~Salvy and J.~Shackell,
``Symbolic asymptotics: functions of two variables, implicit functions,''
\emph{Journal of Symbolic Computation} \textbf{25} (1998), no.~3, 329--349.
\href{https://doi.org/10.1006/jsco.1997.0179}{doi:10.1006/jsco.1997.0179}.

\bibitem{salvy-shackell-1999}
B.~Salvy and J.~Shackell,
``Symbolic asymptotics: multiseries of inverse functions,''
\emph{Journal of Symbolic Computation} \textbf{27} (1999), no.~6, 543--563.
\href{https://doi.org/10.1006/jsco.1999.0281}{doi:10.1006/jsco.1999.0281}.

\bibitem{vdDMM-1997}
L.~van den Dries, A.~Macintyre, and D.~Marker,
``Logarithmic-exponential power series,''
\emph{Journal of the London Mathematical Society} (2) \textbf{56} (1997),
no.~3, 417--434.

\bibitem{vdDMM-2001}
L.~van den Dries, A.~Macintyre, and D.~Marker,
``Logarithmic-exponential series,''
\emph{Annals of Pure and Applied Logic} \textbf{111} (2001), nos.~1--2,
61--113.

\bibitem{vanderhoeven-book}
J.~van der Hoeven,
\emph{Transseries and Real Differential Algebra},
Lecture Notes in Mathematics, vol.~1888, Springer, Berlin, 2006.
\href{https://doi.org/10.1007/3-540-35590-1}{doi:10.1007/3-540-35590-1}.

\end{thebibliography}

\end{document}
